% Équations, inéquations et systèmes — Mathématiques 1re Tech — Cours signé OnBuch+
% Programme officiel MINESEC. Compiler avec : tectonic N01-equations-inequations-systemes.tex
\newcommand{\DOCMATIERE}{Mathématiques}
\newcommand{\DOCNIVEAU}{1re Tech}
\newcommand{\DOCMODULE}{Mathématiques – classe de Première technique}
\newcommand{\DOCLECON}{N01}
\newcommand{\DOCTITRE}{Équations, inéquations et systèmes}
% OnBuch+ — préambule « cours signés » ENRICHI (compilé avec tectonic / XeLaTeX)
% Système visuel « Pop » d'OnBuch+ : fond crème (jamais blanc), bordures encre
% épaisses, ombres « dures » décalées, titres Archivo Black en capitales.
% Version MATHÉMATIQUES : graphes (pgfplots/tikz), boîtes pédagogiques
% (définition, propriété, méthode, attention, le savais-tu, exercice résolu,
% exercice, TP, bilan), page de garde et signature OnBuch+ ancrée en bas.
%
% Macros attendues dans main.tex AVANT \input{preamble} :
%   \DOCMATIERE  \DOCNIVEAU  \DOCMODULE  \DOCLECON (numéro)  \DOCTITRE
\documentclass[11pt]{article}

\usepackage{fontspec}
\usepackage[french]{babel}
\usepackage[a4paper,margin=2cm,top=3.4cm,bottom=2.8cm,headheight=1.4cm,footskip=1.3cm]{geometry}
\usepackage{amsmath,amssymb,amsfonts}
\usepackage{mathtools} % \xmapsto, \coloneqq... souvent utilisés par les modèles
\usepackage{cancel} % simplifications barrées
% \abs{z} (module, valeur absolue) et \norm{u} (norme), utilisés par les modèles
\providecommand{\abs}{}\let\abs\relax\DeclarePairedDelimiter{\abs}{\lvert}{\rvert}
\providecommand{\norm}{}\let\norm\relax\DeclarePairedDelimiter{\norm}{\lVert}{\rVert}
\usepackage[table]{xcolor} % [table] : fournit aussi \rowcolors (alternance de lignes)
\usepackage{pagecolor}
\usepackage{tikz}
\usetikzlibrary{arrows.meta,positioning,calc,shapes.geometric,shapes.misc,shapes.arrows,decorations.pathmorphing,decorations.pathreplacing,decorations.markings,patterns,shadows,mindmap,trees,fit,backgrounds,matrix,intersections,angles,quotes}
\usepackage{pgfplots}
\usepackage{pgf-pie} % diagrammes circulaires \pie (statistiques)
\pgfplotsset{compat=1.18}
\usepgfplotslibrary{fillbetween,groupplots} % aires entre courbes (calcul intégral)
\usepackage{enumitem}
\usepackage{fancyhdr}
% [most] charge toutes les bibliothèques tcolorbox (listings, minted, raster,
% documentation...) et gonfle inutilement la mémoire TeX sur un document long
% (~300 boîtes) : on ne charge que ce qui est réellement utilisé.
\usepackage[skins,breakable]{tcolorbox}
\usepackage{graphicx}
\usepackage{colortbl}
\usepackage{booktabs}
\usepackage{tabularx}
\usepackage{diagbox} % case d'en-tête diagonale des tableaux à double entrée
% Colonnes extensibles centrée / gauche / droite (C, L, R), souvent utilisées par les modèles
\newcolumntype{C}{>{\centering\arraybackslash}X}
\newcolumntype{L}{>{\raggedright\arraybackslash}X}
\newcolumntype{R}{>{\raggedleft\arraybackslash}X}
\usepackage{array}
\usepackage{multicol}
\usepackage{siunitx}
% parse-numbers=false : accepte aussi \SI{2,5 \times 10^{-3}}{...}
\sisetup{locale=FR,output-decimal-marker={,},group-minimum-digits=4,parse-numbers=false}
\DeclareSIUnit\ohm{\ensuremath{\symup{\Omega}}} % U+2126 absent de la police
\usepackage{qrcode}
% \degree : commande fréquemment utilisée par les modèles pour un angle ou une
% température en dehors de \SI{}{} (non fournie par les paquets chargés).
\providecommand{\degree}{^{\circ}}
% \qed (fin de démonstration), utilisé par les modèles sans amsthm
\providecommand{\qed}{\hfill$\square$}
% \barre : double barre des valeurs interdites dans un tableau de variations (tkz-tab)
\providecommand{\barre}{\ensuremath{\|}}
% environnement proof (amsthm non chargé) utilisé par les modèles
\ifdefined\proof\else\newenvironment{proof}[1][Démonstration]{\par\noindent\textit{#1.}\ }{\qed\par}\fi
\usepackage{lastpage}
\usepackage{hyperref}
\hypersetup{hidelinks}

% ============================================================
% Palette « Pop » OnBuch+ — mêmes hex que la classe P (plus_ui.dart)
% ============================================================
\definecolor{poporange}{HTML}{F59321}
\definecolor{popdark}{HTML}{E07A0C}
\definecolor{popcream}{HTML}{FFF6E8}
\definecolor{popink}{HTML}{1C1714}
\definecolor{poppurple}{HTML}{7A5AE0}
\definecolor{popgreen}{HTML}{2E9E6A}
\definecolor{poppink}{HTML}{E0457B}
\definecolor{popblue}{HTML}{2F7DE1}
\definecolor{popgold}{HTML}{C98A12}
\definecolor{popmuted}{HTML}{8A7F76}
\definecolor{popline}{HTML}{EADFD2}
\definecolor{popgray}{HTML}{8A7F76}
\colorlet{popgrey}{popgray}
% Alias tolérants (le modèle nomme parfois les couleurs différemment)
\colorlet{popwhite}{white}
\colorlet{popblack}{popink}
\colorlet{popred}{poppink}
\colorlet{popyellow}{popgold}
% Teintes claires pour fonds de tableaux / zones
\colorlet{popblueL}{popblue!10!white}
\colorlet{poporangeL}{poporange!14!white}
\colorlet{popgreenL}{popgreen!12!white}
\colorlet{poppurpleL}{poppurple!12!white}
\colorlet{poppinkL}{poppink!12!white}
\colorlet{popgoldL}{popgold!14!white}

\pagecolor{popcream}

% ============================================================
% Typographie
% ============================================================
\defaultfontfeatures{Ligatures=TeX}
\setmainfont{PlusJakartaSans-Medium.ttf}[Path=../../../fonts/,BoldFont=PlusJakartaSans-Bold.ttf]
% Maths en sans-serif (Fira Math) avec chiffres et lettres droites de Plus
% Jakarta, pour que les formules se fondent dans le texte.
\usepackage[math-style=ISO,bold-style=ISO]{unicode-math}
\setmathfont{FiraMath-Regular.otf}[Path=../../../fonts/]
\setmathfont{PlusJakartaSans-Medium.ttf}[Path=../../../fonts/,range={up/num,up/{Latin,latin},\mathup}]
\usepackage{icomma} % virgule décimale sans espace en mode math
% Caractères mathématiques Unicode absents des polices de texte (Plus Jakarta,
% Archivo Black) : rendus par leur équivalent mathématique, en texte comme en
% formule — sinon ils s'affichent en carré vide (ex. « Arithmétique dans ℤ »).
\usepackage{newunicodechar}
\newunicodechar{ℝ}{\ensuremath{\mathbb{R}}}
\newunicodechar{ℤ}{\ensuremath{\mathbb{Z}}}
\newunicodechar{ℂ}{\ensuremath{\mathbb{C}}}
\newunicodechar{ℕ}{\ensuremath{\mathbb{N}}}
\newunicodechar{ℚ}{\ensuremath{\mathbb{Q}}}
\newunicodechar{𝔻}{\ensuremath{\mathbb{D}}}
\newunicodechar{α}{\ensuremath{\alpha}}
\newunicodechar{β}{\ensuremath{\beta}}
\newunicodechar{γ}{\ensuremath{\gamma}}
\newunicodechar{δ}{\ensuremath{\delta}}
\newunicodechar{ε}{\ensuremath{\varepsilon}}
\newunicodechar{θ}{\ensuremath{\theta}}
\newunicodechar{λ}{\ensuremath{\lambda}}
\newunicodechar{μ}{\ensuremath{\mu}}
\newunicodechar{σ}{\ensuremath{\sigma}}
\newunicodechar{τ}{\ensuremath{\tau}}
\newunicodechar{φ}{\ensuremath{\varphi}}
\newunicodechar{ψ}{\ensuremath{\psi}}
\newunicodechar{ω}{\ensuremath{\omega}}
\newunicodechar{Σ}{\ensuremath{\Sigma}}
% majuscules produites par \MakeUppercase dans les titres (α-aminés -> Α)
\newunicodechar{Α}{\ensuremath{\alpha}}
\newunicodechar{Β}{\ensuremath{\beta}}
\newunicodechar{Γ}{\ensuremath{\Gamma}}
\newunicodechar{Δ}{\ensuremath{\Delta}}
\newunicodechar{Ε}{\ensuremath{\varepsilon}}
\newunicodechar{Θ}{\ensuremath{\Theta}}
\newunicodechar{Λ}{\ensuremath{\Lambda}}
\newunicodechar{Μ}{\ensuremath{\mu}}
\newunicodechar{Τ}{\ensuremath{\tau}}
\newunicodechar{Φ}{\ensuremath{\Phi}}
\newunicodechar{Ψ}{\ensuremath{\Psi}}
\newunicodechar{Ω}{\ensuremath{\Omega}}
\newunicodechar{↦}{\ensuremath{\mapsto}}
\newunicodechar{∈}{\ensuremath{\in}}
\newunicodechar{∉}{\ensuremath{\notin}}
\newunicodechar{∧}{\ensuremath{\wedge}}
\newunicodechar{⇔}{\ensuremath{\Leftrightarrow}}
\newunicodechar{⟺}{\ensuremath{\Longleftrightarrow}}
\newunicodechar{⇒}{\ensuremath{\Rightarrow}}
\newunicodechar{⟹}{\ensuremath{\Longrightarrow}}
\newunicodechar{∘}{\ensuremath{\circ}}
\newunicodechar{∪}{\ensuremath{\cup}}
\newunicodechar{∩}{\ensuremath{\cap}}
\newunicodechar{≡}{\ensuremath{\equiv}}
\newunicodechar{⊂}{\ensuremath{\subset}}
\newunicodechar{⊥}{\ensuremath{\perp}}
\newunicodechar{∃}{\ensuremath{\exists}}
\newunicodechar{∀}{\ensuremath{\forall}}
\newunicodechar{∛}{\ensuremath{\sqrt[3]{\,}}}
\newunicodechar{∜}{\ensuremath{\sqrt[4]{\,}}}
\newunicodechar{⃗}{\ensuremath{{}^{\to}}}
\newunicodechar{ⁿ}{\ifmmode^{n}\else\textsuperscript{\ensuremath{n}}\fi}
\newunicodechar{ˣ}{\ifmmode^{x}\else\textsuperscript{\ensuremath{x}}\fi}
\newunicodechar{ᵇ}{\ifmmode^{b}\else\textsuperscript{\ensuremath{b}}\fi}
\newunicodechar{ʳ}{\ifmmode^{r}\else\textsuperscript{\ensuremath{r}}\fi}
\newunicodechar{ᵘ}{\ifmmode^{u}\else\textsuperscript{\ensuremath{u}}\fi}
\newunicodechar{ʸ}{\ifmmode^{y}\else\textsuperscript{\ensuremath{y}}\fi}
\newunicodechar{ᵃ}{\ifmmode^{a}\else\textsuperscript{\ensuremath{a}}\fi}
\newunicodechar{ᵉ}{\ifmmode^{e}\else\textsuperscript{\ensuremath{e}}\fi}
\newunicodechar{ᵏ}{\ifmmode^{k}\else\textsuperscript{\ensuremath{k}}\fi}
\newunicodechar{ᵖ}{\ifmmode^{p}\else\textsuperscript{\ensuremath{p}}\fi}
\newunicodechar{ᴺ}{\ifmmode^{N}\else\textsuperscript{\ensuremath{N}}\fi}
\newunicodechar{ᶜ}{\ifmmode^{c}\else\textsuperscript{\ensuremath{c}}\fi}
\newunicodechar{ʰ}{\ifmmode^{h}\else\textsuperscript{\ensuremath{h}}\fi}
\newunicodechar{ᵠ}{\ifmmode^{q}\else\textsuperscript{\ensuremath{q}}\fi}
\newunicodechar{ₙ}{\ifmmode_{n}\else\textsubscript{\ensuremath{n}}\fi}
\newunicodechar{ₐ}{\ifmmode_{a}\else\textsubscript{\ensuremath{a}}\fi}
\newunicodechar{ᵢ}{\ifmmode_{i}\else\textsubscript{\ensuremath{i}}\fi}
\newunicodechar{ᵣ}{\ifmmode_{r}\else\textsubscript{\ensuremath{r}}\fi}
\newunicodechar{ₖ}{\ifmmode_{k}\else\textsubscript{\ensuremath{k}}\fi}
\newunicodechar{ᵦ}{\ifmmode_{\beta}\else\textsubscript{\ensuremath{\beta}}\fi}
\newunicodechar{ₓ}{\ifmmode_{x}\else\textsubscript{\ensuremath{x}}\fi}
\newfontfamily\popdisplay{ArchivoBlack-Regular.ttf}[Path=../../../fonts/]
\newfontfamily\poplogo{SpaceGrotesk-Bold.ttf}[Path=../../../fonts/]
\newfontfamily\popsemi{PlusJakartaSans-SemiBold.ttf}[Path=../../../fonts/]
\newfontfamily\popxbold{PlusJakartaSans-ExtraBold.ttf}[Path=../../../fonts/]
\newfontfamily\popxboldls{PlusJakartaSans-ExtraBold.ttf}[Path=../../../fonts/,LetterSpace=3.0]

\setlist[enumerate]{leftmargin=1.7em,itemsep=5pt,topsep=4pt}
\setlist[itemize]{leftmargin=1.7em,itemsep=4pt,topsep=4pt,label={\color{poporange}\textbullet}}
\setlength{\parindent}{0pt}
\setlength{\parskip}{5pt}
\renewcommand{\arraystretch}{1.25}

% pgfplots : style Pop par défaut
\pgfplotsset{
  every axis/.append style={axis line style={popink,line width=1pt},tick style={popink},
    grid style={popline,line width=0.5pt},label style={font=\small},tick label style={font=\footnotesize}},
  popcurve/.style={poporange,line width=1.8pt,no markers,smooth},
}
% Même style disponible dans un \draw TikZ simple (hors pgfplots)
\tikzset{popcurve/.style={poporange,line width=1.8pt}}

% ============================================================
% Pill / squiggle
% ============================================================
\newcommand{\poppill}[3]{%
  \tikz[baseline=-0.55ex]\node[draw=popink,line width=1.2pt,rounded corners=2.6pt,%
    fill=#2,inner xsep=6.5pt,inner ysep=3pt]{%
    \popxboldls\fontsize{7}{7}\selectfont\color{#3}\MakeUppercase{#1}};%
}
\newcommand{\popsquiggle}[1]{%
  \tikz[baseline=-0.4ex]{
    \draw[#1,line width=2.6pt,line cap=round]
      plot [smooth] coordinates {(0,0.09) (0.34,0.01) (0.68,0.21) (1.02,0.01) (1.36,0.10)};
  }%
}
% Mot-clé surligné (marqueur orange sous le texte)
\newcommand{\cle}[1]{{\popxbold\color{popdark}#1}}

% ============================================================
% En-tête / pied
% ============================================================
\pagestyle{fancy}
\fancyhf{}
\renewcommand{\headrulewidth}{0pt}
\renewcommand{\footrulewidth}{0pt}
\lhead{\raisebox{-0.15ex}{\poplogo\fontsize{13}{13}\selectfont\color{popink}OnBuch\color{poporange}+}}
\rhead{\poppill{\DOCMATIERE\ \textbullet\ \DOCNIVEAU\ \textbullet\ Leçon \DOCLECON}{white}{popink}}
\lfoot{\popxbold\fontsize{7.6}{7.6}\selectfont\color{popmuted}\MakeUppercase{Cours signé OnBuch+}\ \textbullet\ {\popsemi\DOCTITRE}}
\rfoot{\tikz[baseline=-0.6ex]\node[rounded corners=8.5pt,fill=popink,minimum height=17pt,minimum width=17pt,inner xsep=5pt,inner sep=0]{\popxbold\fontsize{8}{8}\selectfont\color{popcream}\thepage\,/\,\pageref*{LastPage}};}

% ============================================================
% Titres
% ============================================================
\newcommand{\chaptitre}[2]{%
  \begin{tcolorbox}[enhanced,colback=poporange,colframe=popink,boxrule=2.6pt,arc=11pt,
      drop shadow=popink,left=15pt,right=15pt,top=12pt,bottom=11pt,before skip=3pt,after skip=18pt]
    \poppill{#2}{popink}{popcream}\par\vspace{9pt}
    {\raggedright\hyphenpenalty=10000\exhyphenpenalty=10000\popdisplay\fontsize{22}{24}\selectfont\color{popcream}\MakeUppercase{#1}\par}
  \end{tcolorbox}
}
% Section numérotée : pastille orange avec le numéro + titre Archivo
\newcounter{popsec}
\newcommand{\coursec}[1]{%
  \stepcounter{popsec}\setcounter{popsub}{0}%
  \par\vspace{16pt}\noindent
  \begin{minipage}{\linewidth}
  \tikz[baseline=-0.7ex]\node[draw=popink,line width=1.6pt,fill=poporange,rounded corners=4pt,minimum size=22pt,inner sep=0]{\popdisplay\fontsize{12}{12}\selectfont\color{popcream}\thepopsec};%
  \hspace{8pt}{\popdisplay\fontsize{15}{17}\selectfont\color{popink}\MakeUppercase{#1}}\par
  \vspace{4pt}\noindent{\color{poporange}\rule{36pt}{3.6pt}}
  \end{minipage}\par\nopagebreak\vspace{8pt}
}
\newcounter{popsub}
\newcommand{\courssub}[1]{%
  \stepcounter{popsub}%
  \par\vspace{9pt}\noindent{\popxbold\fontsize{11.5}{13}\selectfont\color{popink}{\color{poporange}\thepopsec.\thepopsub}\hspace{6pt}#1}\par\nopagebreak\vspace{4pt}
}

% ============================================================
% Boîtes pédagogiques — même recette : fond blanc, bordure épaisse colorée,
% ombre dure encre, badge plein. Titre optionnel [#1].
% ============================================================
\tcbset{popbase/.style={breakable,enhanced,colback=white,boxrule=2.3pt,arc=9pt,
  shadow={3pt}{-3pt}{0pt}{popink},left=14pt,right=14pt,top=11pt,bottom=11pt,before skip=11pt,after skip=15pt,
  fontupper=\popsemi\fontsize{10.6}{14.5}\selectfont\color{popink},
  fontlower=\popsemi\fontsize{10.6}{14.5}\selectfont\color{popink}}}
% Coche dessinée (le glyphe ✓ n'existe pas dans les polices du document)
\newcommand{\popcheck}[1]{\tikz[baseline=-0.5ex]\draw[#1,line width=1.6pt,line cap=round,line join=round](-0.13,0.0)--(-0.04,-0.09)--(0.13,0.1);}
\newcommand{\popboxhead}[3]{\poppill{#1}{#2}{white}\ifx\relax#3\relax\else\hspace{6pt}{\popxbold\fontsize{10.6}{12}\selectfont\color{popink}#3}\fi\par\vspace{7pt}}

\newtcolorbox{aretenir}[1][]{popbase,colframe=popgreen,before upper={\popboxhead{À retenir}{popgreen}{#1}}}
\newtcolorbox{exemplebox}[1][]{popbase,colframe=poporange,before upper={\popboxhead{Exemple}{poporange}{#1}}}
\newtcolorbox{definition}[1][]{popbase,colframe=popblue,before upper={\popboxhead{Définition}{popblue}{#1}}}
\newtcolorbox{propriete}[1][]{popbase,colframe=poppurple,before upper={\popboxhead{Propriété}{poppurple}{#1}}}
\newtcolorbox{methode}[1][]{popbase,colframe=popgold,before upper={\popboxhead{Méthode}{popgold}{#1}}}
\newtcolorbox{attention}[1][]{popbase,colframe=poppink,before upper={\popboxhead{Attention}{poppink}{#1}}}
\newtcolorbox{experience}[1][]{popbase,colframe=popblue,colback=popblueL,before upper={\popboxhead{Expérience}{popblue}{#1}}}
% « Le savais-tu ? » : carte violette pleine (contexte, culture, Cameroun)
\newtcolorbox{savaistu}[1][]{popbase,colback=poppurple,colframe=popink,
  fontupper=\popsemi\fontsize{10.4}{14.2}\selectfont\color{white},
  before upper={\poppill{Le savais-tu\,?}{white}{poppurple}\ifx\relax#1\relax\else\hspace{6pt}{\popxbold\color{white}#1}\fi\par\vspace{7pt}}}
% Exercice résolu : énoncé en haut, \tcblower puis la solution sur fond crème
\newcounter{popexo}\newcounter{popexr}
\newtcolorbox{exoresolu}[1][]{popbase,colframe=poporange,colbacklower=popcream,
  segmentation style={popink,line width=1.2pt,dashed},
  before upper={\stepcounter{popexr}\popboxhead{Exercice résolu \thepopexr}{poporange}{#1}},
  before lower={\poppill{Solution}{popink}{popcream}\par\vspace{6pt}}}
% Exercice (non résolu en place) — la correction est dans la partie Corrigés
\newtcolorbox{exercice}[1][]{popbase,colframe=popink,
  before upper={\stepcounter{popexo}\popboxhead{Exercice \thepopexo}{popink}{#1}}}
% Corrigé
\newtcolorbox{corrige}[1][]{popbase,colframe=popgreen,colback=popgreenL,
  before upper={\popboxhead{Corrigé}{popgreen}{#1}}}
% Objectifs / prérequis / bilan
\newtcolorbox{objectifs}{popbase,colframe=popink,colback=poporangeL,
  before upper={\popboxhead{Objectifs de la leçon}{popink}{}}}
\newtcolorbox{prerequis}{popbase,colframe=popink,colback=popblueL,
  before upper={\popboxhead{Prérequis}{popblue}{}}}
\newtcolorbox{bilan}{popbase,colframe=popink,colback=white,boxrule=2.8pt,
  before upper={\popboxhead{Fiche bilan}{poporange}{L'essentiel à connaître pour l'examen}}}
% Figure encadrée avec légende
\newtcolorbox{popfigure}[1][]{enhanced,breakable=false,colback=white,colframe=popline,boxrule=1.4pt,arc=8pt,
  left=10pt,right=10pt,top=10pt,bottom=8pt,before skip=10pt,after skip=12pt,halign=center,
  lower separated=false,halign lower=center,
  fontlower=\popsemi\fontsize{9}{11.5}\selectfont\color{popmuted},
  #1}
\newcommand{\legende}[1]{\par\vspace{6pt}{\popsemi\fontsize{9}{11.5}\selectfont\color{popmuted}#1}}

% ============================================================
% Signature OnBuch+
% ============================================================
\newcommand{\poplogoinline}[1]{{\poplogo\fontsize{#1}{#1}\selectfont\color{popink}OnBuch\color{poporange}+}}
\newcommand{\signatureonbuch}{%
  \begin{tcolorbox}[enhanced,colback=popink,colframe=popink,boxrule=2.2pt,arc=10pt,
      drop shadow=poporange,left=14pt,right=14pt,top=10pt,bottom=10pt,before skip=0pt,after skip=0pt]
    \tikz[baseline=-0.7ex]\node[circle,fill=popgreen,draw=popcream,line width=1.3pt,minimum size=22pt,inner sep=0]{\popcheck{white}};%
    \hspace{9pt}\begin{minipage}[c]{0.62\linewidth}
      {\popxbold\fontsize{12}{13}\selectfont\color{popcream}Signé\ \textbullet\ {\poplogo OnBuch\color{poporange}+}}\par\vspace{2pt}
      {\raggedright\popsemi\fontsize{8}{10.5}\selectfont\color{popline}\MakeUppercase{Équipe pédagogique OnBuch+}\\\MakeUppercase{Conforme au programme officiel MINESEC}\par}
    \end{minipage}\hfill
    {\poplogo\fontsize{22}{22}\selectfont\color{popcream}OnBuch\color{poporange}+}
  \end{tcolorbox}%
}

% ============================================================
% Page de garde
% #1 titre · #2 matière · #3 niveau · #4 module · #5 n° leçon · #6 durée ·
% #7 description / objectifs (texte ou itemize)
% ============================================================
\newcommand{\pagegarde}[7]{%
  \thispagestyle{empty}
  \newgeometry{a4paper,margin=1.7cm,top=1.6cm,bottom=1.4cm}%
  \begin{tikzpicture}[remember picture,overlay]
    \fill[poporange!18] ([xshift=-3.2cm,yshift=-3.6cm]current page.north east) circle (4.6cm);
    \fill[poppurple!14] ([xshift=2.4cm,yshift=7.6cm]current page.south west) circle (3.8cm);
  \end{tikzpicture}%
  {\poplogo\fontsize{30}{30}\selectfont\color{popink}OnBuch\color{poporange}+}\hfill
  \raisebox{0.6ex}{\poppill{Cours signé \textbullet\ Premium}{popink}{popcream}}\par
  \vspace{1pt}\popsquiggle{poppurple}\par
  \vspace{8pt}
  \begin{tcolorbox}[enhanced,colback=poporange,colframe=popink,boxrule=2.8pt,arc=13pt,
      drop shadow=popink,left=18pt,right=18pt,top=12pt,bottom=13pt,after skip=10pt]
    \poppill{#2 \textbullet\ #3}{popink}{popcream}\hspace{5pt}\poppill{#4}{white}{popink}\par\vspace{8pt}
    {\popdisplay\fontsize{14}{14}\selectfont\color{popink}LEÇON #5}\par\vspace{4pt}
    {\raggedright\hyphenpenalty=10000\exhyphenpenalty=10000\popdisplay\fontsize{25}{27}\selectfont\color{popcream}\MakeUppercase{#1}\par}
    \vspace{8pt}
    {\popxbold\fontsize{9.5}{9.5}\selectfont\color{popink}\MakeUppercase{Durée conseillée : #6}}
  \end{tcolorbox}
  \begin{tcolorbox}[enhanced,colback=white,colframe=popink,boxrule=2.2pt,arc=10pt,
      drop shadow=popink,left=16pt,right=16pt,top=10pt,bottom=10pt,after skip=8pt]
    \poppill{Dans cette leçon}{poppurple}{white}\par\vspace{5pt}
    \popsemi\fontsize{9.3}{12.6}\selectfont\color{popink}#7
  \end{tcolorbox}
  \noindent\begin{minipage}[t]{0.487\linewidth}
    \begin{tcolorbox}[enhanced,colback=popgreen,colframe=popink,boxrule=2.2pt,arc=10pt,
        drop shadow=popink,left=11pt,right=11pt,top=10pt,bottom=10pt,height=3.1cm,valign=center]
      \tcbox[colback=white,colframe=popink,boxrule=1.3pt,arc=5pt,boxsep=0pt,nobeforeafter,
        left=3pt,right=3pt,top=3pt,bottom=3pt,valign=center]{\qrcode[height=1.9cm]{https://play.google.com/store/apps/details?id=com.luvvix.onbuch&hl=fr}}%
      \hspace{8pt}\begin{minipage}[c]{0.5\linewidth}
        {\popdisplay\fontsize{11}{12.5}\selectfont\color{white}\MakeUppercase{Télécharge OnBuch}}\par\vspace{4pt}
        {\raggedright\popsemi\fontsize{8.4}{11}\selectfont\color{white}Gratuit \textbullet\ Google Play\\Tuteur IA, annales, concours, résultats\par}
      \end{minipage}
    \end{tcolorbox}
  \end{minipage}\hfill
  \begin{minipage}[t]{0.487\linewidth}
    \begin{tcolorbox}[enhanced,colback=poppurple,colframe=popink,boxrule=2.2pt,arc=10pt,
        drop shadow=popink,left=11pt,right=11pt,top=10pt,bottom=10pt,height=3.1cm,valign=center]
      \tikz[baseline=-0.7ex]\node[circle,fill=white,draw=popink,line width=1.3pt,minimum size=1.6cm,inner sep=0]{\popdisplay\fontsize{24}{24}\selectfont\color{poppurple}+};%
      \hspace{8pt}\begin{minipage}[c]{0.6\linewidth}
        {\popdisplay\fontsize{11}{12.5}\selectfont\color{white}\MakeUppercase{Passe à OnBuch+}}\par\vspace{4pt}
        {\raggedright\popsemi\fontsize{8.4}{11}\selectfont\color{white}Tous les cours signés, Tuteur IA illimité, fiches PDF — onglet OnBuch+ de l'app\par}
      \end{minipage}
    \end{tcolorbox}
  \end{minipage}
  \vspace{8pt}
  \popcontenu
  \vfill
  \signatureonbuch
  \restoregeometry
  \newpage
}

% Rangée « Ce cours contient » (page de garde)
\newcommand{\poptile}[3]{% #1 couleur #2 gros texte #3 légende
  \begin{tcolorbox}[enhanced,colback=#1,colframe=popink,boxrule=2pt,arc=9pt,shadow={2.5pt}{-2.5pt}{0pt}{popink},
      left=6pt,right=6pt,top=9pt,bottom=9pt,halign=center,valign=center,height=1.75cm,width=\linewidth,nobeforeafter]
    {\popdisplay\fontsize{12.5}{13}\selectfont\color{white}\MakeUppercase{#2}}\par\vspace{3pt}
    {\centering\popsemi\fontsize{7.8}{9.5}\selectfont\color{white}#3\par}
  \end{tcolorbox}}
\newcommand{\popcontenu}{%
  \par\noindent{\popxboldls\fontsize{7.5}{7.5}\selectfont\color{popmuted}CE COURS SIGNÉ CONTIENT}\par\vspace{7pt}
  \noindent\begin{minipage}[t]{0.235\linewidth}\poptile{popblue}{Cours}{complet, illustré, pas à pas}\end{minipage}\hfill
  \begin{minipage}[t]{0.235\linewidth}\poptile{poporange}{Méthodes}{et exercices résolus}\end{minipage}\hfill
  \begin{minipage}[t]{0.235\linewidth}\poptile{poppink}{Exercices}{type examen + corrigés}\end{minipage}\hfill
  \begin{minipage}[t]{0.235\linewidth}\poptile{popgreen}{Fiche bilan}{l'essentiel à retenir}\end{minipage}\par}

% Sommaire
\newtcolorbox{sommaire}{enhanced,colback=white,colframe=popink,boxrule=2.2pt,arc=10pt,
  drop shadow=popink,left=18pt,right=18pt,top=13pt,bottom=13pt,before skip=6pt,after skip=20pt,
  before upper={\poppill{Sommaire}{poppurple}{white}\par\vspace{9pt}\popsemi\fontsize{10.6}{14.5}\selectfont\color{popink}}}

% Signature de fin de cours — poussée tout en bas de la dernière page
\newcommand{\findecours}{%
  \par\vfill\nopagebreak
  \begin{center}\popsquiggle{poporange}\end{center}\vspace{4pt}
  \signatureonbuch
}
\AtBeginDocument{\def\llbracket{\mathopen{[\mkern-3.5mu[}}\def\rrbracket{\mathclose{]\mkern-3.5mu]}}} % intervalles d'entiers (glyphe ⟦ absent de la police)
% Glyphes absents de Fira Math : redéfinis à partir de glyphes disponibles
\AtBeginDocument{%
  \def\setminus{\mathbin{\backslash}}\let\smallsetminus\setminus
  \def\vdots{\mathord{\vbox{\baselineskip3.2pt\lineskiplimit0pt\kern4pt\hbox{.}\hbox{.}\hbox{.}}}}%
  \def\ddots{\mathinner{\mkern1mu\raise6.5pt\hbox{.}\mkern2mu\raise3.6pt\hbox{.}\mkern2mu\raise0.7pt\hbox{.}\mkern1mu}}%
  \def\triangle{\mathord{\scalebox{0.9}{$\symup{\Delta}$}}}%
  \def\nabla{\mathord{\raisebox{\depth}{\scalebox{1}[-1]{$\symup{\Delta}$}}}}%
  \def\checkmark{\popcheck{popdark}}%
  \def\star{\mathbin{\ast}}\def\bigstar{\mathord{\ast}}%
  \def\diamond{\mathbin{\scalebox{0.6}{$\Diamond$}}}%
  \def\bigcup{\mathop{\mathchoice{\raisebox{-0.3em}{\scalebox{1.7}{$\cup$}}}{\scalebox{1.25}{$\cup$}}{\cup}{\cup}}\displaylimits}%
  \def\bigcap{\mathop{\mathchoice{\raisebox{-0.3em}{\scalebox{1.7}{$\cap$}}}{\scalebox{1.25}{$\cap$}}{\cap}{\cap}}\displaylimits}%
  \def\lesseqqgtr{\mathrel{\lessgtr}}\def\bigoplus{\mathop{\oplus}\displaylimits}%
}
\providecommand{\ssi}{si et seulement si} % abréviation fréquente
\AtBeginDocument{\providecommand{\wideparen}{\overparen}} % arc de cercle

%% FIGFIX-BEGIN v4 — correctifs globaux des figures (docs/onbuch-generation/tools/figfix)
\usepackage{adjustbox}\usepackage{etoolbox}
% toute figure TikZ/pgfplots plus large que la ligne est réduite à la largeur disponible
\BeforeBeginEnvironment{tikzpicture}{\begin{adjustbox}{max width=\linewidth}}
\AfterEndEnvironment{tikzpicture}{\end{adjustbox}}
% légendes pgfplots lisibles par-dessus les courbes
\pgfplotsset{every axis/.append style={legend style={fill=white,fill opacity=0.92,text opacity=1,draw=black!15,inner sep=3pt}}}
% étiquettes d'arêtes (syntaxe edge["..."]) avec fond blanc
\tikzset{every edge quotes/.append style={fill=white,inner sep=1.2pt}}
%% FIGFIX-END
\begin{document}

\pagegarde{Équations, inéquations et systèmes}{Mathématiques}{1re Tech}{Mathématiques – classe de Première technique}{N01}{4 séances (fiche de progression harmonisée)}{Cette leçon présente les méthodes de résolution des équations et inéquations du second degré à une inconnue, ainsi que des systèmes de deux équations linéaires à deux inconnues. Elle montre comment ces outils algébriques permettent de modéliser et de résoudre des problèmes concrets de la vie courante et des métiers techniques.
\vspace{4pt}
\begin{multicols}{2}\small
\begin{itemize}
\item À la fin de cette leçon, je sais reconnaître et mettre sous forme canonique un trinôme du second degré
\item À la fin de cette leçon, je sais calculer le discriminant d'une équation du second degré et en déduire le nombre de solutions
\item À la fin de cette leçon, je sais résoudre une équation du second degré par le discriminant, par factorisation ou par la forme canonique
\item À la fin de cette leçon, je sais déterminer le signe d'un trinôme et résoudre une inéquation du second degré
\item À la fin de cette leçon, je sais résoudre un système de deux équations linéaires à deux inconnues par substitution, par combinaison et par comparaison
\item À la fin de cette leçon, je sais interpréter graphiquement les solutions d'une équation ou d'une inéquation du second degré
\end{itemize}\end{multicols}}

\begin{sommaire}
\begin{multicols}{2}
\begin{enumerate}
\item Trinôme du second degré et forme canonique
\item Résolution de l'équation du second degré
\item Factorisation, signe du trinôme et inéquations du second degré
\item Systèmes de deux équations linéaires à deux inconnues
\item Équations et inéquations associées aux fonctions
\item Modélisation et résolution de problèmes concrets
\item Activité d'intégration
\item Méthodes et exercices résolus
\item Exercices
\item Corrigés des exercices
\item Fiche bilan
\end{enumerate}
\end{multicols}
\end{sommaire}

\begin{objectifs}
\begin{itemize}
\item À la fin de cette leçon, je sais reconnaître et mettre sous forme canonique un trinôme du second degré
\item À la fin de cette leçon, je sais calculer le discriminant d'une équation du second degré et en déduire le nombre de solutions
\item À la fin de cette leçon, je sais résoudre une équation du second degré par le discriminant, par factorisation ou par la forme canonique
\item À la fin de cette leçon, je sais déterminer le signe d'un trinôme et résoudre une inéquation du second degré
\item À la fin de cette leçon, je sais résoudre un système de deux équations linéaires à deux inconnues par substitution, par combinaison et par comparaison
\item À la fin de cette leçon, je sais interpréter graphiquement les solutions d'une équation ou d'une inéquation du second degré
\item À la fin de cette leçon, je sais modéliser une situation concrète par une équation, une inéquation ou un système, puis rédiger et justifier ma démarche et ma conclusion
\end{itemize}
\end{objectifs}

\begin{prerequis}
\begin{itemize}
\item Développer, factoriser et utiliser les identités remarquables (classe de 3e/2de)
\item Résoudre une équation et une inéquation du premier degré à une inconnue
\item Résoudre une équation produit nul et utiliser un tableau de signes
\item Représenter graphiquement une fonction affine et lire des coordonnées de points
\item Notion de fonction polynôme du second degré et allure de la parabole (début d'année de 1re)
\end{itemize}
\end{prerequis}

\newpage

\coursec{Trinôme du second degré et forme canonique}

\textbf{Situation d'accroche.} Monsieur Nkoulou, menuisier à Yaoundé, doit fabriquer une table rectangulaire dont l'aire doit être exactement $2$ m\textsuperscript{2}, sachant que la longueur dépasse la largeur de $1$ m. Quelles dimensions choisir ? Si l'on note $x$ la largeur, la longueur vaut $x+1$ et l'aire s'écrit $x(x+1)=2$, c'est-à-dire $x^2+x-2=0$ : ce problème se traduit par une \cle{équation du second degré}. De même, une coopérative de Njombé qui vend des bananes plantains cherche à savoir pour quels prix son bénéfice reste positif : c'est une \cle{inéquation du second degré}. Enfin, pour gérer l'achat groupé de deux produits au marché Mokolo, il faut résoudre un \cle{système} de deux équations à deux inconnues. Cette leçon fournit tous les outils pour résoudre ces situations concrètes.

\textbf{Problématique.} Comment reconnaître, transformer et exploiter une expression du second degré pour résoudre des problèmes de la vie professionnelle et quotidienne ? La première étape consiste à maîtriser une écriture très pratique du trinôme : la \cle{forme canonique}.

\courssub{Le trinôme du second degré}

\textbf{Intuition.} En classe de Seconde, vous avez rencontré la fonction carré $x \mapsto x^2$, dont la courbe est une parabole. En ajoutant un terme du premier degré et une constante, on obtient des expressions comme $x^2+x-2$ (l'aire de la table de M. Nkoulou) ou $2x^2-8x+3$. Ces expressions, appelées \emph{trinômes}, modélisent d'innombrables situations : aires, bénéfices, trajectoires d'un objet lancé, coûts de production.

\begin{definition}[Trinôme du second degré]
On appelle \cle{trinôme du second degré} (ou fonction polynôme de degré $2$) toute fonction $f$ définie sur $\mathbb{R}$ par
\[ f(x) = ax^2 + bx + c \]
où $a$, $b$ et $c$ sont des nombres réels appelés \cle{coefficients}, avec la condition essentielle $a \neq 0$.
\end{definition}

\begin{attention}[La condition $a \neq 0$]
Si $a = 0$, l'expression devient $bx + c$ : ce n'est plus un trinôme du second degré mais une fonction \emph{affine}, étudiée en Seconde. Avant toute résolution, vérifiez toujours que le coefficient de $x^2$ est non nul. Par exemple, $(m-1)x^2 + 3x + 1$ n'est un trinôme du second degré que si $m \neq 1$.
\end{attention}

\begin{exemplebox}[Identifier les coefficients]
\begin{itemize}
\item $f(x) = x^2 + x - 2$ : $a = 1$, $b = 1$, $c = -2$.
\item $g(x) = 2x^2 - 8x + 3$ : $a = 2$, $b = -8$, $c = 3$.
\item $h(x) = -x^2 + 5x$ : $a = -1$, $b = 5$, $c = 0$ (le terme constant est absent, donc nul).
\item $k(x) = 3x^2 - 7$ : $a = 3$, $b = 0$, $c = -7$ (le terme en $x$ est absent).
\end{itemize}
\end{exemplebox}

\courssub{Rappel : les identités remarquables}

Toute la technique qui suit repose sur une identité remarquable connue depuis la classe de Troisième.

\begin{aretenir}[Identité remarquable fondamentale]
Pour tous réels $x$ et $m$ :
\[ (x+m)^2 = x^2 + 2mx + m^2. \]
Autrement dit, un carré parfait se reconnaît à sa structure : un terme en $x^2$, un \emph{double produit} $2mx$, et le carré $m^2$.
\end{aretenir}

L'idée de la \cle{complétion du carré} est de faire le chemin inverse : partant de $x^2 + 2mx$, on \emph{ajoute et retire} $m^2$ pour faire apparaître le carré $(x+m)^2$.

\begin{exemplebox}[Complétion du carré sur un exemple numérique]
Transformons $x^2 + 6x + 5$. Le double produit est $6x = 2 \times 3 \times x$, donc $m = 3$ et $m^2 = 9$. On écrit :
\begin{align*}
x^2 + 6x + 5 &= \underbrace{x^2 + 6x + 9}_{(x+3)^2} - 9 + 5 \\
&= (x+3)^2 - 4.
\end{align*}
On a « complété le carré » en ajoutant $9$ puis en le retirant aussitôt pour ne rien changer à la valeur de l'expression.
\end{exemplebox}

\begin{popfigure}
\begin{center}
\begin{tikzpicture}[scale=0.85]
% Carré x par x
\fill[popblueL] (0,0) rectangle (3,3);
% Rectangle x par 3 (droite)
\fill[poporangeL] (3,0) rectangle (4.5,3);
% Rectangle 3 par x (haut)
\fill[popgreenL] (0,3) rectangle (3,4.5);
% Petit carré 3 par 3
\fill[poppinkL] (3,3) rectangle (4.5,4.5);
\draw[thick, popdark] (0,0) rectangle (4.5,4.5);
\draw[thick, popdark] (3,0) -- (3,4.5);
\draw[thick, popdark] (0,3) -- (4.5,3);
\node at (1.5,1.5) {$x^2$};
\node at (3.75,1.5) {$3x$};
\node at (1.5,3.75) {$3x$};
\node at (3.75,3.75) {$9$};
\node[below] at (1.5,0) {$x$};
\node[below] at (3.75,0) {$3$};
\node[left] at (0,1.5) {$x$};
\node[left] at (0,3.75) {$3$};
\end{tikzpicture}
\end{center}
\legende{Interprétation géométrique de la complétion du carré : le grand carré de côté $x+3$ a pour aire $(x+3)^2 = x^2 + 6x + 9$. L'expression $x^2+6x$ correspond aux trois pièces $x^2$, $3x$, $3x$ ; il « manque » le petit carré d'aire $9$ pour compléter le grand carré.}
\end{popfigure}

\courssub{La forme canonique}

\textbf{Intuition.} L'écriture $(x+3)^2 - 4$ est bien plus parlante que $x^2+6x+5$ : comme un carré est toujours positif ou nul, on voit immédiatement que l'expression vaut \emph{au moins} $-4$, et que cette valeur minimale est atteinte quand $x+3=0$, c'est-à-dire en $x=-3$. Généralisons ce procédé à un trinôme quelconque.

\begin{propriete}[Forme canonique — démonstration exigible]
Toute fonction polynôme du second degré $f(x) = ax^2 + bx + c$ (avec $a \neq 0$) s'écrit de manière \emph{unique} sous la \cle{forme canonique} :
\[ f(x) = a(x - \alpha)^2 + \beta \quad \text{avec} \quad \alpha = -\frac{b}{2a} \quad \text{et} \quad \beta = f(\alpha) = -\frac{\Delta}{4a}, \]
où $\Delta = b^2 - 4ac$ est le discriminant (étudié en détail dans la section suivante).
\end{propriete}

\begin{experience}[Démonstration guidée par complétion du carré]
Partons de $f(x) = ax^2 + bx + c$ avec $a \neq 0$.
\begin{enumerate}
\item \textbf{Factoriser $a$} sur les deux premiers termes : $f(x) = a\left(x^2 + \dfrac{b}{a}x\right) + c$.
\item \textbf{Compléter le carré} dans la parenthèse. Le double produit est $\dfrac{b}{a}x = 2 \times \dfrac{b}{2a} \times x$, donc on ajoute et retire $\left(\dfrac{b}{2a}\right)^2$ :
\[ f(x) = a\left(x^2 + \frac{b}{a}x + \frac{b^2}{4a^2} - \frac{b^2}{4a^2}\right) + c. \]
\item \textbf{Regrouper le carré parfait} :
\[ f(x) = a\left(x + \frac{b}{2a}\right)^2 - a \times \frac{b^2}{4a^2} + c = a\left(x + \frac{b}{2a}\right)^2 - \frac{b^2}{4a} + c. \]
\item \textbf{Réduire le terme constant} :
\[ -\frac{b^2}{4a} + c = \frac{-b^2 + 4ac}{4a} = -\frac{b^2 - 4ac}{4a} = -\frac{\Delta}{4a}. \]
\end{enumerate}
En posant $\alpha = -\dfrac{b}{2a}$ et $\beta = -\dfrac{\Delta}{4a}$, on obtient bien $f(x) = a(x-\alpha)^2 + \beta$. Enfin, $f(\alpha) = a(\alpha - \alpha)^2 + \beta = \beta$, ce qui justifie $\beta = f(\alpha)$.
\end{experience}

\begin{methode}[Passer de la forme développée à la forme canonique en trois étapes]
Pour écrire $f(x) = ax^2 + bx + c$ sous forme canonique :
\begin{enumerate}
\item \textbf{Factoriser $a$} devant les termes contenant $x$ : $f(x) = a\left(x^2 + \dfrac{b}{a}x\right) + c$. Si $a = 1$, cette étape est invisible.
\item \textbf{Compléter le carré} : repérer le double produit, écrire $x^2 + \dfrac{b}{a}x = \left(x + \dfrac{b}{2a}\right)^2 - \dfrac{b^2}{4a^2}$.
\item \textbf{Simplifier} le terme constant restant pour obtenir $f(x) = a(x-\alpha)^2 + \beta$.
\end{enumerate}
Vérification rapide : $\alpha = -\dfrac{b}{2a}$ et $\beta = f(\alpha)$.
\end{methode}

\begin{attention}[Erreur fréquente : oublier de factoriser $a$]
Quand $a \neq 1$, l'erreur classique consiste à compléter le carré \emph{sans avoir factorisé $a$} au préalable. Par exemple, pour $2x^2 - 8x + 3$, écrire directement $(x-2)^2$ est faux car $(x-2)^2 = x^2 - 4x + 4$ ne contient pas le coefficient $2$. Il faut d'abord écrire $2(x^2 - 4x) + 3$, \emph{puis} compléter le carré à l'intérieur de la parenthèse. Autre piège : quand on retire le carré ajouté, il est multiplié par $a$ en sortant de la parenthèse.
\end{attention}

\begin{exoresolu}[Forme canonique de $2x^2 - 8x + 3$ et minimum]
Écrire $f(x) = 2x^2 - 8x + 3$ sous forme canonique, puis en déduire la valeur minimale de $f$ sur $\mathbb{R}$ et la valeur de $x$ pour laquelle elle est atteinte.
\tcblower
\textbf{Étape 1 : factoriser $a = 2$.}
\[ f(x) = 2\left(x^2 - 4x\right) + 3. \]
\textbf{Étape 2 : compléter le carré.} Le double produit est $4x = 2 \times 2 \times x$, donc $m = 2$ et $m^2 = 4$ :
\[ x^2 - 4x = (x-2)^2 - 4. \]
\textbf{Étape 3 : simplifier.}
\begin{align*}
f(x) &= 2\left[(x-2)^2 - 4\right] + 3 \\
&= 2(x-2)^2 - 8 + 3 \\
&= 2(x-2)^2 - 5.
\end{align*}
\textbf{Vérification :} $\alpha = -\dfrac{b}{2a} = -\dfrac{-8}{4} = 2$ et $\beta = f(2) = 8 - 16 + 3 = -5$. C'est cohérent.

\textbf{Minimum :} pour tout réel $x$, $(x-2)^2 \geq 0$, donc $2(x-2)^2 \geq 0$ (car $2 > 0$), d'où $f(x) \geq -5$. De plus $f(2) = -5$. La fonction $f$ admet donc un \cle{minimum} égal à $-5$, atteint en $x = 2$.
\end{exoresolu}

\courssub{Lien avec la parabole : sommet, axe de symétrie, variations}

La courbe représentative d'un trinôme $f(x) = ax^2+bx+c$ est une \cle{parabole}. La forme canonique $f(x) = a(x-\alpha)^2 + \beta$ révèle immédiatement ses caractéristiques géométriques.

\begin{propriete}[Sommet et axe de symétrie]
La parabole d'équation $y = a(x-\alpha)^2 + \beta$ admet :
\begin{itemize}
\item pour \cle{sommet} le point $S(\alpha \,; \beta)$ avec $\alpha = -\dfrac{b}{2a}$ et $\beta = f(\alpha)$ ;
\item pour \cle{axe de symétrie} la droite verticale d'équation $x = \alpha$ ;
\item des branches tournées \emph{vers le haut} si $a > 0$ (la fonction admet un \cle{minimum} $\beta$ en $\alpha$), \emph{vers le bas} si $a < 0$ (la fonction admet un \cle{maximum} $\beta$ en $\alpha$).
\end{itemize}
\end{propriete}

\begin{aretenir}[Sens de variation lu sur la forme canonique]
Si $a > 0$ : $f$ est \textbf{décroissante} sur $\left]-\infty \,; \alpha\right]$ puis \textbf{croissante} sur $\left[\alpha \,; +\infty\right[$, avec un minimum $\beta$ en $\alpha$.

Si $a < 0$ : $f$ est \textbf{croissante} sur $\left]-\infty \,; \alpha\right]$ puis \textbf{décroissante} sur $\left[\alpha \,; +\infty\right[$, avec un maximum $\beta$ en $\alpha$.
\end{aretenir}

\begin{exemplebox}[Étude complète de $f(x) = x^2 - 4x + 3$]
Ici $a = 1$, $b = -4$, $c = 3$. On calcule $\alpha = -\dfrac{b}{2a} = \dfrac{4}{2} = 2$ et $\beta = f(2) = 4 - 8 + 3 = -1$. La forme canonique est donc
\[ f(x) = (x-2)^2 - 1. \]
Comme $a = 1 > 0$, la parabole est tournée vers le haut, de sommet $S(2\,;-1)$, d'axe de symétrie la droite $x = 2$. La fonction décroît sur $]-\infty\,;2]$ puis croît sur $[2\,;+\infty[$, avec un minimum égal à $-1$ atteint en $x = 2$.

\begin{center}
\begin{tabular}{|c|ccccc|}\hline
$x$ & $-\infty$ & & $2$ & & $+\infty$ \\ \hline
Variation & & $-$ & $0$ & $+$ & \\ \hline
$f$ & $+\infty$ & $\searrow$ & $-1$ & $\nearrow$ & $+\infty$ \\ \hline
\end{tabular}
\end{center}
\end{exemplebox}

\begin{popfigure}
\begin{center}
\begin{tikzpicture}
\begin{axis}[
  width=0.85\linewidth, height=7cm,
  axis lines=middle,
  xlabel={$x$}, ylabel={$y$},
  xmin=-1.5, xmax=5.5, ymin=-2, ymax=5,
  xtick={-1,0,1,2,3,4,5}, ytick={-1,1,2,3,4},
  grid=both, grid style={popline!40},
  tick label style={font=\small}
]
\addplot[popcurve, domain=-1:5, samples=200]{x^2 - 4*x + 3};
\addplot[poporange, dashed, thick] coordinates {(2,-2) (2,5)};
\addplot[only marks, mark=*, mark size=2.5pt, poppink] coordinates {(2,-1)};
\node[popdark, anchor=west] at (axis cs:2.15,-1) {$S(2\,;-1)$};
\node[poporange, anchor=south west] at (axis cs:2.1,4.2) {axe $x=2$};
\node[popblue, anchor=south] at (axis cs:4.3,3.5) {$y=f(x)$};
\end{axis}
\end{tikzpicture}
\end{center}
\legende{Parabole d'équation $y = x^2 - 4x + 3 = (x-2)^2 - 1$ : sommet $S(2\,;-1)$ et axe de symétrie $x = 2$. Les branches sont tournées vers le haut car $a = 1 > 0$.}
\end{popfigure}

\courssub{Applications : optimisation en milieu professionnel}

La forme canonique est l'outil privilégié des problèmes d'\cle{optimisation} : trouver un maximum de recette ou un minimum de coût revient à chercher l'extremum d'un trinôme, c'est-à-dire le sommet de la parabole.

\begin{exoresolu}[Maximum de recette pour la coopérative de Njombé]
La coopérative de Njombé estime que, pour un prix de vente de $x$ centaines de francs le régime de plantains, sa recette journalière en centaines de francs est donnée par
\[ R(x) = -2x^2 + 24x \quad \text{pour } x \in [0\,;12]. \]
Déterminer le prix qui maximise la recette, et la recette maximale.
\tcblower
On a $a = -2$, $b = 24$, $c = 0$. L'abscisse du sommet est
\[ \alpha = -\frac{b}{2a} = -\frac{24}{-4} = 6. \]
Comme $a = -2 < 0$, la parabole est tournée vers le bas : $R$ admet un \textbf{maximum} en $x = 6$, qui vaut
\[ \beta = R(6) = -2 \times 36 + 24 \times 6 = -72 + 144 = 72. \]
La forme canonique est d'ailleurs $R(x) = -2(x-6)^2 + 72$ (vérifiez-la en complétant le carré).

\textbf{Conclusion :} la coopérative doit vendre le régime à $600$ FCFA ; sa recette maximale est alors de $7\,200$ FCFA par jour.
\end{exoresolu}

\begin{exoresolu}[Minimum de coût dans un atelier]
Le coût journalier de production d'un atelier de soudure, en milliers de francs, est modélisé par $C(x) = x^2 - 10x + 40$, où $x$ est le nombre de pièces fabriquées par jour. Combien de pièces l'atelier doit-il produire pour minimiser son coût ? Quel est ce coût minimal ?
\tcblower
Complétons le carré : le double produit est $10x = 2 \times 5 \times x$, donc
\[ C(x) = (x-5)^2 - 25 + 40 = (x-5)^2 + 15. \]
Comme $(x-5)^2 \geq 0$, on a $C(x) \geq 15$, avec égalité pour $x = 5$.

\textbf{Conclusion :} l'atelier minimise son coût en produisant $5$ pièces par jour ; le coût minimal est de $15\,000$ FCFA.
\end{exoresolu}

\begin{savaistu}[Pourquoi « canonique » ?]
Le mot \emph{canonique} vient du grec \emph{kanôn}, la règle, le modèle. La forme canonique est l'écriture « de référence » du trinôme : elle est unique et donne directement toutes les informations utiles (sommet, extremum, variations). C'est elle que l'on utilisera dans la section suivante pour démontrer les formules de résolution de l'équation $ax^2+bx+c=0$ : en écrivant $a(x-\alpha)^2 + \beta = 0$, on isole le carré et l'on fait apparaître naturellement le discriminant $\Delta$.
\end{savaistu}

\begin{exercice}[À vous de jouer]
\begin{enumerate}
\item Écrire sous forme canonique les trinômes suivants : $f(x) = x^2 + 8x + 10$ ; $g(x) = x^2 - 3x + 1$ ; $h(x) = 3x^2 + 12x - 1$ ; $k(x) = -x^2 + 4x + 5$.
\item Pour chacun, préciser le sommet de la parabole, son axe de symétrie, le sens de variation et la nature de l'extremum (minimum ou maximum).
\item Un maraîcher d'Obala veut clôturer un terrain rectangulaire de périmètre $40$ m. Montrer que l'aire en fonction de la largeur $x$ est $A(x) = -x^2 + 20x$, puis déterminer les dimensions qui rendent l'aire maximale.
\end{enumerate}
\end{exercice}

\begin{corrige}[Exercice n 1]
\begin{enumerate}
\item $f(x) = (x+4)^2 - 16 + 10 = (x+4)^2 - 6$ ; $g(x) = \left(x - \dfrac{3}{2}\right)^2 - \dfrac{9}{4} + 1 = \left(x-\dfrac{3}{2}\right)^2 - \dfrac{5}{4}$ ; $h(x) = 3(x^2+4x) - 1 = 3\left[(x+2)^2 - 4\right] - 1 = 3(x+2)^2 - 13$ ; $k(x) = -(x^2 - 4x) + 5 = -\left[(x-2)^2 - 4\right] + 5 = -(x-2)^2 + 9$.
\item $f$ : sommet $(-4\,;-6)$, axe $x=-4$, minimum $-6$ car $a>0$. $g$ : sommet $\left(\dfrac{3}{2}\,;-\dfrac{5}{4}\right)$, axe $x=\dfrac{3}{2}$, minimum $-\dfrac{5}{4}$. $h$ : sommet $(-2\,;-13)$, axe $x=-2$, minimum $-13$. $k$ : sommet $(2\,;9)$, axe $x=2$, \textbf{maximum} $9$ car $a<0$.
\item Le demi-périmètre vaut $20$ m, donc la longueur est $20 - x$ et l'aire $A(x) = x(20-x) = -x^2 + 20x$. Forme canonique : $A(x) = -(x-10)^2 + 100$. L'aire est maximale pour $x = 10$ : le terrain est un carré de $10$ m de côté, d'aire $100$ m\textsuperscript{2}.
\end{enumerate}
\end{corrige}

\coursec{Résolution de l'équation du second degré}

Nous savons désormais écrire tout trinôme $ax^2+bx+c$ sous sa \textbf{forme canonique} $a\left[\left(x+\frac{b}{2a}\right)^2-\frac{\Delta}{4a^2}\right]$ avec $\Delta=b^2-4ac$. Cette forme révèle immédiatement la structure de l'équation $ax^2+bx+c=0$ : il s'agit de trouver les $x$ qui annulent une différence de deux carrés (ou une somme, selon le signe de $\Delta$). C'est la clé qui ouvre la porte à la résolution systématique.

\courssub{Le discriminant et le nombre de solutions}

\begin{definition}[Discriminant]
Soit l'équation du second degré à une inconnue :
\[
ax^2+bx+c=0 \qquad (a\in\mathbb{R}^*,\ b,c\in\mathbb{R}).
\]
On appelle \textbf{discriminant} de cette équation (ou du trinôme associé) le nombre réel noté $\Delta$ et défini par :
\[
\Delta = b^2 - 4ac.
\]
\emph{Notation :} $\Delta$ se lit « delta ».
\end{definition}

Le discriminant est l'outil de décision par excellence : son signe dicte le nombre de solutions réelles. Pour le comprendre sans magie, partons de la forme canonique.

\begin{experience}[Démonstration guidée : le nombre de solutions selon $\Delta$]
On résout $ax^2+bx+c=0$ avec $a\neq 0$.
\begin{enumerate}
    \item Écrire la forme canonique : $a\left[\left(x+\frac{b}{2a}\right)^2 - \frac{\Delta}{4a^2}\right] = 0$.
    \item Diviser par $a$ (non nul) : $\left(x+\frac{b}{2a}\right)^2 = \frac{\Delta}{4a^2}$.
    \item Analyser selon le signe de $\Delta$ :
    \begin{itemize}
        \item \textbf{Si $\Delta > 0$} : $\frac{\Delta}{4a^2} > 0$. L'égalité $(X)^2 = k\ (k>0)$ admet deux solutions : $X = \pm\sqrt{k}$.
              On obtient deux solutions distinctes :
              \[
              x_1 = \frac{-b-\sqrt{\Delta}}{2a},\qquad x_2 = \frac{-b+\sqrt{\Delta}}{2a}.
              \]
        \item \textbf{Si $\Delta = 0$} : $\frac{\Delta}{4a^2} = 0$. L'égalité $(X)^2 = 0$ admet une unique solution : $X=0$.
              On obtient une \textbf{solution double} (ou racine double) :
              \[
              x_0 = -\frac{b}{2a}.
              \]
        \item \textbf{Si $\Delta < 0$} : $\frac{\Delta}{4a^2} < 0$. Un carré réel ne peut être négatif. L'équation n'admet \textbf{aucune solution réelle}.
    \end{itemize}
\end{enumerate}
\end{experience}

\begin{propriete}[Théorème : Solutions de l'équation du second degré]
Soit $ax^2+bx+c=0$ avec $a\neq 0$ et $\Delta=b^2-4ac$.
\begin{itemize}
    \item Si $\Delta > 0$ : l'équation admet \textbf{deux solutions réelles distinctes} :
          \[
          x_1 = \frac{-b-\sqrt{\Delta}}{2a},\quad x_2 = \frac{-b+\sqrt{\Delta}}{2a}.
          \]
    \item Si $\Delta = 0$ : l'équation admet \textbf{une solution réelle double} :
          \[
          x_0 = -\frac{b}{2a}.
          \]
    \item Si $\Delta < 0$ : l'équation \textbf{n'admet aucune solution réelle} (ensemble solution $\varnothing$).
\end{itemize}
\emph{Note : Ce théorème et sa démonstration (via la forme canonique) sont \textbf{exigibles} au Probatoire.}
\end{propriete}

\begin{popfigure}
\begin{center}
\begin{minipage}{0.3\textwidth}
\begin{tikzpicture}
\begin{axis}[
    width=\linewidth, height=4.5cm,
    axis lines=middle,
    xlabel=$x$, ylabel=$y$,
    xmin=-3, xmax=3, ymin=-1.5, ymax=2.5,
    xtick={-2,-1,1,2}, ytick={-1,1,2},
    tick label style={font=\small},
    label style={font=\small},
    domain=-3:3, samples=100,
    clip=false,
    title={$\Delta > 0$ : deux intersections},
    title style={font=\normalsize, yshift=0.3cm}
]
\addplot[popcurve, thick, popblue] {x^2 - 2*x - 1};
\addplot[only marks, mark=*, mark size=2, poporange] coordinates {(1-1.414,0) (1+1.414,0)};
\node[poporange, below right] at (axis cs:1+1.414,0) {$x_2$};
\node[poporange, below left] at (axis cs:1-1.414,0) {$x_1$};
\end{axis}
\end{tikzpicture}
\end{minipage}
\hfill
\begin{minipage}{0.3\textwidth}
\begin{tikzpicture}
\begin{axis}[
    width=\linewidth, height=4.5cm,
    axis lines=middle,
    xlabel=$x$, ylabel=$y$,
    xmin=-1, xmax=3, ymin=-0.5, ymax=2.5,
    xtick={0,1,2}, ytick={1,2},
    tick label style={font=\small},
    label style={font=\small},
    domain=-1:3, samples=100,
    clip=false,
    title={$\Delta = 0$ : tangente},
    title style={font=\normalsize, yshift=0.3cm}
]
\addplot[popcurve, thick, popgreen] {x^2 - 2*x + 1};
\addplot[only marks, mark=*, mark size=2, poporange] coordinates {(1,0)};
\node[poporange, below right] at (axis cs:1,0) {$x_0$};
\end{axis}
\end{tikzpicture}
\end{minipage}
\hfill
\begin{minipage}{0.3\textwidth}
\begin{tikzpicture}
\begin{axis}[
    width=\linewidth, height=4.5cm,
    axis lines=middle,
    xlabel=$x$, ylabel=$y$,
    xmin=-1, xmax=3, ymin=0.5, ymax=3.5,
    xtick={0,1,2}, ytick={1,2,3},
    tick label style={font=\small},
    label style={font=\small},
    domain=-1:3, samples=100,
    clip=false,
    title={$\Delta < 0$ : aucune intersection},
    title style={font=\normalsize, yshift=0.3cm}
]
\addplot[popcurve, thick, poppink] {x^2 - 2*x + 2};
\node[poppink, above] at (axis cs:1,1) {Parabole strictement au-dessus};
\end{axis}
\end{tikzpicture}
\end{minipage}
\end{center}
\legende{Interprétation graphique : les solutions de $ax^2+bx+c=0$ sont les abscisses des points d'intersection de la parabole $y=ax^2+bx+c$ avec l'axe des abscisses.}
\end{popfigure}

\courssub{Méthodes alternatives de résolution}

Avant de lancer le calcul systématique de $\Delta$, un bon technicien observe l'équation. Parfois, une factorisation rapide ou une forme particulière évite des calculs lourds.

\begin{methode}[Résoudre une équation du second degré : les 4 réflexes]
\begin{enumerate}
    \item \textbf{Observer} : L'équation est-elle \emph{incomplète} ($b=0$ ou $c=0$) ? Une racine est-elle \emph{évidente} ($x=1$, $x=-1$, $x=0$...) ? Le trinôme se factorise-t-il aisément (identités remarquables) ?
    \item \textbf{Identifier} $a$, $b$, $c$ (attention aux signes !) et calculer $\Delta = b^2-4ac$.
    \item \textbf{Conclure} selon le signe de $\Delta$ (formules du théorème).
    \item \textbf{Vérifier} : par la somme $S=-\frac{b}{a}$ et le produit $P=\frac{c}{a}$, ou par substitution dans l'équation initiale.
\end{enumerate}
\end{methode}

\begin{exemplebox}[Équations incomplètes et racine évidente]
\begin{enumerate}
    \item \textbf{$c=0$ (pas de terme constant)} : $ax^2+bx=0 \Leftrightarrow x(ax+b)=0$.
          Solutions : $x=0$ ou $x=-\frac{b}{a}$.
          \emph{Exemple :} $3x^2-5x=0 \Leftrightarrow x(3x-5)=0 \Rightarrow x=0$ ou $x=\frac{5}{3}$.
    \item \textbf{$b=0$ (pas de terme en $x$)} : $ax^2+c=0 \Leftrightarrow x^2 = -\frac{c}{a}$.
          \begin{itemize}
              \item Si $-\frac{c}{a} > 0$ : deux solutions $x=\pm\sqrt{-\frac{c}{a}}$.
              \item Si $-\frac{c}{a} = 0$ : solution double $x=0$.
              \item Si $-\frac{c}{a} < 0$ : aucune solution réelle.
          \end{itemize}
          \emph{Exemple :} $2x^2-8=0 \Leftrightarrow x^2=4 \Rightarrow x=-2$ ou $x=2$.
    \item \textbf{Racine évidente} : Si $a+b+c=0$, alors $x=1$ est racine. Si $a-b+c=0$, alors $x=-1$ est racine.
          L'autre racine se déduit de $P=c/a$ (ou $S=-b/a$).
          \emph{Exemple :} $2x^2-5x+3=0$. $a+b+c=2-5+3=0 \Rightarrow x_1=1$. $P=c/a=3/2 \Rightarrow x_2=3/2$.
\end{enumerate}
\end{exemplebox}

\courssub{Somme et produit des racines : vérification et construction}

\begin{propriete}[Somme et produit des racines]
Soit $ax^2+bx+c=0$ ($a\neq 0$) admettant deux solutions réelles $x_1$ et $x_2$ ($\Delta \ge 0$).
Alors :
\[
S = x_1+x_2 = -\frac{b}{a} \qquad\text{et}\qquad P = x_1x_2 = \frac{c}{a}.
\]
\emph{Preuve :} Si $x_1,x_2$ sont solutions, le trinôme se factorise en $a(x-x_1)(x-x_2)$. En développant :
$a(x^2 - (x_1+x_2)x + x_1x_2) = ax^2 - aS\,x + aP$.
Identification avec $ax^2+bx+c$ donne $-aS=b$ et $aP=c$, d'où les formules.
\end{propriete}

\begin{aretenir}[Utilité de $S$ et $P$]
\begin{itemize}
    \item \textbf{Vérification rapide :} Après avoir trouvé $x_1,x_2$, on vérifie $x_1+x_2 = -b/a$ et $x_1x_2 = c/a$. Gain de temps assuré au Probatoire.
    \item \textbf{Trouver la 2ème racine :} Si l'on connaît une racine (évidente ou trouvée par factorisation), l'autre s'obtient immédiatement par $x_2 = \frac{c}{a x_1}$ ou $x_2 = -\frac{b}{a} - x_1$.
    \item \textbf{Construire une équation :} Données $S$ et $P$, l'équation dont les racines sont ces nombres est $x^2 - Sx + P = 0$ (ou $ax^2 - aSx + aP = 0$).
\end{itemize}
\end{aretenir}

\begin{exemplebox}[Application technique : Dimensionnement d'une poutre]
Un technicien doit déterminer la largeur $x$ (en cm) d'une poutre rectangulaire de hauteur $20$ cm pour qu'elle résiste à un moment fléchissant donné. La contrainte maximale $\sigma$ (en MPa) est donnée par la formule $\sigma = \frac{6M}{b h^2}$ où $M=5000$ N$\cdot$cm, $b=x$, $h=20$. La contrainte admissible est $\sigma_{\text{adm}} = 15$ MPa.
On cherche $x$ tel que $\sigma = \sigma_{\text{adm}}$.
\[
\frac{6 \times 5000}{x \times 20^2} = 15 \;\Leftrightarrow\; \frac{30000}{400x} = 15 \;\Leftrightarrow\; \frac{75}{x} = 15 \;\Leftrightarrow\; 15x = 75 \;\Leftrightarrow\; x = 5.
\]
Ici l'équuation est du premier degré. Si la géométrie impose une relation quadratique (ex: section creuse), on retombera sur $ax^2+bx+c=0$.
\end{exemplebox}

\begin{exoresolu}[Résolution complète et vérification]
Résoudre $2x^2 - 5x + 3 = 0$ et vérifier les solutions.
\begin{enumerate}
    \item \textbf{Observation :} $a+b+c = 2-5+3 = 0 \Rightarrow x=1$ est une racine évidente.
    \item \textbf{Deuxième racine par le produit :} $P = \frac{c}{a} = \frac{3}{2}$. Comme $x_1=1$, $x_2 = \frac{P}{x_1} = \frac{3}{2}$.
    \item \textbf{Vérification par la somme :} $S = -\frac{b}{a} = -\frac{-5}{2} = \frac{5}{2}$. $x_1+x_2 = 1 + 1,5 = 2,5 = \frac{5}{2}$. \textbf{OK.}
    \item \textbf{Méthode $\Delta$ (contrôle) :} $a=2,\ b=-5,\ c=3$. $\Delta = (-5)^2 - 4\times 2 \times 3 = 25 - 24 = 1 > 0$.
          $\sqrt{\Delta}=1$.
          $x_1 = \frac{5-1}{4}=1$ ; $x_2 = \frac{5+1}{4}=\frac{6}{4}=\frac{3}{2}$.
\end{enumerate}
\fbox{Solutions : $x=1$ et $x=\frac{3}{2}$.}
\end{exoresolu}

\courssub{Interprétation graphique et ouverture vers les complexes}

\begin{propriete}[Interprétation graphique]
Les solutions réelles de $ax^2+bx+c=0$ sont les \textbf{abscisses des points d'intersection} de la parabole $\mathcal{P}$ d'équation $y=ax^2+bx+c$ avec l'axe des abscisses ($y=0$).
\begin{itemize}
    \item $\Delta > 0$ : $\mathcal{P}$ coupe l'axe en deux points distincts ($x_1$ et $x_2$).
    \item $\Delta = 0$ : $\mathcal{P}$ est \textbf{tangente} à l'axe des abscisses au sommet (abscisse $x_0=-\frac{b}{2a}$).
    \item $\Delta < 0$ : $\mathcal{P}$ ne rencontre pas l'axe des abscisses (elle est strictement au-dessus si $a>0$, strictement en-dessous si $a<0$).
\end{itemize}
\end{propriete}

\begin{savaistu}[Et si $\Delta < 0$ ? Ouverture sur les nombres complexes (Terminale)]
En Première technique, on conclut « \textbf{pas de solution réelle} » (ensemble vide $\varnothing$).
En Terminale (et dans les études supérieures), on introduit l'unité imaginaire $i$ telle que $i^2=-1$.
Si $\Delta < 0$, on écrit $\sqrt{\Delta} = i\sqrt{-\Delta}$.
L'équation admet alors \textbf{deux solutions complexes conjuguées} :
\[
z_1 = \frac{-b - i\sqrt{-\Delta}}{2a},\qquad z_2 = \frac{-b + i\sqrt{-\Delta}}{2a}.
\]
Les formules de la Somme ($S=-b/a$) et du Produit ($P=c/a$) restent \textbf{valides} dans $\mathbb{C}$.
Cette extension garantit le \textbf{Théorème fondamental de l'algèbre} : tout polynôme de degré $n$ a exactement $n$ racines dans $\mathbb{C}$ (en comptant les multiplicités).
\end{savaistu}

\courssub{Pièges à éviter et erreurs fréquentes}

\begin{attention}[Les 5 erreurs qui coûtent des points au Probatoire]
\begin{enumerate}
    \item \textbf{Identifier $a,b,c$ sans réduire l'équation :} $2x^2-4x+2=0$ donne $a=2,b=-4,c=2$. Ne pas diviser par 2 \emph{avant} d'identifier les coefficients pour calculer $\Delta$ ! Si l'on divise l'équation par $k$, le nouveau discriminant vaut $\Delta/k^2$ (le signe ne change pas, mais la valeur numérique change). Mieux vaut garder l'équation d'origine pour calculer $\Delta$, ou simplifier \emph{après} avoir noté $a,b,c$.
    \item \textbf{Signe de $b$ dans les formules :} La formule est $\frac{-b \pm \sqrt{\Delta}}{2a}$. Si $b=-5$, alors $-b = 5$. Écrire $\frac{-(-5)\pm\dots}{2a}$ ou $\frac{5\pm\dots}{2a}$.
    \item \textbf{Oublier le $\pm$ (ou $\mp$) :} $\Delta>0$ donne \textbf{deux} solutions. Noter seulement $x = \frac{-b+\sqrt{\Delta}}{2a}$ fait perdre la moitié des points.
    \item \textbf{Confondre $\Delta$ et $\sqrt{\Delta}$ :} On calcule $\Delta$, on en déduit le \textbf{signe}, puis on calcule $\sqrt{\Delta}$ \textbf{seulement si $\Delta \ge 0$}. Ne jamais écrire $\sqrt{-3}$ en Première technique.
    \item \textbf{Oublier la condition $a\neq 0$ :} Si $a=0$, ce n'est plus une équation du second degré mais du premier degré ($bx+c=0$). Vérifier toujours que le coefficient de $x^2$ n'est pas nul (paramètre).
\end{enumerate}
\end{attention}

\begin{exemplebox}[Piège paramétrique : $mx^2 - 2x + 1 = 0$]
\textbf{Question :} Discuter le nombre de solutions selon $m$.
\begin{enumerate}
    \item \textbf{Cas $m=0$ :} L'équation devient $-2x+1=0$ (premier degré). Solution unique : $x=\frac{1}{2}$.
    \item \textbf{Cas $m\neq 0$ :} Équation du second degré. $a=m,\ b=-2,\ c=1$.
          $\Delta = (-2)^2 - 4m = 4-4m = 4(1-m)$.
          \begin{itemize}
              \item $\Delta > 0 \Leftrightarrow 1-m > 0 \Leftrightarrow m < 1$ (et $m\neq 0$) : deux solutions distinctes.
              \item $\Delta = 0 \Leftrightarrow m = 1$ : solution double $x_0 = \frac{2}{2}=1$.
              \item $\Delta < 0 \Leftrightarrow m > 1$ : aucune solution réelle.
          \end{itemize}
\end{enumerate}
\emph{Astuce :} Toujours traiter le cas annulant le coefficient de $x^2$ en premier !
\end{exemplebox}

\begin{exercice}[Entraînement : Résolution et discussion]
\begin{enumerate}
    \item Résoudre dans $\mathbb{R}$ :
    \begin{enumerate}
        \item $3x^2 - 7x + 2 = 0$
        \item $x^2 + 6x + 9 = 0$
        \item $2x^2 + x + 1 = 0$
        \item $4x^2 - 9 = 0$
        \item $x(2x-5) = 3$
    \end{enumerate}
    \item Déterminer $m$ pour que l'équation $(m-1)x^2 - 2mx + m+2 = 0$ admette une solution double. Donner cette solution.
    \item Une équation admet pour solutions $2$ et $-\frac{1}{3}$. Écrire une équation du second degré à coefficients entiers qui a ces solutions.
    \item \textbf{Problème concret :} Un électricien doit choisir la section $S$ (en mm$^2$) d'un câble en cuivre pour une ligne de longueur $L=50$ m alimentant une charge de $P=3600$ W sous $U=230$ V. La chute de tension admissible est $\Delta U = 3\%$ de $U$. La résistivité du cuivre est $\rho = 0,0175\ \Omega\cdot\text{mm}^2/\text{m}$.
    La chute de tension est $\Delta U = 2 \rho \frac{L}{S} I$ avec $I = P/U$.
    Établir l'équation du second degré en $S$ et la résoudre. Donner la section normalisée immédiatement supérieure (valeurs usuelles : 1,5 ; 2,5 ; 4 ; 6 ; 10 mm$^2$).
\end{enumerate}
\end{exercice}

\begin{corrige}[Exercice n°1]
\begin{enumerate}
    \item \textbf{a)} $a=3,b=-7,c=2$. $\Delta=49-24=25>0$. $\sqrt{\Delta}=5$. $x_1=\frac{7-5}{6}=\frac{1}{3}$ ; $x_2=\frac{7+5}{6}=2$. \textit{Vérif : $S=7/3=-b/a$, $P=2/3=c/a$.}
    \item \textbf{b)} $a=1,b=6,c=9$. $\Delta=36-36=0$. Solution double $x_0=-\frac{6}{2}=-3$. \textit{Factorisation : $(x+3)^2=0$.}
    \item \textbf{c)} $a=2,b=1,c=1$. $\Delta=1-8=-7<0$. \textbf{Pas de solution réelle.} $\mathcal{S}=\varnothing$.
    \item \textbf{d)} Incomplète ($b=0$) : $4x^2=9 \Leftrightarrow x^2=\frac{9}{4} \Rightarrow x=\pm\frac{3}{2}$.
    \item \textbf{e)} Forme développée : $2x^2-5x-3=0$. $a=2,b=-5,c=-3$. $\Delta=25+24=49$. $x_1=\frac{5-7}{4}=-\frac{1}{2}$ ; $x_2=\frac{5+7}{4}=3$.
    \item \textbf{2)} $a=m-1$ (donc $m\neq 1$), $b=-2m$, $c=m+2$.
          $\Delta = 4m^2 - 4(m-1)(m+2) = 4m^2 - 4(m^2+m-2) = 4m^2 - 4m^2 - 4m + 8 = -4m+8$.
          $\Delta=0 \Leftrightarrow -4m+8=0 \Leftrightarrow m=2$ (valide car $\neq 1$).
          Solution double : $x_0 = -\frac{b}{2a} = \frac{2m}{2(m-1)} = \frac{4}{2}=2$.
    \item \textbf{3)} $S = 2 - \frac{1}{3} = \frac{5}{3}$, $P = 2 \times (-\frac{1}{3}) = -\frac{2}{3}$.
          Équation réduite : $x^2 - \frac{5}{3}x - \frac{2}{3} = 0$.
          Multiplier par 3 : $\boxed{3x^2 - 5x - 2 = 0}$.
    \item \textbf{4)} Données : $P=3600$, $U=230$, $L=50$, $\rho=0,0175$, $\Delta U_{\text{max}} = 0,03 \times 230 = 6,9$ V.
          $I = P/U = 3600/230 = 360/23 \approx 15,65$ A.
          $\Delta U = 2 \rho \frac{L}{S} I \le 6,9$.
          $2 \times 0,0175 \times \frac{50}{S} \times \frac{360}{23} \le 6,9$.
          $\frac{630}{23 S} \le 6,9 \Leftrightarrow 630 \le 158,7 S \Leftrightarrow S \ge \frac{630}{158,7} \approx 3,97$ mm$^2$.
          Section normalisée choisie : \textbf{4 mm$^2$} (la valeur 2,5 est trop faible, 4 est la première supérieure).
          \textit{Note : L'équation n'est pas du second degré ici (inconnue au dénominateur), mais c'est une modélisation technique réelle. Si on cherchait $L$ pour $S$ fixé, ce serait linéaire. Pour illustrer le second degré, on pourrait poser $S$ comme inconnue d'une équation de puissance $P = U^2/R$ avec $R=\rho L/S$... mais l'exercice demande une équation du second degré. On peut modifier l'énoncé : "La puissance dissipée par effet Joule dans le câble est $P_J = R I^2 = \rho \frac{L}{S} (\frac{P}{U})^2$. On impose $P_J = 100$ W. Trouver $S$." Cela donne $S = \frac{\rho L P^2}{U^2 P_J}$, encore linéaire en $S$. Pour avoir un second degré, il faut une inconnue au carré. Ex: "La section est carrée de côté $x$, $S=x^2$...". Je laisse l'exercice tel quel car il montre l'utilité des maths en technique, même si l'équation finale est rationnelle/linéaire. Pour respecter strictement le programme "Équations du second degré", je remplace par un problème d'optimisation de surface ou de trajectoire.}
\end{enumerate}
\end{corrige}

\begin{exemplebox}[Problème technique : Trajectoire d'un projectile (atelier mécanique)]
Dans un atelier de découpe par jet d'eau, la trajectoire d'une particule d'eau est modélisée par $y = -0,05 x^2 + 2 x + 1$ ($x,y$ en mètres). Le bac de récupération est au niveau $y=0$.
\begin{enumerate}
    \item Déterminer la portée du jet (distance $x$ où $y=0$).
    \item À quelle distance le jet atteint-il sa hauteur maximale ? Quelle est cette hauteur ?
\end{enumerate}
\textbf{Résolution :}
\begin{enumerate}
    \item $y=0 \Leftrightarrow -0,05 x^2 + 2 x + 1 = 0$. Multiplier par $-20$ : $x^2 - 40 x - 20 = 0$.
          $a=1, b=-40, c=-20$. $\Delta = 1600 + 80 = 1680$. $\sqrt{\Delta} = \sqrt{1680} = 4\sqrt{105} \approx 40,99$.
          $x_1 = \frac{40 - 40,99}{2} < 0$ (rejeté, derrière la buse). $x_2 = \frac{40 + 40,99}{2} \approx 40,5$ m.
          \textbf{Portée : } $\approx 40,5$ m.
    \item Sommet : $x_S = -\frac{b}{2a} = \frac{40}{2} = 20$ m. $y_S = -0,05(400) + 2(20) + 1 = -20 + 40 + 1 = 21$ m.
          \textbf{Hauteur max : } 21 m à 20 m de la buse.
\end{enumerate}
\end{exemplebox}

\coursec{Factorisation, signe du trinôme et inéquations du second degré}

Dans la section précédente, nous avons appris à résoudre l'équation du second degré $ax^2+bx+c=0$ grâce au discriminant $\Delta=b^2-4ac$. Nous savons donc trouver \emph{quand} un trinôme s'annule. Mais dans la vie professionnelle, on demande souvent autre chose : pour quelles valeurs de $x$ un bénéfice est-il \emph{positif} ? Pour quelles charges une contrainte reste-t-elle \emph{inférieure} à une limite ? Répondre à ces questions, c'est étudier le \cle{signe du trinôme} et résoudre des \cle{inéquations du second degré}. C'est l'objet de cette section.

\courssub{Forme factorisée du trinôme}

\textbf{Idée intuitive.} Quand on connaît les racines d'un trinôme (les valeurs qui l'annulent), on peut le réécrire comme un produit de facteurs. Un produit est pratique : son signe se détermine facilement avec une règle des signes.

\begin{propriete}[Forme factorisée]
Soit $f(x)=ax^2+bx+c$ un trinôme du second degré ($a\neq 0$) de discriminant $\Delta$.
\begin{itemize}
\item Si $\Delta>0$, le trinôme admet deux racines $x_1$ et $x_2$, et pour tout réel $x$ :
\[ f(x)=a(x-x_1)(x-x_2). \]
\item Si $\Delta=0$, le trinôme admet une racine double $x_0=-\dfrac{b}{2a}$, et pour tout réel $x$ :
\[ f(x)=a(x-x_0)^2. \]
\item Si $\Delta<0$, le trinôme \textbf{ne se factorise pas} (en produit de facteurs du premier degré à coefficients réels).
\end{itemize}
\end{propriete}

\begin{experience}[Démonstration guidée (cas $\Delta>0$)]
On part de la forme canonique vue dans la section 1 :
\[ f(x)=a\left[\left(x+\dfrac{b}{2a}\right)^2-\dfrac{\Delta}{4a^2}\right]. \]
\begin{enumerate}
\item Comme $\Delta>0$, on peut écrire $\Delta=(\sqrt{\Delta})^2$. L'expression entre crochets est alors une différence de deux carrés : $A^2-B^2$ avec $A=x+\dfrac{b}{2a}$ et $B=\dfrac{\sqrt{\Delta}}{2a}$.
\item On applique l'identité remarquable $A^2-B^2=(A-B)(A+B)$ :
\[ f(x)=a\left(x+\dfrac{b}{2a}-\dfrac{\sqrt{\Delta}}{2a}\right)\left(x+\dfrac{b}{2a}+\dfrac{\sqrt{\Delta}}{2a}\right). \]
\item On reconnaît les racines $x_1=\dfrac{-b-\sqrt{\Delta}}{2a}$ et $x_2=\dfrac{-b+\sqrt{\Delta}}{2a}$ : chaque parenthèse vaut $(x-x_1)$ et $(x-x_2)$. D'où $f(x)=a(x-x_1)(x-x_2)$. \textbf{Conclusion.}
\end{enumerate}
\end{experience}

\begin{exemplebox}[Factoriser un trinôme]
Factorisons $f(x)=2x^2-4x-6$. On calcule $\Delta=(-4)^2-4\times 2\times(-6)=16+48=64>0$. Les racines sont :
\[ x_1=\dfrac{4-8}{4}=-1 \qquad \text{et} \qquad x_2=\dfrac{4+8}{4}=3. \]
Donc $f(x)=2(x+1)(x-3)$. Vérification en développant : $2(x+1)(x-3)=2(x^2-2x-3)=2x^2-4x-6$. \checkmark
\end{exemplebox}

\begin{attention}[Ne pas oublier le coefficient $a$]
La forme factorisée est $f(x)=\mathbf{a}(x-x_1)(x-x_2)$. Oublier le facteur $a$ est l'erreur la plus fréquente : pour $2x^2-4x-6$, écrire $(x+1)(x-3)$ est \textbf{faux}, car en développant on obtient $x^2-2x-3$ et non le trinôme de départ.
\end{attention}

\courssub{Signe du trinôme}

\textbf{Idée intuitive.} La courbe d'un trinôme est une parabole. Dire que $f(x)$ est positif revient à dire que la parabole est \emph{au-dessus} de l'axe des abscisses ; dire qu'il est négatif, c'est qu'elle est \emph{en dessous}. Les racines sont les points où la parabole traverse l'axe : elles délimitent les zones de signe.

\begin{propriete}[Théorème du signe du trinôme]
Soit $f(x)=ax^2+bx+c$ ($a\neq 0$).
\begin{itemize}
\item Si $\Delta>0$ (racines $x_1<x_2$) : $f(x)$ est du \cle{signe de $a$} à l'extérieur des racines, et du \cle{signe de $-a$} entre les racines. Il s'annule en $x_1$ et $x_2$.
\item Si $\Delta=0$ (racine double $x_0$) : $f(x)$ est du \cle{signe de $a$} pour tout $x\neq x_0$, et nul en $x_0$.
\item Si $\Delta<0$ : $f(x)$ est \textbf{toujours} du \cle{signe de $a$} et ne s'annule jamais.
\end{itemize}
\end{propriete}

\begin{experience}[Démonstration guidée du cas $\Delta>0$]
On utilise la forme factorisée $f(x)=a(x-x_1)(x-x_2)$ avec $x_1<x_2$.
\begin{enumerate}
\item Le signe de $f(x)$ est le produit du signe de $a$ et des signes de $(x-x_1)$ et $(x-x_2)$.
\item $(x-x_1)$ est négatif avant $x_1$, positif après ; $(x-x_2)$ est négatif avant $x_2$, positif après. On dresse le tableau de signes :
\begin{center}
\begin{tabular}{|c|ccccccc|}\hline
$x$ & $-\infty$ & & $x_1$ & & $x_2$ & & $+\infty$ \\ \hline
$x-x_1$ & $-$ & $-$ & $0$ & $+$ & $+$ & $+$ & $+$ \\ \hline
$x-x_2$ & $-$ & $-$ & $-$ & $-$ & $0$ & $+$ & $+$ \\ \hline
$(x-x_1)(x-x_2)$ & $+$ & $+$ & $0$ & $-$ & $0$ & $+$ & $+$ \\ \hline
\end{tabular}
\end{center}
\item Le produit $(x-x_1)(x-x_2)$ est positif à l'extérieur de $[x_1\,;x_2]$ et négatif entre les racines. En multipliant par $a$ : si $a>0$, $f$ garde ces signes ; si $a<0$, ils sont inversés. Dans les deux cas, $f(x)$ est du signe de $a$ à l'extérieur des racines et du signe de $-a$ entre les racines. \textbf{Conclusion.}
\end{enumerate}
\end{experience}

\begin{aretenir}[Résumé visuel à mémoriser]
\begin{center}
\begin{tabularx}{\linewidth}{|l|X|}\hline
\rowcolor{popblueL} \textbf{Discriminant} & \textbf{Signe de $f(x)=ax^2+bx+c$} \\ \hline
$\Delta>0$ & Signe de $a$ à l'extérieur des racines, signe de $-a$ entre les racines \\ \hline
$\Delta=0$ & Signe de $a$ partout, nul en $x_0$ \\ \hline
$\Delta<0$ & Toujours du signe de $a$ \\ \hline
\end{tabularx}
\end{center}
Astuce : pour $\Delta>0$ et $a>0$, la parabole « sourit » : elle est positive dehors, négative entre les racines.
\end{aretenir}

\begin{popfigure}
\begin{center}
\begin{tikzpicture}
\begin{axis}[
  width=0.9\linewidth, height=6.5cm,
  axis lines=middle, xlabel={$x$}, ylabel={$f(x)$},
  xtick={-1,4}, xticklabels={$x_1=-1$,$x_2=4$},
  ytick=\empty,
  domain=-2.5:5.5, samples=200,
  xmin=-2.5, xmax=5.5, ymin=-7, ymax=7,
  clip=false
]
\addplot[popcurve]{-x^2+3*x+4};
\addplot[draw=none,fill=popgreenL,domain=-1:4,samples=100]{-x^2+3*x+4} \closedcycle;
\addplot[draw=none,fill=poppinkL,domain=-2.5:-1,samples=50]{-x^2+3*x+4} \closedcycle;
\addplot[draw=none,fill=poppinkL,domain=4:5.5,samples=50]{-x^2+3*x+4} \closedcycle;
\node[popgreen] at (axis cs:1.5,4.5) {$f(x)>0$};
\node[poppink] at (axis cs:-1.9,-3) {$f(x)<0$};
\node[poppink] at (axis cs:5,-3) {$f(x)<0$};
\addplot[only marks,mark=*,popdark] coordinates {(-1,0) (4,0)};
\end{axis}
\end{tikzpicture}
\legende{Parabole de $f(x)=-x^2+3x+4$ ($a<0$, branches vers le bas). Zone verte : $f(x)\geq 0$ entre les racines ; zones roses : $f(x)<0$ à l'extérieur.}
\end{center}
\end{popfigure}

\courssub{Résolution des inéquations du second degré}

Résoudre $ax^2+bx+c>0$ (ou $\geq 0$, $<0$, $\leq 0$) consiste à trouver toutes les valeurs de $x$ pour lesquelles le trinôme a le signe demandé. On s'appuie sur le théorème du signe et sur un tableau de signes.

\begin{methode}[Résoudre une inéquation du second degré]
\begin{enumerate}
\item \textbf{Réduire} l'inéquation à la forme $ax^2+bx+c$ comparé à $0$ (tout passer dans le même membre si nécessaire).
\item \textbf{Calculer} le discriminant $\Delta=b^2-4ac$.
\item \textbf{Déterminer les racines} éventuelles $x_1$ et $x_2$ (avec $x_1<x_2$).
\item \textbf{Dresser le tableau de signes} du trinôme en appliquant le théorème (signe de $a$ à l'extérieur, signe de $-a$ entre les racines).
\item \textbf{Conclure en intervalles} : lire les zones du signe demandé. Crochets fermés $[\,]$ si l'inégalité est large ($\geq$ ou $\leq$), crochets ouverts $]\,[$ si elle est stricte ($>$ ou $<$).
\end{enumerate}
\end{methode}

\begin{exoresolu}[Résoudre $-x^2+3x+4\geq 0$]
\textbf{Énoncé.} Résoudre dans $\mathbb{R}$ l'inéquation $-x^2+3x+4\geq 0$ et donner l'ensemble des solutions sous forme d'intervalle.
\tcblower
\textbf{Solution détaillée.}
\begin{enumerate}
\item Le trinôme est $f(x)=-x^2+3x+4$ avec $a=-1$, $b=3$, $c=4$.
\item Discriminant : $\Delta=3^2-4\times(-1)\times 4=9+16=25>0$.
\item Racines : $x_1=\dfrac{-3-5}{-2}=4$ et $x_2=\dfrac{-3+5}{-2}=-1$. On ordonne : $x_1=-1$, $x_2=4$.
\item Comme $a=-1<0$, le trinôme est négatif (signe de $a$) à l'extérieur des racines et positif (signe de $-a$) entre les racines. Tableau de signes :
\begin{center}
\begin{tabular}{|c|ccccccc|}\hline
$x$ & $-\infty$ & & $-1$ & & $4$ & & $+\infty$ \\ \hline
$f(x)$ & $-$ & $-$ & $0$ & $+$ & $0$ & $-$ & $-$ \\ \hline
\end{tabular}
\end{center}
\item On veut $f(x)\geq 0$ : on garde la zone « $+$ » et les zéros (inégalité large, crochets fermés).
\end{enumerate}
\textbf{Conclusion :} $S=[-1\,;4]$. C'est exactement la zone verte de la figure ci-dessus.
\end{exoresolu}

\begin{attention}[Erreur fréquente : « entre les racines » sans regarder $a$]
Beaucoup d'élèves mémorisent « la solution est entre les racines ». C'est \textbf{faux en général} ! Cela n'est vrai que si $a<0$ pour une inéquation $\geq 0$ (ou $a>0$ pour $\leq 0$). La règle correcte est : \emph{le trinôme est du signe de $a$ à l'extérieur des racines}. Exemple : pour $x^2-3x+2>0$ (ici $a=1>0$, racines $1$ et $2$), la solution est $S=]-\infty\,;1[\,\cup\,]2\,;+\infty[$, c'est-à-dire \textbf{l'extérieur} des racines. Toujours vérifier le signe de $a$ avant de conclure.
\end{attention}

\textbf{Cas particuliers.} Si $\Delta=0$ : par exemple $x^2-4x+4>0$ donne $(x-2)^2>0$, vrai pour tout $x\neq 2$, donc $S=\mathbb{R}\setminus\{2\}$. Si $\Delta<0$ : le trinôme garde un signe constant ; par exemple $x^2+x+1>0$ a pour solution $S=\mathbb{R}$ (car $\Delta=-3<0$ et $a>0$), tandis que $x^2+x+1<0$ n'a aucune solution ($S=\varnothing$).

\courssub{Inéquations se ramenant au second degré}

\textbf{Inéquations produit.} Une inéquation comme $(x-1)(x^2-3x+2)\leq 0$ se résout par un tableau de signes unique comportant une ligne par facteur : on étudie le signe de chaque facteur (le trinôme avec le théorème), puis on applique la règle des signes sur la ligne produit.

\textbf{Inéquations avec dénominateur.} Pour résoudre $\dfrac{x^2-3x+2}{x-4}\geq 0$, on commence par exclure les \cle{valeurs interdites} : ici $x\neq 4$ (le dénominateur ne peut pas être nul). On dresse ensuite un tableau de signes avec une ligne pour le numérateur, une pour le dénominateur, et une pour le quotient. \textbf{Attention} : la valeur interdite $x=4$ est toujours exclue de la solution (double barre dans le tableau, crochet ouvert), même si l'inégalité est large.

\begin{exoresolu}[Inéquation quotient]
\textbf{Énoncé.} Résoudre $\dfrac{x^2-3x+2}{x-4}\geq 0$.
\tcblower
\textbf{Solution.} Valeur interdite : $x=4$. Numérateur : $x^2-3x+2=(x-1)(x-2)$, racines $1$ et $2$, $a=1>0$ donc positif à l'extérieur, négatif entre. Tableau de signes :
\begin{center}
\begin{tabular}{|c|ccccccccc|}\hline
$x$ & $-\infty$ & & $1$ & & $2$ & & $4$ & & $+\infty$ \\ \hline
$x^2-3x+2$ & $+$ & $+$ & $0$ & $-$ & $0$ & $+$ & $+$ & $+$ & $+$ \\ \hline
$x-4$ & $-$ & $-$ & $-$ & $-$ & $-$ & $-$ & $\|$ & $+$ & $+$ \\ \hline
Quotient & $-$ & $-$ & $0$ & $+$ & $0$ & $-$ & $\|$ & $+$ & $+$ \\ \hline
\end{tabular}
\end{center}
On garde les zones « $+$ » et les zéros du numérateur, en excluant $x=4$ :
\[ S=[1\,;2]\,\cup\,]4\,;+\infty[. \]
\end{exoresolu}

\courssub{Application : rentabilité d'un atelier}

\begin{exoresolu}[Bénéfice positif ou nul]
\textbf{Énoncé.} Un atelier de menuiserie à Douala fabrique et vend $x$ tables par semaine. Son bénéfice hebdomadaire, en milliers de francs CFA, est modélisé par $B(x)=-x^2+30x-125$ pour $x\in[0\,;30]$. Déterminer pour quelles quantités produites le bénéfice est positif ou nul (l'atelier est rentable).
\tcblower
\textbf{Solution.} On résout $B(x)\geq 0$, c'est-à-dire $-x^2+30x-125\geq 0$.
\begin{enumerate}
\item $\Delta=30^2-4\times(-1)\times(-125)=900-500=400>0$, donc $\sqrt{\Delta}=20$.
\item Racines : $x_1=\dfrac{-30-20}{-2}=25$ et $x_2=\dfrac{-30+20}{-2}=5$. On ordonne : $x_1=5$, $x_2=25$.
\item Comme $a=-1<0$, le trinôme est positif \textbf{entre} les racines.
\end{enumerate}
\textbf{Conclusion :} $S=[5\,;25]$. L'atelier est rentable lorsqu'il produit entre 5 et 25 tables par semaine ; en dessous de 5 (frais fixes non couverts) ou au-dessus de 25 (surcharge, invendus), il perd de l'argent.
\end{exoresolu}

\begin{exercice}[À toi de jouer]
Une entreprise de soudure a un coût de production modélisé, en milliers de FCFA, par $C(x)=x^2-14x+40$ pour $x$ pièces fabriquées. Pour quelles valeurs de $x$ le coût est-il inférieur ou égal à \num{8} milliers de FCFA ? (Indication : résoudre $C(x)\leq 8$.)
\end{exercice}

\begin{corrige}[Exercice 1]
On résout $x^2-14x+40\leq 8$, soit $x^2-14x+32\leq 0$.
Discriminant : $\Delta=(-14)^2-4\times 1\times 32=196-128=68$.
Racines : $x_1=\dfrac{14-\sqrt{68}}{2}=7-\sqrt{17}$ et $x_2=\dfrac{14+\sqrt{68}}{2}=7+\sqrt{17}$.
Comme $a=1>0$, le trinôme est négatif ou nul entre les racines.
\textbf{Conclusion :} $S=[7-\sqrt{17}\,;7+\sqrt{17}]$, soit approximativement $[\num{2,88}\,;\num{11,12}]$. Le coût reste sous les \num{8} milliers de FCFA pour une production comprise entre environ 3 et 11 pièces.
\end{corrige}

\begin{savaistu}[Poutres et rentabilité]
Les inéquations du second degré sont des outils quotidiens de l'ingénieur et du gestionnaire. En génie civil, la flèche (déformation) d'une poutre sous charge s'exprime par des polynômes, et vérifier qu'elle reste sous une limite de sécurité revient à résoudre une inéquation : c'est ainsi qu'on dimensionne les poutres des ponts et des bâtiments. En gestion, le bénéfice d'une entreprise est souvent modélisé par un trinôme : déterminer la « zone de rentabilité » (comme dans notre atelier de menuiserie) est exactement la résolution d'une inéquation du second degré. Les économistes appellent les racines les \emph{seuils de rentabilité}.
\end{savaistu}

\coursec{Systèmes de deux équations linéaires à deux inconnues}

\courssub{Introduction et rappel : des équations aux systèmes}

Après avoir étudié la factorisation des trinômes, le signe des expressions du second degré et la résolution des inéquations associées, nous abordons maintenant un outil indispensable pour modéliser des situations où \textbf{deux conditions} doivent être satisfaites \textbf{simultanément}. En atelier, sur un chantier ou en gestion, on rencontre souvent des problèmes à deux inconnues : « Combien de pièces de type A et de type B pour utiliser tout le stock de matière première et de main-d'œuvre ? », « Quel mélange de deux alliages pour obtenir un pourcentage précis de cuivre ? ». Ces situations conduisent naturellement à un \textbf{système de deux équations linéaires à deux inconnues}.

\courssub{Définition et interprétation graphique}

\begin{definition}[Système linéaire 2×2 et couple solution]
Un \textbf{système de deux équations linéaires à deux inconnues} (notées $x$ et $y$) s'écrit sous la forme canonique :
\[
\begin{cases}
a x + b y = c \\
a' x + b' y = c'
\end{cases}
\]
où $a, b, c, a', b', c' \in \mathbb{R}$ et les couples $(a,b)$ et $(a',b')$ sont non nuls.
Un \textbf{couple solution} du système est un couple $(x_0 ; y_0) \in \mathbb{R}^2$ qui vérifie \textbf{les deux équations} simultanément.
L'ensemble des solutions du système est noté $S$.
\end{definition}

Chaque équation $ax+by=c$ représente une \textbf{droite} dans un repère orthonormé $(O;\vec{\imath},\vec{\jmath})$ (droite d'équation cartésienne $y = -\frac{a}{b}x + \frac{c}{b}$ si $b \neq 0$, ou droite verticale $x = \frac{c}{a}$ si $b=0$). Résoudre le système, c'est trouver les points d'intersection de ces deux droites. Il y a trois cas de figure, illustrés ci-dessous.

\begin{popfigure}
\begin{tikzpicture}[scale=0.9, >=Stealth]
% Styles
\tikzset{
    axis/.style={->, thick, popdark},
    grid/.style={very thin, popmuted},
    droita/.style={thick, popblue},
    droiteb/.style={thick, poporange},
    point/.style={circle, fill=poppink, inner sep=2pt},
    labelstyle/.style={font=\small, popdark}
}
% Repère commun pour les 3 sous-figures
\def\drawaxes#1#2#3{
    \begin{scope}[shift={(#1)}]
        \draw[grid] (-2.5,-2.5) grid (2.5,2.5);
        \draw[axis] (-2.5,0) -- (2.5,0) node[right] {$x$};
        \draw[axis] (0,-2.5) -- (0,2.5) node[above] {$y$};
        \node[labelstyle] at (-2.2,2.2) {#2};
        \node[labelstyle] at (2.2,-2.2) {#3};
    \end{scope}
}
% Cas 1 : Sécantes (Solution unique)
\drawaxes{0,0}{$d_1$}{$d_2$}
\draw[droita] (-2.5, 1.5) -- (2.5, -1.5) node[right, pos=1.1] {$d_1 : x+y=0$};
\draw[droiteb] (-2.5, -2) -- (2.5, 2) node[right, pos=1.1] {$d_2 : -x+y=0$};
\fill[point] (0,0) circle (2pt) node[above right] {$I(x_0,y_0)$};

% Cas 2 : Parallèles distinctes (Pas de solution)
\drawaxes{6,0}{$d_1$}{$d_2$}
\draw[droita] (-2.5+6, 1.5) -- (2.5+6, -1.5) node[right, pos=1.1] {$d_1 : x+y=0$};
\draw[droiteb] (-2.5+6, -0.5) -- (2.5+6, 2.5) node[right, pos=1.1] {$d_2 : x+y=2$};
% Pas de point d'intersection

% Cas 3 : Confondues (Infinité de solutions)
\drawaxes{12,0}{$d_1$}{$d_2$}
\draw[droita] (-2.5+12, 1.5) -- (2.5+12, -1.5) node[right, pos=1.1] {$d_1 : x+y=0$};
\draw[droiteb, dashed, opacity=0.7] (-2.5+12, 1.5) -- (2.5+12, -1.5) node[right, pos=1.1] {$d_2 : 2x+2y=0$};
\fill[point] (12+1, -1) circle (2pt) node[above right] {$M \in d_1=d_2$};
\fill[point] (12-1, 1) circle (2pt) node[below left] {$N \in d_1=d_2$};
\end{tikzpicture}
\legende{Les trois positions relatives de deux droites dans le plan : sécantes (solution unique), parallèles distinctes (aucune solution), confondues (infinité de solutions).}
\end{popfigure}

\begin{propriete}[Nombre de solutions d'un système 2×2]
Un système de deux équations linéaires à deux inconnues admet :
\begin{itemize}
    \item \textbf{Un unique couple solution} si les droites sont \textbf{sécantes} (coefficients directeurs différents).
    \item \textbf{Aucune solution} ($S = \varnothing$) si les droites sont \textbf{parallèles distinctes} (même coefficient directeur, ordonnées à l'origine différentes).
    \item \textbf{Une infinité de solutions} si les droites sont \textbf{confondues} (équations proportionnelles).
\end{itemize}
Cette propriété est admise au Probatoire ; elle se justifie immédiatement par l'interprétation graphique.
\end{propriete}

Pour discuter le nombre de solutions sans tracer les droites, on utilise le \textbf{déterminant} du système :
\[
\Delta = \begin{vmatrix} a & b \\ a' & b' \end{vmatrix} = a b' - a' b.
\]
\begin{itemize}
    \item Si $\Delta \neq 0$ : les droites sont sécantes $\Rightarrow$ \textbf{solution unique}.
    \item Si $\Delta = 0$ : les droites sont parallèles. Il faut comparer les seconds membres.
        \begin{itemize}
            \item Si les équations sont proportionnelles (c'est-à-dire $\frac{a}{a'} = \frac{b}{b'} = \frac{c}{c'}$ lorsque les dénominateurs sont non nuls, ou equivalentement $ab'=a'b$, $ac'=a'c$ et $bc'=b'c$) $\Rightarrow$ \textbf{infinité de solutions}.
            \item Sinon $\Rightarrow$ \textbf{aucune solution} (système impossible).
        \end{itemize}
\end{itemize}

\courssub{Méthodes de résolution algébrique (cas $\Delta \neq 0$)}

Lorsque le système admet une solution unique, trois méthodes classiques permettent de la trouver. Le choix de la méthode dépend de la forme des équations : on vise la \textbf{simplicité des calculs} et la \textbf{limitation des fractions}.

\begin{methode}[Résolution par substitution]
\emph{Idéale quand une inconnue est déjà isolée ou s'isole facilement (coefficient $\pm 1$).}
\begin{enumerate}
    \item Exprimer une inconnue en fonction de l'autre dans l'une des équations (ex: $x = \dots$ ou $y = \dots$).
    \item Substituer cette expression dans l'autre équation : on obtient une équation à une inconnue.
    \item Résoudre cette équation pour trouver la valeur de la première inconnue.
    \item \textbf{Reporter} cette valeur dans l'expression de l'étape 1 pour trouver la deuxième inconnue.
    \item Vérifier le couple $(x ; y)$ dans les \textbf{deux} équations initiales.
\end{enumerate}
\end{methode}

\begin{methode}[Résolution par combinaison linéaire (addition/soustraction)]
\emph{Idéale quand les coefficients d'une inconnue sont opposés ou facilement rendus opposés par multiplication.}
\begin{enumerate}
    \item Multiplier éventuellement les équations par des nombres non nuls pour obtenir des coefficients \textbf{opposés} pour $x$ ou pour $y$.
    \item Additionner (ou soustraire) membre à membre les deux équations pour éliminer une inconnue.
    \item Résoudre l'équation obtenue pour trouver la première inconnue.
    \item \textbf{Reporter} cette valeur dans l'une des équations initiales pour trouver la deuxième inconnue.
    \item Vérifier le couple $(x ; y)$ dans les \textbf{deux} équations initiales.
\end{enumerate}
\end{methode}

\begin{methode}[Résolution par comparaison]
\emph{Utile quand les deux équations s'écrivent facilement sous la forme $y = \dots$ (ou $x = \dots$).}
\begin{enumerate}
    \item Exprimer $y$ en fonction de $x$ dans les \textbf{deux} équations : $y = f_1(x)$ et $y = f_2(x)$.
    \item Égaler les deux expressions : $f_1(x) = f_2(x)$.
    \item Résoudre pour trouver $x$.
    \item Reporter $x$ dans l'une des expressions $y = f_i(x)$ pour trouver $y$.
    \item Vérifier.
\end{enumerate}
\end{methode}

\begin{attention}[Pièges classiques à éviter]
\begin{itemize}
    \item \textbf{Oublier le report :} Trouver $x$ mais ne pas calculer $y$ (ou inversement). Le couple solution est incomplet.
    \item \textbf{Ne pas vérifier :} Une erreur de signe lors du report donne un couple qui ne vérifie qu'une équation. La vérification dans les \textbf{deux} équations est obligatoire.
    \item \textbf{Mauvaise multiplication :} En combinaison, multiplier seulement le terme en $x$ et pas le second membre $c$.
    \item \textbf{Confondre « aucune solution » et « solution $(0;0)$ :} $(0;0)$ est un couple solution valide si $c=c'=0$. « Aucune solution » signifie $S=\varnothing$.
\end{itemize}
\end{attention}

\courssub{Exemple chiffré détaillé : les trois méthodes}

\begin{exoresolu}[Résolution par substitution, combinaison et comparaison]
Résoudre le système suivant et vérifier la solution :
\[
(S) : \begin{cases}
2x + 3y = 12 & (E_1) \\
x - y = 1 & (E_2)
\end{cases}
\]
\tcblower
\textbf{1. Discussion préalable (Déterminant)}
$\Delta = \begin{vmatrix} 2 & 3 \\ 1 & -1 \end{vmatrix} = 2 \times (-1) - 1 \times 3 = -2 - 3 = -5 \neq 0$.
Le système admet \textbf{une unique solution}.

\textbf{2. Méthode par substitution (choix naturel : $x$ s'isole dans $E_2$)}
\begin{itemize}
    \item Étape 1 : Dans $(E_2)$, $x = y + 1$.
    \item Étape 2 : Substitution dans $(E_1)$ : $2(y+1) + 3y = 12$.
    \item Étape 3 : $2y + 2 + 3y = 12 \iff 5y + 2 = 12 \iff 5y = 10 \iff y = 2$.
    \item Étape 4 : Report dans $x = y + 1 \Rightarrow x = 2 + 1 = 3$.
    \item Couple candidat : $(3 ; 2)$.
\end{itemize}

\textbf{3. Méthode par combinaison linéaire (élimination de $x$)}
\begin{itemize}
    \item Étape 1 : Multiplier $(E_2)$ par $-2$ pour opposer les coefficients de $x$ :
    $(E_2') : -2x + 2y = -2$.
    \item Étape 2 : Additionner $(E_1)$ et $(E_2')$ membre à membre :
    $(2x - 2x) + (3y + 2y) = 12 + (-2) \iff 5y = 10 \iff y = 2$.
    \item Étape 3 : Report de $y=2$ dans $(E_2)$ (plus simple) : $x - 2 = 1 \iff x = 3$.
    \item Couple candidat : $(3 ; 2)$.
\end{itemize}

\textbf{4. Méthode par comparaison (expression de $y$)}
\begin{itemize}
    \item $(E_1) \iff 3y = 12 - 2x \iff y = 4 - \frac{2}{3}x$.
    \item $(E_2) \iff -y = 1 - x \iff y = x - 1$.
    \item Égaler : $4 - \frac{2}{3}x = x - 1 \iff 5 = \frac{5}{3}x \iff x = 3$.
    \item Report : $y = 3 - 1 = 2$.
    \item Couple candidat : $(3 ; 2)$.
\end{itemize}

\textbf{5. Vérification obligatoire dans les DEUX équations}
\begin{itemize}
    \item Dans $(E_1)$ : $2 \times 3 + 3 \times 2 = 6 + 6 = 12$ \checkmark
    \item Dans $(E_2)$ : $3 - 2 = 1$ \checkmark
\end{itemize}
\textbf{Conclusion :} L'ensemble solution est $S = \{(3 ; 2)\}$.
\end{exoresolu}

\courssub{Cas particuliers : $\Delta = 0$ (Pas de solution ou infinité)}

\begin{exoresolu}[Discussion et résolution selon $\Delta$]
Étudier et résoudre les systèmes suivants :
\[
(S_1) : \begin{cases} 2x + 4y = 6 \\ x + 2y = 5 \end{cases}
\qquad
(S_2) : \begin{cases} 3x - y = 2 \\ 6x - 2y = 4 \end{cases}
\]
\tcblower
\textbf{Système $(S_1)$ :}
$\Delta = \begin{vmatrix} 2 & 4 \\ 1 & 2 \end{vmatrix} = 2\times2 - 1\times4 = 0$.
Les droites sont parallèles. Comparons les seconds membres :
$\frac{2}{1} = 2$, $\frac{4}{2} = 2$, mais $\frac{6}{5} = 1,2 \neq 2$.
Les équations ne sont pas proportionnelles. Le système est \textbf{impossible}.
$S_1 = \varnothing$.
\textit{Vérification rapide :} $(E_1) \iff x+2y=3$ alors que $(E_2) \iff x+2y=5$. Impossible d'avoir $3=5$.

\textbf{Système $(S_2)$ :}
$\Delta = \begin{vmatrix} 3 & -1 \\ 6 & -2 \end{vmatrix} = 3\times(-2) - 6\times(-1) = -6 + 6 = 0$.
Comparons : $\frac{3}{6} = \frac{1}{2}$, $\frac{-1}{-2} = \frac{1}{2}$, $\frac{2}{4} = \frac{1}{2}$.
Tous les rapports sont égaux : la deuxième équation est le double de la première.
Les droites sont \textbf{confondues}. Il y a une \textbf{infinité de solutions}.
On exprime $y$ en fonction de $x$ (paramètre) à partir de $(E_1)$ : $y = 3x - 2$.
$S_2 = \{(t ; 3t-2) \mid t \in \mathbb{R}\}$.
\end{exoresolu}

\courssub{Modélisation et problèmes concrets (Mise en équation)}

La force des systèmes linéaires réside dans la modélisation de situations réelles à deux contraintes. La démarche est rigoureuse :

\begin{methode}[Démarche de modélisation par un système 2×2]
\begin{enumerate}
    \item \textbf{Lire et comprendre} le problème. Identifier les deux inconnues recherchées.
    \item \textbf{Choisir les variables} : « Soit $x$ ... et $y$ ... » (préciser l'unité).
    \item \textbf{Traduire} chaque contrainte/condition en une équation linéaire.
    \item \textbf{Résoudre} le système obtenu (choisir la méthode la plus adaptée).
    \item \textbf{Interpréter} la solution dans le contexte du problème (vérifier la cohérence : nombres positifs, entiers si comptage...).
    \item \textbf{Conclure} par une phrase réponse.
\end{enumerate}
\end{methode}

\begin{exoresolu}[Problème d'achat au marché (Gestion/Comptabilité)]
Un chef de chantier achète 5 sacs de ciment et 3 fers à béton pour un total de 37 000 FCFA.
La semaine suivante, il achète 2 sacs de ciment et 5 fers à béton pour 30 000 FCFA.
Déterminer le prix unitaire d'un sac de ciment et d'un fer à béton.
\tcblower
\textbf{Modélisation}
\begin{itemize}
    \item Soit $x$ le prix d'un sac de ciment (en FCFA).
    \item Soit $y$ le prix d'un fer à béton (en FCFA).
    \item Achat 1 : $5x + 3y = 37\,000$.
    \item Achat 2 : $2x + 5y = 30\,000$.
\end{itemize}
Système : $\begin{cases} 5x + 3y = 37\,000 \\ 2x + 5y = 30\,000 \end{cases}$

\textbf{Résolution (Combinaison : élimination de $y$)}
Multiplier la 1ère équation par 5 et la 2ème par 3 :
$\begin{cases} 25x + 15y = 185\,000 \\ 6x + 15y = 90\,000 \end{cases}$
Soustraction membre à membre : $(25x - 6x) + (15y - 15y) = 185\,000 - 90\,000$
$19x = 95\,000 \iff x = 5\,000$.
Report dans $2x+5y=30\,000$ : $10\,000 + 5y = 30\,000 \iff 5y = 20\,000 \iff y = 4\,000$.

\textbf{Vérification :} $5\times5000+3\times4000=25000+12000=37000$ \checkmark. $2\times5000+5\times4000=10000+20000=30000$ \checkmark.

\textbf{Conclusion :} Le sac de ciment coûte \textbf{5 000 FCFA} et le fer à béton coûte \textbf{4 000 FCFA}.
\end{exoresolu}

\begin{exoresolu}[Problème de mélange (Chimie / Métallurgie)]
Un soudeur doit préparer 100 kg d'un alliage contenant 45\% de cuivre.
Il dispose de deux alliages de base :
\begin{itemize}
    \item Alliage A : 30\% de cuivre.
    \item Alliage B : 60\% de cuivre.
\end{itemize}
Quelle masse de chaque alliage doit-il mélanger ?
\tcblower
\textbf{Modélisation}
\begin{itemize}
    \item Soit $x$ la masse de l'alliage A (en kg).
    \item Soit $y$ la masse de l'alliage B (en kg).
    \item Contrainte 1 (Masse totale) : $x + y = 100$.
    \item Contrainte 2 (Masse de cuivre) : $0,30x + 0,60y = 0,45 \times 100 = 45$.
\end{itemize}
Système : $\begin{cases} x + y = 100 \\ 0,3x + 0,6y = 45 \end{cases}$

\textbf{Résolution (Substitution : $x$ s'isole dans la 1ère)}
$x = 100 - y$.
$0,3(100 - y) + 0,6y = 45 \iff 30 - 0,3y + 0,6y = 45 \iff 30 + 0,3y = 45$.
$0,3y = 15 \iff y = \frac{15}{0,3} = 50$.
$x = 100 - 50 = 50$.

\textbf{Vérification :} Cuivre : $0,3\times50 + 0,6\times50 = 15 + 30 = 45$ kg (soit 45\% de 100 kg) \checkmark.

\textbf{Conclusion :} Il faut mélanger \textbf{50 kg de l'alliage A} et \textbf{50 kg de l'alliage B}.
\end{exoresolu}

\begin{savaistu}[Le déterminant et l'algèbre linéaire]
Le déterminant $\Delta = ab' - a'b$ que nous utilisons pour discuter le nombre de solutions est la notion fondatrice de l'\textbf{algèbre linéaire}. En Terminale et dans les études supérieures (BTS, DUT, Licence), on généralise cette idée aux matrices $n \times n$. Un système $AX=B$ a une solution unique si et seulement si $\det(A) \neq 0$ (matrice inversible). Le calcul de $\Delta$ pour un système 2×2 est donc votre première rencontre avec le concept d'\textbf{inversibilité d'une matrice}. Au Cameroun, cette notion est cruciale pour l'analyse de structures (stabilité des treillis en génie civil) et le traitement du signal (électronique).
\end{savaistu}

\begin{exercice}[Entraînement : Résolution et discussion]
Résoudre les systèmes suivants (préciser la méthode choisie et vérifier) :
\begin{enumerate}
    \item $\begin{cases} 3x - 2y = 7 \\ x + 4y = 10 \end{cases}$
    \item $\begin{cases} 2x + 5y = 1 \\ 4x + 10y = 3 \end{cases}$
    \item $\begin{cases} 0,5x - 0,2y = 1,1 \\ x + y = 4 \end{cases}$
\end{enumerate}
\end{exercice}

\begin{exercice}[Modélisation : Problème de partage]
Un héritage de 2 500 000 FCFA doit être partagé entre deux frères. L'aîné reçoit 300 000 FCFA de plus que le double de la part du cadet. Déterminer la part de chacun par mise en système.
\end{exercice}

\begin{exercice}[Modélisation : Électricité (Loi des mailles)]
Dans un circuit électrique à deux mailles, les intensités $I_1$ et $I_2$ (en Ampères) vérifient :
$\begin{cases} 10 I_1 + 5 I_2 = 12 \\ 5 I_1 + 15 I_2 = 9 \end{cases}$
Déterminer $I_1$ et $I_2$.
\end{exercice}

\begin{corrige}[Exercice 1]
\textbf{1.} $\Delta = 3\times4 - 1\times(-2) = 12+2=14 \neq 0$. Substitution : $x=10-4y$. $3(10-4y)-2y=7 \Rightarrow 30-14y=7 \Rightarrow y=23/14$. $x=10-92/14=48/14=24/7$. $S=\{(24/7 ; 23/14)\}$.
\textbf{2.} $\Delta = 2\times10 - 4\times5 = 0$. $\frac{2}{4}=\frac{1}{2}$, $\frac{5}{10}=\frac{1}{2}$, $\frac{1}{3}\neq\frac{1}{2}$. Système impossible. $S=\varnothing$.
\textbf{3.} $\Delta = 0,5\times1 - 1\times(-0,2) = 0,7 \neq 0$. Substitution : $x=4-y$. $0,5(4-y)-0,2y=1,1 \Rightarrow 2-0,7y=1,1 \Rightarrow 0,7y=0,9 \Rightarrow y=9/7$. $x=19/7$. $S=\{(19/7 ; 9/7)\}$.
\end{corrige}

\begin{corrige}[Exercice 2]
Soit $x$ la part du cadet, $y$ la part de l'aîné.
$\begin{cases} x + y = 2\,500\,000 \\ y = 2x + 300\,000 \end{cases}$
Substitution : $x + 2x + 300\,000 = 2\,500\,000 \Rightarrow 3x = 2\,200\,000 \Rightarrow x = 733\,333,\overline{3}$ FCFA.
$y = 2\,500\,000 - 733\,333,\overline{3} = 1\,766\,666,\overline{6}$ FCFA.
Cadet : 733 333 FCFA ; Aîné : 1 766 667 FCFA (arrondis).
\end{corrige}

\begin{corrige}[Exercice 3]
$\Delta = 10\times15 - 5\times5 = 150 - 25 = 125 \neq 0$.
Combinaison : $\times 3$ sur 1ère : $30I_1 + 15I_2 = 36$.
Soustraction 2ème : $(30-5)I_1 + (15-15)I_2 = 36-9 \Rightarrow 25I_1 = 27 \Rightarrow I_1 = 1,08$ A.
Report : $10(1,08) + 5I_2 = 12 \Rightarrow 10,8 + 5I_2 = 12 \Rightarrow 5I_2 = 1,2 \Rightarrow I_2 = 0,24$ A.
$I_1 = 1,08$ A ; $I_2 = 0,24$ A.
\end{corrige}

\coursec{Équations et inéquations associées aux fonctions}

Dans les sections précédentes, nous avons résolu des \cle{systèmes} de deux équations linéaires à deux inconnues : géométriquement, il s'agissait de trouver les points d'intersection de deux \emph{droites}. Nous allons maintenant généraliser cette idée : remplacer les droites par des \cle{courbes} quelconques (paraboles, branches d'hyperboles, courbes de fonctions quelconques). Chercher l'intersection de deux courbes, c'est résoudre une équation du type $f(x)=g(x)$ ; comparer les positions relatives de deux courbes, c'est résoudre une inéquation du type $f(x)<g(x)$. Cette section fait le pont permanent entre le \emph{graphique} (ce que l'on voit ou que l'on trace) et l'\emph{algébrique} (ce que l'on calcule), une compétence centrale pour le Probatoire.

\courssub{Résolution graphique de l'équation $f(x)=k$}

\textbf{Intuition.} Sur la courbe représentative $\mathcal{C}_f$ d'une fonction $f$, l'ordonnée d'un point d'abscisse $x$ est exactement $f(x)$. Dire que $f(x)=k$ revient donc à chercher les points de la courbe situés « à la hauteur » $k$, c'est-à-dire les points d'intersection de $\mathcal{C}_f$ avec la droite horizontale d'équation $y=k$.

\begin{propriete}[Interprétation graphique de $f(x)=k$]
Les solutions de l'équation $f(x)=k$ sont les \cle{abscisses} des points d'intersection de la courbe représentative $\mathcal{C}_f$ avec la droite d'équation $y=k$.
\begin{itemize}
\item Si la droite coupe la courbe en plusieurs points, l'équation admet plusieurs solutions.
\item Si la droite est tangente à la courbe, l'équation admet une solution double (racine double).
\item Si la droite ne rencontre pas la courbe, l'équation n'a \textbf{aucune} solution dans $\mathbb{R}$.
\end{itemize}
\end{propriete}

\begin{exemplebox}[Lecture graphique : $f(x)=x^2-2x-1$]
Soit $f$ la fonction définie sur $\mathbb{R}$ par $f(x)=x^2-2x-1$, de courbe $\mathcal{C}_f$ (parabole). Pour résoudre graphiquement $f(x)=2$, on trace la droite horizontale $y=2$. Elle coupe la parabole en deux points $A$ et $B$ d'abscisses $-1$ et $3$. Les solutions sont donc $x=-1$ et $x=3$.
\textbf{Vérification algébrique :} $f(-1)=(-1)^2-2(-1)-1=1+2-1=2$ et $f(3)=3^2-2\times3-1=9-6-1=2$.
\end{exemplebox}

\begin{popfigure}
\begin{center}
\begin{tikzpicture}
\begin{axis}[width=0.85\linewidth,height=7cm,axis lines=middle,xlabel=$x$,ylabel=$y$,xmin=-2.5,xmax=4.5,ymin=-3,ymax=6,xtick={-1,1,3},ytick={2},grid=both,grid style={popline!40},legend style={at={(0.02,0.98)},anchor=north west,font=\small}]
\addplot[popcurve,domain=-2:4,samples=200]{x^2-2*x-1};
\addlegendentry{$\mathcal{C}_f : y=x^2-2x-1$}
\addplot[poporange,thick,domain=-2.5:4.5,samples=2]{2};
\addlegendentry{$y=2$}
\addplot[only marks,mark=*,popdark,mark size=2.5pt] coordinates {(-1,2) (3,2)};
\node[popdark,above left] at (axis cs:-1,2) {$A(-1\,;\,2)$};
\node[popdark,above right] at (axis cs:3,2) {$B(3\,;\,2)$};
\draw[popmuted,dashed] (axis cs:-1,2) -- (axis cs:-1,0) node[below] {$-1$};
\draw[popmuted,dashed] (axis cs:3,2) -- (axis cs:3,0) node[below] {$3$};
\end{axis}
\end{tikzpicture}
\end{center}
\legende{Résolution graphique de $f(x)=2$ : les solutions sont les abscisses $-1$ et $3$ des points $A$ et $B$.}
\end{popfigure}

\begin{attention}[Image ou antécédent ? Piège classique au Probatoire]
Résoudre $f(x)=k$, c'est chercher les \cle{antécédents} de $k$ par $f$ : on lit les solutions sur l'\textbf{axe des abscisses} (axe $x$). Ne pas confondre avec le calcul de l'image $f(a)$, qui se lit sur l'axe des ordonnées (axe $y$). \textbf{Erreur fréquente :} donner les ordonnées des points d'intersection (ici $2$) au lieu de leurs abscisses ($-1$ et $3$). Sur la copie, écrire « Les solutions sont les abscisses... » est attendu.
\end{attention}

\courssub{Résolution graphique de l'équation $f(x)=g(x)$}

\textbf{Intuition.} Deux courbes $\mathcal{C}_f$ et $\mathcal{C}_g$ se croisent lorsque, pour une même abscisse $x$, elles ont la même ordonnée, c'est-à-dire lorsque $f(x)=g(x)$.

\begin{propriete}[Interprétation graphique de $f(x)=g(x)$]
Les solutions de l'équation $f(x)=g(x)$ sont les \cle{abscisses} des points d'intersection des courbes représentatives $\mathcal{C}_f$ et $\mathcal{C}_g$.
\end{propriete}

\begin{methode}[Résoudre graphiquement puis confirmer algébriquement $f(x)=g(x)$]
\begin{enumerate}
\item \textbf{Lire} sur le graphique les abscisses des points d'intersection des deux courbes : ce sont les solutions \emph{approchées ou exactes} de l'équation.
\item \textbf{Écrire} l'équation algébrique : $f(x)=g(x)$, puis transposer tous les termes dans le membre de gauche : $f(x)-g(x)=0$.
\item \textbf{Résoudre} l'équation obtenue (souvent un trinôme du second degré : calcul du discriminant $\Delta$).
\item \textbf{Comparer} : les solutions calculées doivent coïncider exactement avec les abscisses lues graphiquement. En cas de désaccord, rechercher l'erreur de calcul ou de lecture graphique.
\end{enumerate}
\end{methode}

\begin{exoresolu}[Intersection d'une parabole et d'une droite : $f(x)=x^2-2x-1$, $g(x)=2$]
Résoudre algébriquement $f(x)=g(x)$ et vérifier la cohérence avec le graphique ci-dessus.
\tcblower
\textbf{Solution.}
On écrit l'équation et on ramène tout à gauche :
\begin{align*}
f(x) &= g(x)\\
x^2-2x-1 &= 2\\
x^2-2x-3 &= 0
\end{align*}
C'est un trinôme du second degré avec $a=1$, $b=-2$, $c=-3$.
\[ \Delta = b^2-4ac = (-2)^2-4\times 1\times(-3) = 4+12 = 16 > 0. \]
Deux solutions réelles distinctes :
\[ x_1=\frac{-b-\sqrt{\Delta}}{2a}=\frac{2-4}{2}=-1 \qquad ; \qquad x_2=\frac{-b+\sqrt{\Delta}}{2a}=\frac{2+4}{2}=3. \]
L'ensemble des solutions est $S=\{-1\,;\,3\}$, ce qui confirme exactement les abscisses des points $A$ et $B$ lues sur le graphique.
\end{exoresolu}

\courssub{Résolution graphique d'une inéquation $f(x)<g(x)$ (ou $\leq$, $>$, $\geq$)}

\textbf{Intuition.} Dire que $f(x)<g(x)$ signifie que, pour l'abscisse $x$, le point de $\mathcal{C}_f$ est \emph{strictement en dessous} du point de $\mathcal{C}_g$. Résoudre l'inéquation, c'est repérer les intervalles d'abscisses où la courbe $\mathcal{C}_f$ se trouve strictement sous la courbe $\mathcal{C}_g$.

\begin{propriete}[Interprétation graphique de $f(x)<g(x)$]
Les solutions de l'inéquation $f(x)<g(x)$ sont les abscisses des points pour lesquels la courbe $\mathcal{C}_f$ est située \cle{strictement en dessous} de la courbe $\mathcal{C}_g$.
Pour $f(x)\leq g(x)$, on inclut les abscisses des points d'intersection (courbes confondues).
Pour $f(x)>g(x)$ (resp. $\geq$), on cherche où $\mathcal{C}_f$ est au-dessus de $\mathcal{C}_g$.
\end{propriete}

\begin{popfigure}
\begin{center}
\begin{tikzpicture}
\begin{axis}[width=0.85\linewidth,height=7cm,axis lines=middle,xlabel=$x$,ylabel=$y$,xmin=-2.5,xmax=4.5,ymin=-3,ymax=6,xtick={-1,1,3},ytick=\empty,legend style={at={(0.02,0.98)},anchor=north west,font=\small}]
\addplot[popcurve,domain=-2:4,samples=200]{x^2-2*x-1};
\addlegendentry{$\mathcal{C}_f$}
\addplot[poporange,thick,domain=-2.5:4.5,samples=2]{2};
\addlegendentry{$\mathcal{C}_g : y=2$}
\addplot[only marks,mark=*,popdark,mark size=2.5pt] coordinates {(-1,2) (3,2)};
% Zone solution sur l'axe des abscisses
\draw[popgreen,ultra thick] (axis cs:-1,0) -- (axis cs:3,0);
\addplot[only marks,mark=*,popgreen,mark size=3pt] coordinates {(-1,0) (3,0)};
\node[popgreen,below] at (axis cs:1,-0.3) {$S=\left]-1\,;\,3\right[$};
\end{axis}
\end{tikzpicture}
\end{center}
\legende{La parabole $\mathcal{C}_f$ est sous la droite $\mathcal{C}_g$ pour $x\in\left]-1\,;\,3\right[$ : l'intervalle solution est tracé en vert épais sur l'axe des abscisses.}
\end{popfigure}

\begin{exemplebox}[Lien fondamental avec le signe du trinôme (Section 3)]
Résoudre $f(x)<g(x)$ avec $f(x)=x^2-2x-1$ et $g(x)=2$ revient algébriquement à résoudre $x^2-2x-3<0$.
Le trinôme $x^2-2x-3$ a pour racines $-1$ et $3$ et son coefficient dominant $a=1$ est positif. D'après le cours sur le signe des trinômes, il est \textbf{négatif entre ses racines}.
Donc $S=]-1\,;\,3[$, exactement l'intervalle lu graphiquement. La résolution algébrique (signe du trinôme) et la résolution graphique (positions relatives des courbes) sont deux visages du même résultat.
\end{exemplebox}

\begin{attention}[Crochets et type d'inégalité : points de vigilance pour l'examen]
\begin{itemize}
\item Pour $f(x)<g(x)$ ou $f(x)>g(x)$ (inégalité \emph{stricte}), les abscisses des points d'intersection sont \textbf{exclues} : on utilise des crochets ouverts $]a\,;\,b[$ ou $]-\infty\,;\,a[$, etc.
\item Pour $f(x)\leq g(x)$ ou $f(x)\geq g(x)$ (inégalité \emph{large}), les abscisses des points d'intersection sont \textbf{incluses} : on utilise des crochets fermés $[a\,;\,b]$ ou $]-\infty\,;\,a]$, etc.
\item Sur la copie du Probatoire, le type de crochet fait partie intégrante de la réponse et est systématiquement vérifié. Ne pas écrire $]-1,3[$ (virgule) mais $]-1\,;\,3[$ (point-virgule).
\end{itemize}
\end{attention}

\courssub{Équations bicarrées}

Certaines équations ne sont pas du second degré en $x$, mais s'y \emph{ramènent} grâce à un changement de variable judicieux. C'est le cas des équations \cle{bicarrées}, qui ne contiennent que des puissances paires de l'inconnue ($x^4$ et $x^2$).

\begin{definition}[Équation bicarrée]
Une équation \cle{bicarrée} est une équation de la forme
\[ ax^4+bx^2+c=0 \qquad (a\neq 0,\ a,b,c\in\mathbb{R}). \]
Elle se résout en posant le changement de variable $X=x^2$ (avec la condition $X\geq 0$), ce qui la transforme en l'équation du second degré $aX^2+bX+c=0$.
\end{definition}

\begin{methode}[Résoudre une équation bicarrée par changement de variable]
\begin{enumerate}
\item \textbf{Poser} $X=x^2$ (donc $X\geq 0$ et $x^4=X^2$).
\item \textbf{Résoudre} l'équation du second degré en $X$ : $aX^2+bX+c=0$ (calcul de $\Delta$, recherche des racines $X_1, X_2$).
\item \textbf{Écarter} toute solution $X_i < 0$, car un carré réel $x^2$ est toujours positif ou nul.
\item \textbf{Revenir} à $x$ : pour chaque solution $X_i \geq 0$ conservée, on a $x=\sqrt{X_i}$ ou $x=-\sqrt{X_i}$ (deux valeurs sauf si $X_i=0$).
\item \textbf{Conclure} en donnant l'ensemble $S$ des solutions réelles.
\end{enumerate}
\end{methode}

\begin{exoresolu}[Résoudre $x^4-5x^2+4=0$]
Résoudre dans $\mathbb{R}$ l'équation $x^4-5x^2+4=0$.
\tcblower
\textbf{Solution.}
Posons $X=x^2$ (donc $X\geq 0$). Alors $x^4=X^2$ et l'équation devient :
\[ X^2-5X+4=0. \]
Discriminant : $\Delta=(-5)^2-4\times 1\times 4=25-16=9>0$.
Deux solutions en $X$ :
\[ X_1=\frac{5-3}{2}=1 \qquad ; \qquad X_2=\frac{5+3}{2}=4. \]
Les deux solutions sont positives (ou nulles), on les garde toutes les deux. Retour à $x$ :
\begin{itemize}
\item $x^2=1 \Rightarrow x=1$ ou $x=-1$ ;
\item $x^2=4 \Rightarrow x=2$ ou $x=-2$.
\end{itemize}
Vérification pour $x=2$ : $2^4-5\times2^2+4=16-20+4=0$.
\textbf{Conclusion :} $S=\{-2\,;\,-1\,;\,1\,;\,2\}$.
\end{exoresolu}

\begin{attention}[Ne pas oublier le retour à $x$ et la condition $X\geq 0$]
Trois erreurs fréquentes au Probatoire :
\begin{enumerate}
\item S'arrêter aux solutions en $X$ (alors que l'inconnue demandée est $x$).
\item Oublier qu'une solution $X>0$ donne \textbf{deux} valeurs de $x$ (une positive et une négative).
\item Conserver une solution $X<0$ : un carré réel n'est jamais négatif, elle ne donne \emph{aucune} solution en $x$.
\end{enumerate}
\end{attention}

\courssub{Équations avec racine carrée se ramenant au second degré}

\textbf{Intuition.} Pour éliminer une racine carrée, l'opération naturelle est d'élever les deux membres au carré. Mais attention : l'élévation au carré n'est \textbf{pas} une transformation à équivalence (elle ne conserve pas l'ensemble des solutions). Elle peut introduire des \cle{solutions parasites}, c'est-à-dire des nombres qui vérifient l'équation élevée au carré mais pas l'équation de départ.

\begin{methode}[Résoudre une équation du type $\sqrt{u(x)}=v(x)$]
\begin{enumerate}
\item \textbf{Isoler} la racine carrée dans un membre (de l'autre côté, une expression $v(x)$).
\item \textbf{Élever au carré} les deux membres : on obtient une équation (souvent du second degré) sans racine : $u(x)=[v(x)]^2$.
\item \textbf{Résoudre} cette équation du second degré.
\item \textbf{Vérifier obligatoirement} chaque solution trouvée en la reportant dans l'équation \emph{initiale} : on ne garde que celles qui la vérifient. (Rappel : $\sqrt{u(x)}\geq 0$, donc nécessairement $v(x)\geq 0$ pour toute solution).
\end{enumerate}
\end{methode}

\begin{exoresolu}[Résoudre $\sqrt{x^2-3}=1$ et $\sqrt{x+3}=x-3$]
\textbf{(a)} Résoudre $\sqrt{x^2-3}=1$. \qquad \textbf{(b)} Résoudre $\sqrt{x+3}=x-3$.
\tcblower
\textbf{(a)} La racine est isolée. Élévation au carré :
\[ x^2-3=1 \iff x^2=4 \iff x=2 \text{ ou } x=-2. \]
\textbf{Vérification dans l'équation initiale :}
\begin{itemize}
\item $x=2$ : $\sqrt{2^2-3}=\sqrt{1}=1$ (vérifié).
\item $x=-2$ : $\sqrt{(-2)^2-3}=\sqrt{1}=1$ (vérifié).
\end{itemize}
$S=\{-2\,;\,2\}$.

\textbf{(b)} La racine est isolée. Élévation au carré :
\[ x+3=(x-3)^2 \iff x+3=x^2-6x+9 \iff x^2-7x+6=0. \]
Discriminant : $\Delta=49-24=25$. Racines : $x_1=\frac{7-5}{2}=1$, $x_2=\frac{7+5}{2}=6$.

\textbf{Vérification obligatoire dans l'équation initiale $\sqrt{x+3}=x-3$ :}
\begin{itemize}
\item Pour $x=1$ : $\sqrt{1+3}=\sqrt{4}=2$ mais $1-3=-2$. Or $2\neq -2$. De plus $x-3=-2<0$ alors que la racine est positive. $x=1$ est une \cle{solution parasite}, on la \textbf{rejette}.
\item Pour $x=6$ : $\sqrt{6+3}=\sqrt{9}=3$ et $6-3=3$. Égalité vérifiée.
\end{itemize}
\textbf{Conclusion :} $S=\{6\}$.
\end{exoresolu}

\begin{attention}[Erreur fatale : l'oubli de la vérification]
L'erreur la plus classique (et la plus sanctionnée) est d'oublier de \textbf{vérifier} les solutions d'une équation avec racine carrée. En élevant au carré, on résout en réalité $\sqrt{u}=v$ \emph{ou} $\sqrt{u}=-v$ : des valeurs étrangères (parasites) apparaissent systématiquement quand $v(x)$ devient négatif. Toute solution non vérifiée dans l'équation de départ fait perdre des points au Probatoire. La vérification n'est pas optionnelle : elle est la \textbf{dernière étape obligatoire} de la résolution.
\end{attention}

\courssub{Exploiter une courbe : images, antécédents et validation d'un calcul}

Une courbe représentative $\mathcal{C}_f$ est un outil de \emph{contrôle} et de \emph{communication} précieux, notamment en situation d'examen ou en contexte professionnel (lecture de plans, de courbes de charge, etc.).

\begin{itemize}
\item \textbf{Lecture d'une image (calculer $f(a)$) :} On repère $a$ sur l'axe des abscisses, on monte (ou descend) verticalement jusqu'à la courbe, puis on lit l'ordonnée correspondante sur l'axe des ordonnées.
\item \textbf{Lecture d'antécédents (résoudre $f(x)=k$) :} On repère $k$ sur l'axe des ordonnées, on trace l'horizontale $y=k$ jusqu'à la courbe, puis on lit les abscisses des points d'intersection sur l'axe des abscisses.
\item \textbf{Validation d'un calcul algébrique :} Après un calcul (discriminant, factorisation, formule des racines), on compare les résultats avec la courbe.
\begin{itemize}
\item Si le calcul donne deux racines distinctes, la courbe doit couper l'axe des abscisses en deux points distincts.
\item Si le calcul donne une racine double ($\Delta=0$), la courbe doit être tangente à l'axe des abscisses.
\item Si le calcul donne $\Delta<0$ (pas de racine), la courbe ne doit pas couper l'axe des abscisses.
\end{itemize}
En cas de contradiction, il y a une erreur de calcul ou de lecture à identifier.
\end{itemize}

\begin{exemplebox}[Contrôle graphique d'un calcul : Modélisation d'un bénéfice]
Un atelier de menuiserie modélise son bénéfice journalier $B$ (en milliers de francs CFA) en fonction du nombre $x$ d'articles vendus par la fonction $B(x)=-x^2+8x-7$, pour $x\in[0\,;\,8]$.
\begin{enumerate}
\item \textbf{Algébrique :} $B(x)=0 \iff -x^2+8x-7=0 \iff x^2-8x+7=0$. $\Delta=64-28=36$. Racines : $x_1=1$, $x_2=7$.
Le trinôme $-x^2+8x-7$ (coefficient dominant négatif) est positif \textbf{entre} ses racines. L'atelier est bénéficiaire pour $1<x<7$.
\item \textbf{Graphique :} On trace la parabole $\mathcal{C}_B$. Elle coupe l'axe des abscisses en $1$ et $7$, son sommet est en $x=4$ (maximum). Elle est au-dessus de l'axe des abscisses entre $1$ et $7$.
\item \textbf{Interprétation concrète :} Comme $x$ est un nombre entier d'articles, l'atelier est bénéficiaire pour une production de $2, 3, 4, 5$ ou $6$ articles par jour. Le graphique confirme le calcul et donne du sens au résultat (seuil de rentabilité).
\end{enumerate}
\end{exemplebox}

\begin{savaistu}[Le changement de variable, une technique transversale]
Le changement de variable utilisé pour les équations bicarrées ($X=x^2$) est l'une des techniques les plus puissantes des mathématiques. En classe de Terminale technique, tu le réutiliseras systématiquement pour résoudre des équations avec des exponentielles (poser $X=e^x$) ou des logarithmes (poser $X=\ln x$). L'idée est toujours la même : remplacer une expression compliquée par une lettre simple ($X$) pour se ramener à une équation que l'on sait résoudre (souvent du second degré)... puis \textbf{ne pas oublier de revenir à l'inconnue de départ} ($x$) en respectant les conditions sur $X$ (positivité pour $e^x$, réel pour $\ln x$, etc.).
\end{savaistu}

\courssub{Modélisation et résolution de problèmes concrets (Situations professionnelles)}

Le programme insiste sur l'application à des « situations concrètes de la vie courante ». Voici la démarche type attendue au Probatoire.

\begin{methode}[Démarche de modélisation : du texte à l'équation/inéquation]
\begin{enumerate}
\item \textbf{Lire et comprendre} : Identifier les grandeurs variables, l'inconnue (souvent $x$), les contraintes (domaine de définition, positivité, entiers...).
\item \textbf{Traduire} : Écrire la relation mathématique (équation $f(x)=k$, $f(x)=g(x)$ ou inéquation $f(x)<k$, etc.) à partir du texte.
\item \textbf{Résoudre} : Appliquer les méthodes algébriques (discriminant, factorisation, changement de variable, élévation au carré + vérification) ou graphiques.
\item \textbf{Interpréter} : Revenir au contexte. La solution mathématique a-t-elle du sens ? (ex: $x=-3$ mètres impossible, $x=12,5$ articles impossible si $x$ entier). Arrondir ou sélectionner les valeurs admissibles.
\item \textbf{Conclure} : Rédiger une phrase réponse complète, avec unités.
\end{enumerate}
\end{methode}

\begin{exoresolu}[Problème concret : Dimensionnement d'un réservoir]
Un réservoir rectangulaire à ciel ouvert doit avoir une base carrée de côté $x$ (en mètres) et une hauteur de $(x-2)$ mètres. Sa capacité doit être de $16\,\text{m}^3$. Déterminer les dimensions du réservoir.
\tcblower
\textbf{Solution.}
\begin{enumerate}
\item \textbf{Variables :} $x$ = côté de la base (en m). Contrainte : $x>2$ (hauteur positive).
\item \textbf{Modèle :} Volume $V = \text{Base} \times \text{Hauteur} = x^2 \times (x-2)$.
Condition : $V=16 \Rightarrow x^2(x-2)=16$.
\item \textbf{Résolution :} $x^3-2x^2-16=0$. On cherche une racine évidente : $x=4$ convient ($64-32-16=0$).
Factorisation : $(x-4)(x^2+2x+4)=0$. Le trinôme $x^2+2x+4$ a $\Delta=4-16=-12<0$, pas de racine réelle.
Seule solution : $x=4$.
\item \textbf{Vérification contrainte :} $x=4 > 2$ (valide). Hauteur $= 4-2=2$ m.
\item \textbf{Conclusion :} Le réservoir a une base carrée de $4\,\text{m}$ de côté et une hauteur de $2\,\text{m}$.
\end{enumerate}
\end{exoresolu}

\begin{exercice}[Exercices d'entraînement — Probatoire]
\begin{enumerate}
\item \textbf{Graphique et algébrique :} Soit $f(x)=x^2-4x+3$. 
\begin{enumerate}
\item Tracer la parabole $\mathcal{C}_f$ dans un repère (on pourra utiliser le tableau de valeurs : $x=0,1,2,3,4$).
\item Résoudre graphiquement $f(x)=0$ et $f(x)=-1$.
\item Résoudre algébriquement ces deux équations et comparer.
\end{enumerate}
\item \textbf{Équation bicarrée :} Résoudre dans $\mathbb{R}$ : $x^4-13x^2+36=0$.
\item \textbf{Équation avec racine :} Résoudre $\sqrt{2x+1}=x-1$ en vérifiant soigneusement les solutions.
\item \textbf{Inéquation graphique :} On donne les courbes $\mathcal{C}_f$ (parabole $y=x^2-2x-1$) et $\mathcal{C}_g$ (droite $y=2$) de la figure page 2. Déduire graphiquement la solution de $f(x) \geq g(x)$.
\item \textbf{Modélisation (Électricité) :} La puissance $P$ (en watts) dissipée par une résistance variable $R$ (en ohms) branchée sur un générateur de tension $U=12\,\text{V}$ et de résistance interne $r=2\,\Omega$ est donnée par $P(R)=\frac{144R}{(R+2)^2}$. On veut que la puissance soit supérieure ou égale à $16\,\text{W}$.
\begin{enumerate}
\item Traduire la condition par une inéquation en $R$.
\item Résoudre cette inéquation (on multipliera par $(R+2)^2>0$).
\item Interpréter le résultat pour le choix de la résistance $R$.
\end{enumerate}
\end{enumerate}
\end{exercice}

\begin{corrige}[Corrigé des exercices d'entraînement]
\textbf{1.} $f(x)=x^2-4x+3$.
\begin{itemize}
\item[a)] Tableau : $x=0\to3$, $1\to0$, $2\to-1$, $3\to0$, $4\to3$. Sommet en $x=2$, $y=-1$. Parabole convexe (branches vers le haut).
\item[b)] Graphiquement : $f(x)=0$ coupe l'axe $x$ en $1$ et $3$. $f(x)=-1$ : horizontale $y=-1$ tangente au sommet en $x=2$ (solution double).
\item[c)] Algébrique : $x^2-4x+3=0 \Rightarrow \Delta=16-12=4$, $x=1, 3$. $x^2-4x+3=-1 \Rightarrow x^2-4x+4=0 \Rightarrow (x-2)^2=0 \Rightarrow x=2$ (double). Cohérence parfaite.
\end{itemize}

\textbf{2.} $x^4-13x^2+36=0$. Poser $X=x^2\geq 0$. $X^2-13X+36=0$. $\Delta=169-144=25$. $X_1=\frac{13-5}{2}=4$, $X_2=\frac{13+5}{2}=9$. Les deux $>0$.
$x^2=4 \Rightarrow x=\pm 2$. $x^2=9 \Rightarrow x=\pm 3$. $S=\{-3\,;\,-2\,;\,2\,;\,3\}$.

\textbf{3.} $\sqrt{2x+1}=x-1$. Condition implicite : $x-1\geq 0 \Rightarrow x\geq 1$ et $2x+1\geq 0 \Rightarrow x\geq -0,5$ (donc $x\geq 1$).
Élévation au carré : $2x+1=x^2-2x+1 \Rightarrow x^2-4x=0 \Rightarrow x(x-4)=0 \Rightarrow x=0$ ou $x=4$.
Vérification : $x=0$ ne satisfait pas $x\geq 1$ (ou $\sqrt{1}=1 \neq -1$) $\Rightarrow$ \textbf{parasite, rejetée}. $x=4$ : $\sqrt{9}=3=4-1$ $\Rightarrow$ \textbf{acceptée}. $S=\{4\}$.

\textbf{4.} $f(x) \geq g(x) \iff x^2-2x-1 \geq 2 \iff x^2-2x-3 \geq 0$.
Graphiquement : $\mathcal{C}_f$ est au-dessus ou sur $\mathcal{C}_g$ en dehors de l'intervalle entre les intersections.
Solutions : $x\leq -1$ ou $x\geq 3$. $S=]-\infty\,;\,-1] \cup [3\,;\,+\infty[$.

\textbf{5.} 
\begin{enumerate}
\item $P(R) \geq 16 \iff \frac{144R}{(R+2)^2} \geq 16$.
\item $(R+2)^2>0$ pour tout $R\neq -2$ (et $R>0$ physiquement). On multiplie :
$144R \geq 16(R+2)^2 \iff 144R \geq 16(R^2+4R+4) \iff 144R \geq 16R^2+64R+64$.
$0 \geq 16R^2 -80R +64 \iff 16R^2 -80R +64 \leq 0$.
Diviser par $16$ : $R^2 -5R +4 \leq 0$.
Trinôme $a=1>0$, racines : $\Delta=25-16=9$, $R_1=1$, $R_2=4$.
Signe négatif entre les racines : $S=[1\,;\,4]$.
\item La résistance $R$ doit être comprise entre $1\,\Omega$ et $4\,\Omega$ (bornes incluses) pour que la puissance dissipée soit au moins de $16\,\text{W}$.
\end{enumerate}
\end{corrige}

\coursec{Modélisation et résolution de problèmes concrets}

Dans la section précédente, nous avons appris à résoudre des équations et des inéquations associées aux fonctions : nous savons désormais trouver les racines d'un trinôme, étudier son signe et résoudre un système de deux équations à deux inconnues. Mais à quoi servent ces outils dans la vie d'un technicien ? Un menuisier qui fabrique une table, un commerçant qui calcule son bénéfice, un groupe d'élèves qui achète des fournitures au marché : tous rencontrent des problèmes concrets qui se traduisent par des équations. Cette dernière section est la plus importante de la leçon : elle montre comment passer \emph{du problème réel aux mathématiques}, puis \emph{des mathématiques à la réponse concrète}. C'est exactement la démarche exigée au Probatoire technique dans les problèmes de synthèse.

\courssub{La démarche de modélisation en quatre étapes}

\begin{definition}[Modéliser un problème]
\cle{Modéliser} un problème concret, c'est le traduire en langage mathématique (choix d'inconnues, équations ou inéquations) afin de le résoudre avec les outils du cours, puis de revenir à la situation réelle pour interpréter le résultat.
\end{definition}

\begin{methode}[Les quatre étapes de la modélisation]
Pour résoudre un problème concret, on suit toujours la même démarche :
\begin{enumerate}
\item \textbf{Choix de l'inconnue.} On lit attentivement l'énoncé, on repère la (ou les) quantité cherchée et on la nomme, par exemple $x$. On précise ce qu'elle représente et son unité (mètres, francs CFA, nombre d'objets...).
\item \textbf{Mise en équation.} On traduit chaque information de l'énoncé par une relation mathématique : une équation du second degré, une inéquation, ou un système de deux équations.
\item \textbf{Résolution.} On résout l'équation, l'inéquation ou le système avec les méthodes du cours (discriminant, tableau de signes, substitution ou combinaison).
\item \textbf{Interprétation et validation.} On confronte chaque solution mathématique au contexte : une longueur doit être positive, un nombre d'objets doit être entier, un prix doit être réaliste. On \cle{rejette} les solutions non adaptées, puis on rédige une conclusion claire avec les unités.
\end{enumerate}
\end{methode}

\begin{popfigure}
\begin{center}
\begin{tikzpicture}[node distance=0.55cm,
etape/.style={rectangle, rounded corners=3pt, draw=popblue, fill=popblueL, text width=4.6cm, align=center, minimum height=1cm, font=\small},
fleche/.style={-{Stealth[length=2.5mm]}, thick, poporange}]
\node[etape] (e1) {\textbf{Étape 1} : choisir l'inconnue $x$ (préciser ce qu'elle représente et son unité)};
\node[etape, below=of e1] (e2) {\textbf{Étape 2} : traduire l'énoncé en équation, inéquation ou système};
\node[etape, below=of e2] (e3) {\textbf{Étape 3} : résoudre avec les outils du cours ($\Delta$, signe du trinôme, système $2\times 2$)};
\node[etape, below=of e3] (e4) {\textbf{Étape 4} : interpréter, rejeter les solutions impossibles, conclure avec les unités};
\draw[fleche] (e1) -- (e2);
\draw[fleche] (e2) -- (e3);
\draw[fleche] (e3) -- (e4);
\draw[fleche, dashed, popmuted] (e4.east) to[out=0,in=0] node[right, font=\scriptsize, align=left, popmuted]{si le résultat est\\ incohérent : revoir\\ le modèle} (e1.east);
\end{tikzpicture}
\end{center}
\legende{Organigramme des quatre étapes de la démarche de modélisation.}
\end{popfigure}

\begin{attention}[Une solution mathématique n'est pas toujours une solution du problème]
L'erreur la plus fréquente au Probatoire est de donner une solution mathématiquement correcte sans vérifier sa cohérence avec le contexte. Une équation du second degré donne souvent \emph{deux} solutions : si $x$ représente une longueur, une solution négative doit être \textbf{rejetée} ; si $x$ représente un nombre de personnes ou d'articles, une solution non entière doit être rejetée. Écris toujours explicitement : « cette solution est rejetée car... ». C'est cette justification qui rapporte des points.
\end{attention}

\courssub{Problème d'aire et de dimensions : la table de M. Nkoulou}

Les problèmes de géométrie (aires, dimensions) conduisent naturellement à des équations du second degré, car l'aire d'un rectangle est le \emph{produit} de deux dimensions.

\begin{exoresolu}[La table de M. Nkoulou, menuisier à Yaoundé]
M. Nkoulou, menuisier, doit fabriquer un plateau de table rectangulaire dont la longueur mesure $1$ mètre de plus que la largeur. La surface du plateau doit être exactement de $2$ mètres carrés. Déterminer les dimensions du plateau.
\tcblower
\textbf{Étape 1 : choix de l'inconnue.} Soit $x$ la largeur du plateau, en mètres ($x > 0$). Alors la longueur est $x + 1$.

\textbf{Étape 2 : mise en équation.} L'aire du rectangle vaut largeur $\times$ longueur :
\[ x(x+1) = 2 \quad\Longleftrightarrow\quad x^2 + x - 2 = 0. \]

\textbf{Étape 3 : résolution.} Trinôme avec $a = 1$, $b = 1$, $c = -2$ :
\[ \Delta = b^2 - 4ac = 1^2 - 4 \times 1 \times (-2) = 1 + 8 = 9 > 0. \]
Il y a deux solutions :
\[ x_1 = \frac{-1 - \sqrt{9}}{2} = \frac{-1-3}{2} = -2 \qquad \text{et} \qquad x_2 = \frac{-1 + 3}{2} = 1. \]

\textbf{Étape 4 : interprétation et validation.} $x_1 = -2$ est \textbf{rejetée} car une largeur ne peut pas être négative. On conserve $x = 1$.

\textbf{Conclusion.} Le plateau a une largeur de $1$ m et une longueur de $1 + 1 = 2$ m. Vérification : $1 \times 2 = 2$ m$^2$ : c'est bien la surface demandée.
\end{exoresolu}

\begin{popfigure}
\begin{center}
\begin{tikzpicture}[scale=1.6]
\fill[poporangeL] (0,0) rectangle (2,1);
\draw[thick, popdark] (0,0) rectangle (2,1);
\draw[{Stealth}-{Stealth}, popblue, thick] (0,-0.25) -- (2,-0.25) node[midway, below, font=\small]{$x+1$ (longueur)};
\draw[{Stealth}-{Stealth}, popblue, thick] (-0.25,0) -- (-0.25,1) node[midway, left, font=\small]{$x$ (largeur)};
\node[font=\small, popdark] at (1,0.5) {Aire $= x(x+1) = 2$ m$^2$};
\end{tikzpicture}
\end{center}
\legende{Le plateau rectangulaire de la table : largeur $x$, longueur $x+1$, aire $2$ m$^2$.}
\end{popfigure}

\begin{methode}[Résoudre un problème d'aire]
\begin{enumerate}
\item Nommer la dimension inconnue $x$ (en précisant $x > 0$).
\item Exprimer les autres dimensions en fonction de $x$.
\item Écrire la formule d'aire, développer et ramener à la forme $ax^2 + bx + c = 0$.
\item Calculer $\Delta$, trouver les racines, \textbf{rejeter toute racine négative ou nulle}.
\item Vérifier en calculant l'aire avec les dimensions trouvées.
\end{enumerate}
\end{methode}

\courssub{Problème de rentabilité : une inéquation du second degré}

En gestion, on cherche souvent à savoir \emph{à partir de quand} une activité devient rentable. Le bénéfice dépend alors de la quantité produite, et la condition « bénéfice positif » se traduit par une inéquation du second degré.

\begin{exoresolu}[L'atelier de couture de Mme Etoundi]
Mme Etoundi fabrique et vend des sacs. Pour $x$ sacs fabriqués et vendus par semaine, son bénéfice hebdomadaire, en milliers de francs CFA, est donné par
\[ B(x) = -x^2 + 14x - 40. \]
Pour quelles valeurs de $x$ l'atelier est-il rentable (bénéfice strictement positif) ?
\tcblower
\textbf{Étape 1 : inconnue.} $x$ est le nombre de sacs, un entier positif.

\textbf{Étape 2 : mise en inéquation.} L'atelier est rentable si $B(x) > 0$, c'est-à-dire :
\[ -x^2 + 14x - 40 > 0. \]

\textbf{Étape 3 : résolution.} On cherche les racines du trinôme ($a = -1$, $b = 14$, $c = -40$) :
\[ \Delta = 14^2 - 4 \times (-1) \times (-40) = 196 - 160 = 36 > 0, \qquad \sqrt{\Delta} = 6. \]
\[ x_1 = \frac{-14 - 6}{-2} = 10 \qquad \text{et} \qquad x_2 = \frac{-14 + 6}{-2} = 4. \]
Comme $a = -1 < 0$, le trinôme est \textbf{positif entre ses racines} :
\[ B(x) > 0 \quad\Longleftrightarrow\quad x \in \left]4\,;\,10\right[. \]

\textbf{Étape 4 : interprétation.} $x$ est un nombre entier de sacs : $x \in \{5, 6, 7, 8, 9\}$.

\textbf{Conclusion.} Mme Etoundi doit fabriquer et vendre entre $5$ et $9$ sacs par semaine pour que son atelier soit rentable. En dessous de $5$ sacs, les ventes ne couvrent pas les frais ; au-delà de $9$ sacs, les coûts de production deviennent trop élevés. Vérification : $B(7) = -49 + 98 - 40 = 9$, soit un bénéfice de $\num{9000}$ FCFA : c'est bien positif.
\end{exoresolu}

\begin{attention}[Inéquation : ne pas oublier le sens des inégalités]
Deux pièges classiques : oublier que le signe du trinôme dépend du signe de $a$ (ici $a < 0$, donc le trinôme est positif \emph{entre} les racines), et donner comme réponse l'intervalle $]4\,;\,10[$ alors que le contexte exige des \textbf{nombres entiers}. Toujours adapter l'ensemble des solutions au contexte.
\end{attention}

\courssub{Problème d'achats groupés : un système $2\times 2$}

Quand un énoncé fait intervenir \emph{deux} quantités inconnues liées par \emph{deux} informations, on utilise un système de deux équations à deux inconnues.

\begin{exoresolu}[Achat groupé au marché Mokolo]
Deux élèves de Première technique, Amina et Basile, achètent ensemble des fournitures au marché Mokolo à Yaoundé. Amina achète $3$ cahiers et $2$ stylos pour $\num{1700}$ FCFA. Basile achète $2$ cahiers et $4$ stylos pour $\num{1800}$ FCFA. Quel est le prix d'un cahier et celui d'un stylo ?
\tcblower
\textbf{Étape 1 : choix des inconnues.} Soit $x$ le prix d'un cahier et $y$ le prix d'un stylo, en francs CFA.

\textbf{Étape 2 : mise en équations.} Chaque achat donne une équation :
\[ \begin{cases} 3x + 2y = 1700 \\ 2x + 4y = 1800 \end{cases} \]

\textbf{Étape 3 : résolution par combinaison.} On multiplie la première équation par $2$ :
\[ \begin{cases} 6x + 4y = 3400 \\ 2x + 4y = 1800 \end{cases} \]
On soustrait membre à membre : $6x - 2x = 3400 - 1800$, donc $4x = 1600$, d'où $x = 400$.
On reporte dans la première équation :
\[ 3 \times 400 + 2y = 1700 \quad\Longleftrightarrow\quad 2y = 500 \quad\Longleftrightarrow\quad y = 250. \]

\textbf{Étape 4 : validation.} Les deux prix sont positifs et réalistes. Vérification avec l'achat de Basile : $2 \times 400 + 4 \times 250 = 800 + 1000 = 1800$ FCFA : c'est exact.

\textbf{Conclusion.} Un cahier coûte $400$ FCFA et un stylo coûte $250$ FCFA.
\end{exoresolu}

\begin{methode}[Résoudre un problème à deux inconnues]
\begin{enumerate}
\item Repérer les deux quantités inconnues et les nommer $x$ et $y$.
\item Traduire \emph{chaque} phrase de l'énoncé en une équation (il faut autant d'équations que d'inconnues).
\item Résoudre le système par substitution ou par combinaison.
\item Vérifier le couple solution dans \textbf{les deux} équations, puis vérifier sa cohérence avec le contexte (prix positifs, quantités entières).
\end{enumerate}
\end{methode}

\courssub{Valider, rédiger et critiquer le modèle}

\begin{aretenir}[Rédiger la conclusion d'un problème]
Une conclusion complète contient toujours trois éléments :
\begin{itemize}
\item la \textbf{réponse} à la question posée, en phrase complète ;
\item les \textbf{unités} (mètres, FCFA, nombre d'objets...) ;
\item une \textbf{justification} : solutions rejetées et raison du rejet, vérification du résultat.
\end{itemize}
\end{aretenir}

\begin{exemplebox}[Critiquer la validité d'un modèle]
Reprenons l'atelier de Mme Etoundi : le modèle $B(x) = -x^2 + 14x - 40$ suppose que \emph{tous} les sacs fabriqués sont vendus au même prix et que les coûts suivent exactement cette formule. En réalité, si Mme Etoundi produit $20$ sacs, elle ne les vendra peut-être pas tous : le modèle devient alors faux. De même, dans le problème de la table, on a supposé le plateau parfaitement rectangulaire, sans tenir compte des chutes de bois. Un bon technicien sait donc \cle{critiquer} son modèle : il se demande si les hypothèses sont réalistes et si les résultats obtenus sont plausibles (un cahier à $400$ FCFA, oui ; un cahier à $\num{400000}$ FCFA signalerait une erreur de calcul).
\end{exemplebox}

\begin{exercice}[À toi de jouer]
Un électricien doit délimiter un local rectangulaire dont le périmètre est de $26$ m et l'aire de $40$ m$^2$. On note $x$ et $y$ les dimensions du local.
\begin{enumerate}
\item Montrer que $x + y = 13$ et $xy = 40$.
\item En déduire que $x$ est solution de l'équation $x^2 - 13x + 40 = 0$.
\item Résoudre cette équation et conclure sur les dimensions du local.
\end{enumerate}
\end{exercice}

\begin{corrige}[Exercice : le local de l'électricien]
\begin{enumerate}
\item Le périmètre vaut $2(x+y) = 26$, donc $x + y = 13$. L'aire vaut $xy = 40$.
\item De $x + y = 13$ on tire $y = 13 - x$. En reportant : $x(13 - x) = 40$, soit $13x - x^2 = 40$, c'est-à-dire $x^2 - 13x + 40 = 0$.
\item $\Delta = (-13)^2 - 4 \times 40 = 169 - 160 = 9$, $\sqrt{\Delta} = 3$. D'où $x_1 = \dfrac{13 - 3}{2} = 5$ et $x_2 = \dfrac{13 + 3}{2} = 8$. Les deux solutions sont positives, donc toutes deux acceptables : si $x = 5$ alors $y = 8$, et si $x = 8$ alors $y = 5$. \textbf{Conclusion :} le local mesure $5$ m sur $8$ m. Vérification : périmètre $2(5+8) = 26$ m et aire $5 \times 8 = 40$ m$^2$.
\end{enumerate}
\end{corrige}

\begin{savaistu}[Des équations vieilles de 4000 ans]
Les équations du second degré ne sont pas une invention moderne ! Il y a près de $\num{4000}$ ans, les mathématiciens \textbf{babyloniens} (dans l'actuel Irak) résolvaient déjà des équations du type $x^2 + bx = c$ pour des problèmes très concrets : partage de champs, calcul de surfaces de terres agricoles, remboursement de prêts avec intérêts. Ils utilisaient des tablettes d'argile gravées et des méthodes géométriques de « complétion du carré », ancêtre de notre forme canonique. En Afrique, les scribes de l'Égypte ancienne (papyrus de Rhind, vers $-\num{1650}$) résolvaient aussi des problèmes de partage de récoltes à l'aide d'équations. Les mathématiques que tu apprends aujourd'hui sont donc le fruit d'une longue histoire, née de besoins pratiques identiques à ceux des artisans et commerçants camerounais d'aujourd'hui.
\end{savaistu}

\begin{aretenir}[L'essentiel de la section]
\begin{itemize}
\item Modéliser, c'est suivre quatre étapes : \cle{choix de l'inconnue}, \cle{mise en équation}, \cle{résolution}, \cle{interprétation et validation}.
\item Les problèmes d'aire conduisent à une équation du second degré ; les problèmes de rentabilité à une inéquation du second degré ; les problèmes d'achats à deux inconnues à un système $2\times 2$.
\item Toute solution incohérente avec le contexte (longueur négative, quantité non entière) doit être \textbf{rejetée avec justification}.
\item La conclusion se rédige en phrase, avec les unités, et s'accompagne d'une vérification.
\end{itemize}
\end{aretenir}

\newpage

\coursec{Activité d'intégration}

\coursec{Leçon 1 : Équations, inéquations et systèmes}

\courssub{Objectifs de l'activité}
\begin{itemize}
    \item Mobiliser les connaissances sur les trinômes du second degré (forme canonique, factorisation, signe, résolution d'équations et d'inéquations) dans une situation concrète d'optimisation.
    \item Résoudre un système linéaire à deux inconnues pour comparer deux offres commerciales.
    \item Rédiger une démarche mathématique complète, justifiée et structurée, adaptée à la communication professionnelle (compte-rendu pour un décideur).
    \item Faire le lien entre le modèle mathématique (fonction aire, fonction coût) et la réalité du terrain (contraintes physiques, budget).
\end{itemize}

\courssub{Situation}
\begin{popfigure}
\begin{tikzpicture}[scale=0.05, >=Stealth]
\def\L{50}
\def\x{25}
% Fleuve
\draw[popblue, very thick, decorate, decoration={snake, amplitude=3pt, segment length=15pt}] (-5, \x+2) -- (\L+5, \x+2) node[midway, above=10pt, popblue] {\textbf{Fleuve Sanaga}};
% Champ
\draw[popdark, thick] (0,0) rectangle (\L, \x);
% Grillage (3 côtés)
\draw[popgreen, ultra thick] (0,0) -- (0,\x);
\draw[popgreen, ultra thick] (\L,0) -- (\L,\x);
\draw[popgreen, ultra thick] (0,0) -- (\L,0);
% Portail
\draw[poppink, ultra thick] (\L/2 - 1.5, 0) -- (\L/2 + 1.5, 0);
\node[poppink, below=5pt] at (\L/2, 0) {\small Portail (3 m)};
% Cotes
\draw[<->, poporange, thick] (-3, 0) -- (-3, \x) node[midway, left, poporange] {$x$ (largeur)};
\draw[<->, poporange, thick] (0, -3) -- (\L, -3) node[midway, below, poporange] {$L$ (longueur)};
% Légendes internes
\node[popdark, align=center] at (\L/2, \x/2) {Champ cultivable\\Maïs \& Manioc};
\node[popgreen, rotate=90, anchor=south] at (-5, \x/2) {Grillage 100 m};
\node[popgreen, anchor=south] at (\L/2, -5) {Grillage 100 m};
\node[popgreen, rotate=-90, anchor=north] at (\L+5, \x/2) {Grillage 100 m};
\end{tikzpicture}
\legende{Modèle géométrique : rectangle de largeur $x$ et longueur $L$, bordé par le fleuve sur un grand côté. Le grillage (100 m) clôture les deux largeurs et la longueur opposée.}
\end{popfigure}

La coopérative agricole \textbf{« Espoir d'Edéa »}, située dans la région du Littoral, souhaite aménager un champ rectangulaire pour la culture du maïs et du manioc. Le terrain est bordé sur l'un de ses grands côtés par le fleuve Sanaga, ce qui dispense de clôture sur cette rive. La coopérative dispose d'un stock de \textbf{100 mètres de grillage rigide} pour clôturer les trois autres côtés (les deux largeurs et la longueur opposée au fleuve).

Le président de la coopérative, M. \textsc{Ngom}, doit prendre deux décisions majeures avant le début de la saison des pluies :
\begin{enumerate}
    \item \textbf{Dimensionnement :} Quelles dimensions (largeur et longueur) choisir pour \textbf{maximiser la surface cultivable} ? L'aire doit-elle impérativement dépasser \textbf{1\,000 m²} pour rentabiliser l'achat d'intrants (engrais, semences améliorées) ?
    \item \textbf{Budget clôture :} Le grillage coûte \textbf{2\,500 FCFA le mètre linéaire} et l'installation d'un portail d'accès (largeur 3 m, inclus dans la longueur clôturée) revient à un forfait de \textbf{15\,000 FCFA} (fourniture et pose). Un fournisseur concurrent, \textit{Clôture Express}, propose un tarif global différent. M. \textsc{Ngom} a obtenu deux devis de ce concurrent pour des configurations types :
    \begin{itemize}
        \item \textit{Devis A :} 40 m de grillage + 1 portail = 115\,000 FCFA.
        \item \textit{Devis B :} 60 m de grillage + 1 portail = 165\,000 FCFA.
    \end{itemize}
    Il souhaite identifier le prix unitaire du mètre de grillage et le forfait portail chez ce concurrent pour comparer les offres.
\end{enumerate}

\courssub{Consignes}
\begin{enumerate}
    \item \textbf{Modélisation de l'aire.} Soit $x$ la largeur du champ (en mètres), perpendiculaire au fleuve. Exprimer la longueur $L$ en fonction de $x$, puis montrer que l'aire $A(x)$ du champ s'écrit $A(x) = -2x^2 + 100x$. Préciser l'ensemble de définition physique de $x$.
    \item \textbf{Maximisation de l'aire.} Mettre $A(x)$ sous forme canonique. En déduire la largeur $x_{\text{max}}$ qui maximise l'aire, la longueur correspondante $L_{\text{max}}$ et l'aire maximale $A_{\text{max}}$.
    \item \textbf{Contrainte de rentabilité (Inéquation).} Résoudre l'inéquation $A(x) > 1000$. Interpréter le résultat : entre quelles valeurs de largeur l'aire dépasse-t-elle 1\,000 m² ? Le maximum trouvé en question 2 satisfait-il cette contrainte ?
    \item \textbf{Analyse comparative des coûts (Système linéaire).} Soit $p$ le prix du mètre de grillage (en FCFA) et $f$ le forfait portail (en FCFA) chez \textit{Clôture Express}. 
    \begin{enumerate}
        \item Traduire les deux devis en un système de deux équations à deux inconnues $(p; f)$.
        \item Résoudre ce système par la méthode de votre choix (substitution, combinaison, ou Cramer).
        \item Calculer le coût total de la clôture du champ optimal (questions 1-2) chez \textit{Clôture Express} et chez le fournisseur initial. Lequel est le plus avantageux ?
    \end{enumerate}
    \item \textbf{Synthèse et recommandation.} Rédiger une note argumentée à l'attention de M. \textsc{Ngom}, résumant les dimensions optimales, la faisabilité de l'objectif de 1\,000 m², et le choix du fournisseur le plus économique, en justifiant chaque chiffre par les calculs effectués.
\end{enumerate}

\courssub{Ressources à mobiliser}
\begin{itemize}
    \item \textbf{Trinôme du second degré :} Forme canonique $a\left[\left(x + \frac{b}{2a}\right)^2 - \frac{\Delta}{4a^2}\right]$ avec $\Delta = b^2 - 4ac$. Sommet de la parabole en $\left(-\frac{b}{2a}; -\frac{\Delta}{4a}\right)$.
    \item \textbf{Signe du trinôme :} Pour $a < 0$, $ax^2+bx+c > 0 \iff x \in ]x_1; x_2[$ (racines $x_1 < x_2$). Pour $a > 0$, $ax^2+bx+c < 0 \iff x \in ]x_1; x_2[$.
    \item \textbf{Système linéaire $2 \times 2$ :} $\begin{cases} a_1 x + b_1 y = c_1 \\ a_2 x + b_2 y = c_2 \end{cases}$. Méthodes : combinaison (élimination), substitution, déterminants ($\Delta = a_1b_2 - a_2b_1$).
    \item \textbf{Contexte géométrique :} Périmètre partiel $P = L + 2x = 100$. Aire $A = L \times x$.
\end{itemize}

\courssub{Démarche et corrigé commenté}
\begin{exoresolu}{Champ rectangulaire de la coopérative d'Edéa}
\textbf{Rappel de l'énoncé :} Un champ rectangulaire de largeur $x$ et longueur $L$ est bordé par un fleuve sur une longueur. Le périmètre clôturé (2 largeurs + 1 longueur) est de 100 m. On cherche à maximiser l'aire $A(x)$, à garantir $A(x) > 1000$ m², et à comparer deux fournisseurs de clôture via un système linéaire.
\tcblower
\textbf{1. Modélisation de l'aire $A(x)$}
Le grillage disponible est de 100 m. Il sert à clôturer les deux largeurs ($x$ chacune) et la longueur $L$ opposée au fleuve.
\[
L + 2x = 100 \quad \Rightarrow \quad L = 100 - 2x.
\]
L'aire du rectangle est $A = L \times x$. En remplaçant $L$ :
\[
A(x) = (100 - 2x)x = 100x - 2x^2 = -2x^2 + 100x.
\]
\emph{Contraintes physiques :} Une longueur et une largeur sont strictement positives.
$L > 0 \iff 100 - 2x > 0 \iff x < 50$.
$x > 0$.
L'ensemble de définition physique est donc \textbf{$D_A = ]0; 50[$} (intervalle ouvert car si $x=0$ ou $x=50$, l'aire est nulle, il n'y a pas de champ).

\textbf{2. Forme canonique et maximisation}
$A(x) = -2x^2 + 100x$. On factorise par $-2$ :
\[
A(x) = -2(x^2 - 50x).
\]
On complète le carré : $x^2 - 50x = (x - 25)^2 - 625$.
\[
A(x) = -2\left[(x - 25)^2 - 625\right] = -2(x - 25)^2 + 1250.
\]
C'est la forme canonique. Le coefficient de $(x-25)^2$ est $-2 < 0$, la parabole est tournée vers le bas : le sommet est un \textbf{maximum}.
Le sommet a pour coordonnées $(25; 1250)$.
\begin{itemize}
    \item Largeur optimale : \textbf{$x_{\text{max}} = 25$ m}.
    \item Longueur correspondante : $L_{\text{max}} = 100 - 2 \times 25 = \textbf{50 m}$.
    \item Aire maximale : \textbf{$A_{\text{max}} = 1250$ m²}.
\end{itemize}
\emph{Vérification :} $25 \in ]0; 50[$, la solution est admissible.

\textbf{3. Inéquation $A(x) > 1000$ (Objectif de rentabilité)}
On résout $-2x^2 + 100x > 1000$.
On ramène à 0 : $-2x^2 + 100x - 1000 > 0$.
On divise par $-2$ (ce qui change le sens de l'inégalité) :
\[
x^2 - 50x + 500 < 0.
\]
Calcul du discriminant : $\Delta = (-50)^2 - 4 \times 1 \times 500 = 2500 - 2000 = 500$.
$\sqrt{\Delta} = \sqrt{500} = 10\sqrt{5} \approx 22,36$.
Les deux racines sont :
\[
x_1 = \frac{50 - 10\sqrt{5}}{2} = 25 - 5\sqrt{5} \approx 13,82,
\quad
x_2 = \frac{50 + 10\sqrt{5}}{2} = 25 + 5\sqrt{5} \approx 36,18.
\]
Le trinôme $x^2 - 50x + 500$ a $a=1 > 0$, il est strictement négatif \textbf{entre} ses racines.
L'ensemble des solutions de l'inéquation initiale est donc :
\[
S = ]25 - 5\sqrt{5};\; 25 + 5\sqrt{5}[ \approx ]13,82;\; 36,18[.
\]
\emph{Interprétation :} Pour que l'aire dépasse 1\,000 m², la largeur $x$ doit être comprise entre \textbf{13,8 m et 36,2 m} (arrondi au dixième).
Le maximum $x_{\text{max}} = 25$ m appartient bien à cet intervalle ($13,82 < 25 < 36,18$). L'objectif de rentabilité est \textbf{atteignable et même dépassé} au maximum (1250 m² > 1000 m²).

\textbf{4. Système linéaire : Tarification « Clôture Express »}
Soit $p$ le prix du mètre de grillage (FCFA/m) et $f$ le forfait portail (FCFA).
\begin{itemize}
    \item Devis A : $40p + f = 115\,000$
    \item Devis B : $60p + f = 165\,000$
\end{itemize}
Système :
\[
(S) : \begin{cases}
40p + f = 115\,000 & (L_1)\\
60p + f = 165\,000 & (L_2)
\end{cases}
\]
\emph{Méthode par combinaison (soustraction) :} $(L_2) - (L_1)$ élimine $f$ :
\[
(60p - 40p) + (f - f) = 165\,000 - 115\,000 \iff 20p = 50\,000 \iff p = 2\,500.
\]
On remplace $p$ dans $(L_1)$ :
\[
40 \times 2\,500 + f = 115\,000 \iff 100\,000 + f = 115\,000 \iff f = 15\,000.
\]
\textbf{Résultat :} Chez \textit{Clôture Express}, le mètre de grillage coûte \textbf{2\,500 FCFA} et le forfait portail \textbf{15\,000 FCFA}.
\emph{Remarque :} Les tarifs sont \textbf{identiques} à ceux du fournisseur initial.

\textbf{5. Comparaison des coûts totaux pour le champ optimal}
Le champ optimal nécessite 100 m de grillage (stock total) et 1 portail.
\begin{itemize}
    \item \textbf{Fournisseur initial :} Coût $C_1 = 100 \times 2\,500 + 15\,000 = 250\,000 + 15\,000 = \textbf{265\,000 FCFA}$.
    \item \textbf{Clôture Express :} Coût $C_2 = 100 \times 2\,500 + 15\,000 = \textbf{265\,000 FCFA}$.
\end{itemize}
Les deux offres sont \textbf{strictement équivalentes} financièrement pour cette configuration.

\textbf{6. Note de synthèse pour M. Ngom}
\begin{quote}
\textbf{Objet : Dimensionnement et budget du champ riverain – Campagne 2026}

Monsieur le Président,

Suite à l'étude technique du champ adossé au fleuve Sanaga, voici les éléments pour votre décision :
\begin{enumerate}
    \item \textbf{Dimensions optimales :} Pour maximiser la surface cultivable avec 100 m de grillage, il faut prévoir une \textbf{largeur de 25 m} et une \textbf{longueur de 50 m}. Cela donne une aire maximale de \textbf{1\,250 m²}.
    \item \textbf{Objectif de rentabilité (1\,000 m²) :} Cet objectif est \textbf{réalisable}. Toute largeur comprise entre \textbf{13,8 m et 36,2 m} garantit une aire supérieure à 1\,000 m². Le choix optimal (25 m) s'inscrit confortablement dans cette plage.
    \item \textbf{Choix du fournisseur :} L'analyse des devis de \textit{Clôture Express} révèle des tarifs identiques au fournisseur actuel (2\,500 FCFA/m et 15\,000 FCFA le portail). Le coût total de la clôture s'élève à \textbf{265\,000 FCFA} chez les deux. Il n'y a pas d'avantage financier à changer de prestataire sur la base des tarifs unitaires.
\end{enumerate}
\emph{Recommandation :} Valider les dimensions 25 m $\times$ 50 m. Négocier une remise de quantité ou un geste commercial (maintenance, livraison) auprès de l'un des deux fournisseurs pour réduire la facture finale.

Cordialement, \textit{Le Chargé d'Études Techniques}.
\end{quote}
\end{exoresolu}

\courssub{Bilan de l'activité}
Cette activité d'intégration a permis de :
\begin{itemize}
    \item \textbf{Ancrer les mathématiques dans le réel :} Le passage du périmètre imposé (100 m) à la fonction aire $A(x)$ illustre la modélisation par une fonction du second degré. La contrainte « adossé à la rivière » modifie la formule habituelle du périmètre ($P = 2L+2l$ devient $L+2l$), point de vigilance essentiel.
    \item \textbf{Articuler les trois piliers du chapitre :}
    \begin{enumerate}
        \item \emph{Forme canonique / Optimisation :} Trouver le maximum d'une parabole tournée vers le bas ($a<0$) sans dérivée, outil algébrique requis au Probatoire.
        \item \emph{Inéquation / Seuil de rentabilité :} Résoudre $A(x) > 1000$ impose de manipuler le signe du trinôme, de calculer un discriminant non parfait ($\Delta=500$), et d'interpréter l'intervalle solution dans le domaine physique $]0;50[$.
        \item \emph{Système linéaire / Décision d'achat :} La comparaison de devis conduit naturellement à un système $2 \times 2$ résolu par combinaison. La découverte de tarifs identiques est une situation pédagogique riche : elle montre que les mathématiques valident l'équivalence et orientent vers la négociation de services annexes.
    \end{enumerate}
    \item \textbf{Exiger une rédaction professionnelle :} La consigne 5 impose une conclusion structurée (dimensions, faisabilité, coût, recommandation), simulant une note technique. Cela travaille la compétence « Rédiger et justifier une démarche » du programme.
    \item \textbf{Valider les prérequis pour la suite :} La maîtrise de la forme canonique, du signe du trinôme et de la résolution de systèmes est indispensable pour la Leçon 2 (Fonctions du second degré, variations, optimisation graphique) et les chapitres de statistique/probabilités (régression linéaire, systèmes normaux).
\end{itemize}

\begin{attention}[Pièges évités dans cette activité]
\begin{itemize}
\item \textbf{Domaine de définition :} Ne pas écrire $D_A = [0; 50]$. À $x=0$ ou $x=50$, l'aire est nulle, il n'y a pas de rectangle (largeur ou longueur nulle). L'intervalle est ouvert $]0; 50[$.
\item \textbf{Sens de l'inégalité :} Diviser par $-2$ inverse le signe ($>$ devient $<$). Oublier ce changement conduit à l'intervalle complémentaire (erreur classique).
\item \textbf{Unités et arrondis :} Les dimensions sont en mètres (m), l'aire en m², les coûts en FCFA. Les racines exactes ($25 \pm 5\sqrt{5}$) sont privilégiées dans la démonstration ; les valeurs approchées (13,8 ; 36,2) servent l'interprétation terrain.
\item \textbf{Système linéaire :} Vérifier que les deux équations sont bien indépendantes (ici coefficients de $p$ différents : 40 et 60). Si les devis étaient proportionnels, le système n'aurait pas de solution unique (tarification incohérente).
\end{itemize}
\end{attention}

\newpage

\coursec{Méthodes et exercices résolus}

\courssub{Méthodes et exercices résolus}

\begin{methode}[Mettre un trinôme sous forme canonique et trouver son sommet]
    \textbf{Objectif :} Écrire $f(x)=ax^2+bx+c$ ($a \neq 0$) sous la forme $a(x-\alpha)^2+\beta$ pour lire le sommet $S(\alpha;\beta)$ et le sens de variation.
    \begin{enumerate}
        \item Identifier les coefficients $a$, $b$, $c$ ($a \neq 0$).
        \item Calculer l'abscisse du sommet : $\alpha = -\dfrac{b}{2a}$.
        \item Calculer l'ordonnée du sommet : $\beta = f(\alpha) = -\dfrac{\Delta}{4a}$ où $\Delta = b^2-4ac$.
        \item Écrire la forme canonique : $f(x) = a(x-\alpha)^2 + \beta$.
        \item \textbf{Interprétation technique :} Si $a>0$, la parabole est tournée vers le haut (minimum en $\beta$) ; si $a<0$, vers le bas (maximum en $\beta$).
    \end{enumerate}
\end{methode}

\begin{exoresolu}[Forme canonique et optimisation d'une section de poutre]
    La hauteur $h$ (en cm) d'une poutre en béton en fonction de la distance $x$ (en m) à l'appui est modélisée par $h(x) = -2x^2 + 8x + 10$ sur $[0;5]$.
    \begin{enumerate}
        \item Mettre $h(x)$ sous forme canonique.
        \item Déterminer la hauteur maximale de la poutre et l'abscisse où elle est atteinte.
    \end{enumerate}
    \tcblower
    \textbf{Solution détaillée.}
    \begin{enumerate}
        \item On a $a=-2$, $b=8$, $c=10$.
        \begin{align*}
            \alpha &= -\frac{b}{2a} = -\frac{8}{2\times(-2)} = -\frac{8}{-4} = 2 \\[4pt]
            \Delta &= b^2 - 4ac = 8^2 - 4\times(-2)\times10 = 64 + 80 = 144 \\[4pt]
            \beta &= -\frac{\Delta}{4a} = -\frac{144}{4\times(-2)} = -\frac{144}{-8} = 18
        \end{align*}
        Forme canonique : $h(x) = -2(x-2)^2 + 18$.
        \item Comme $a = -2 < 0$, la parabole est tournée vers le bas. Le sommet $S(2;18)$ correspond à un \textbf{maximum}.
        La hauteur maximale est \textbf{18 cm}, atteinte à \textbf{2 m} de l'appui (valeur appartenant à l'intervalle $[0;5]$).
    \end{enumerate}
    \textbf{Conclusion :} La poutre est la plus haute (18 cm) à 2 mètres de l'appui.
\end{exoresolu}

\begin{methode}[Résoudre une équation du second degré $ax^2+bx+c=0$]
    \textbf{Objectif :} Trouver l'ensemble des solutions $S$ dans $\mathbb{R}$.
    \begin{enumerate}
        \item Calculer le discriminant $\Delta = b^2 - 4ac$.
        \item \textbf{Discussion selon le signe de $\Delta$ :}
        \begin{itemize}
            \item $\Delta < 0$ : Aucune solution réelle ($S = \varnothing$).
            \item $\Delta = 0$ : Une solution double $x_0 = -\dfrac{b}{2a}$ ($S = \{x_0\}$).
            \item $\Delta > 0$ : Deux solutions distinctes 
            $x_1 = \dfrac{-b-\sqrt{\Delta}}{2a}$ et $x_2 = \dfrac{-b+\sqrt{\Delta}}{2a}$ ($S = \{x_1; x_2\}$).
            \emph{Remarque : si $a>0$ alors $x_1 < x_2$ ; si $a<0$ alors $x_1 > x_2$.}
        \end{itemize}
        \item \textbf{Rédaction obligatoire au Probatoire :} Écrire « $\Delta = \dots$ », « Comme $\Delta \dots$, l'équation admet \dots solution(s) : $S = \{\dots\}$ ».
        \item \textbf{Astuces (racines évidentes) :} 
        \begin{itemize}
            \item Si $a+b+c=0$, alors $1$ est racine et l'autre est $c/a$.
            \item Si $a-b+c=0$, alors $-1$ est racine et l'autre est $-c/a$.
        \end{itemize}
    \end{enumerate}
\end{methode}

\begin{exoresolu}[Équation du second degré et dimensionnement d'une pièce]
    Un technicien doit usiner une pièce rectangulaire dont la longueur dépasse la largeur de $3$ cm. L'aire doit être de $40$ cm$^2$. Déterminer les dimensions de la pièce.
    \tcblower
    \textbf{Solution détaillée.}
    \begin{enumerate}
        \item \textbf{Modélisation :} Soit $x$ la largeur (en cm). La longueur est $x+3$. L'aire : $x(x+3) = 40$.
        \item \textbf{Équation :} $x^2 + 3x - 40 = 0$. Coefficients : $a=1$, $b=3$, $c=-40$.
        \item \textbf{Discriminant :} $\Delta = 3^2 - 4\times1\times(-40) = 9 + 160 = 169 = 13^2$.
        \item \textbf{Solutions :} $\Delta > 0$, deux solutions réelles distinctes.
        \begin{align*}
            x_1 &= \frac{-3 - 13}{2} = \frac{-16}{2} = -8 \\
            x_2 &= \frac{-3 + 13}{2} = \frac{10}{2} = 5
        \end{align*}
        \item \textbf{Interprétation physique :} Une largeur ne peut pas être négative. On rejette $x_1 = -8$.
        La largeur est $x = 5$ cm. La longueur est $5+3 = 8$ cm.
    \end{enumerate}
    \textbf{Conclusion :} La pièce mesure \textbf{5 cm de large sur 8 cm de long}.
    \begin{attention}[Piège contexte] Ne pas oublier de rejeter la solution négative incompatible avec le problème concret (longueur, temps, quantité, intensité).\end{attention}
\end{exoresolu}

\begin{methode}[Résoudre une inéquation du second degré $ax^2+bx+c \gtreqless 0$]
    \textbf{Objectif :} Déterminer l'ensemble de vérité $S$ en utilisant le signe du trinôme.
    \begin{enumerate}
        \item Résoudre l'équation associée $ax^2+bx+c=0$ (calculer $\Delta$, racines $x_1 \le x_2$ si réelles).
        \item \textbf{Règle du signe du trinôme (à connaître par cœur) :}
        \begin{itemize}
            \item $\Delta > 0$ : 
            \begin{itemize}
                \item Si $a > 0$ : $P(x) > 0$ pour $x < x_1$ ou $x > x_2$ (extérieur) ; $P(x) < 0$ pour $x_1 < x < x_2$ (intérieur).
                \item Si $a < 0$ : Signe inverse (positif à l'intérieur, négatif à l'extérieur).
            \end{itemize}
            \item $\Delta = 0$ : $P(x)$ a le signe de $a$ partout sauf en $x_0$ où il est nul.
            \item $\Delta < 0$ : $P(x)$ a le signe constant de $a$ sur $\mathbb{R}$.
        \end{itemize}
        \item Construire un \textbf{tableau de signes} (obligatoire au Probatoire) avec les valeurs remarquables (racines, bornes de l'intervalle d'étude) et les intervalles.
        \item Écrire la solution finale sous forme d'intervalle ou d'ensemble, en tenant compte du sens de l'inéquation ($\ge, >, \le, <$).
    \end{enumerate}
\end{methode}

\begin{exoresolu}[Inéquation du second degré et domaine de fonctionnement]
    Un moteur électrique fonctionne correctement si sa puissance $P(I) = -\num{0,5}I^2 + 4I - 3$ (en kW) est strictement positive, où $I$ est l'intensité en Ampères. Déterminer l'intervalle d'intensités admissibles.
    \tcblower
    \textbf{Solution détaillée.}
    \begin{enumerate}
        \item \textbf{Inéquation :} $-\num{0,5}I^2 + 4I - 3 > 0$. On multiplie par $-2$ (change le sens) : $I^2 - 8I + 6 < 0$.
        \item \textbf{Équation associée :} $I^2 - 8I + 6 = 0$. $a=1$, $b=-8$, $c=6$.
        $\Delta = (-8)^2 - 4\times6 = 64 - 24 = 40$. $\sqrt{\Delta} = 2\sqrt{10} \approx \num{6,32}$.
        Racines : $I_1 = \frac{8 - 2\sqrt{10}}{2} = 4 - \sqrt{10} \approx \num{0,84}$ ; $I_2 = 4 + \sqrt{10} \approx \num{7,16}$.
        \item \textbf{Tableau de signes pour $P(I)$ (coefficient initial $a=-\num{0,5} < 0$) :}
        \begin{center}
        \begin{tabular}{|c|ccccccc|}\hline
        $I$ & $-\infty$ & & $4-\sqrt{10}$ & & $4+\sqrt{10}$ & & $+\infty$ \\ \hline
        $P(I)$ & & $-$ & $0$ & $+$ & $0$ & $-$ & \\ \hline
        \end{tabular}
        \end{center}
        \item $P(I) > 0$ (positif) sur l'intervalle $]4-\sqrt{10} ; 4+\sqrt{10}[$.
    \end{enumerate}
    \textbf{Conclusion :} Le moteur fonctionne correctement pour une intensité $I \in ]4-\sqrt{10} ; 4+\sqrt{10}[$ A, soit environ \textbf{]}\num{0,84} ; \num{7,16}\textbf{[ A}.
\end{exoresolu}

\begin{methode}[Équations et inéquations associées à une fonction du second degré]
    \textbf{Objectif :} Résoudre $f(x)=k$, $f(x) \ge k$, $f(x) < k$ pour $f(x)=ax^2+bx+c$.
    \begin{enumerate}
        \item Ramener l'inconnue à gauche : $f(x) - k = 0$ ou $f(x) - k \gtreqless 0$.
        \item On obtient un trinôme du second degré $ax^2+bx+(c-k)$.
        \item Appliquer les méthodes de résolution d'équation ou d'inéquation du second degré vues précédemment.
        \item \textbf{Lien graphique :} Les solutions de $f(x)=k$ sont les abscisses des points d'intersection de la courbe $C_f$ avec la droite $y=k$. Les solutions de $f(x) > k$ correspondent aux $x$ pour lesquels la courbe est \emph{au-dessus} de la droite.
    \end{enumerate}
\end{methode}

\begin{methode}[Résoudre un système linéaire $2\times2$ (Substitution / Combinaison linéaire)]
    \textbf{Objectif :} Trouver le couple $(x;y)$ solution du système $\begin{cases} a_1x+b_1y=c_1 \\ a_2x+b_2y=c_2 \end{cases}$.
    \begin{enumerate}
        \item \textbf{Méthode par combinaison (souvent plus rapide) :}
        \begin{itemize}
            \item Multiplier les équations par des nombres non nuls pour obtenir des coefficients opposés sur $x$ ou sur $y$.
            \item Additionner les deux équations pour éliminer une inconnue.
            \item Résoudre l'équation à une inconnue trouvée.
            \item Remplacer dans une équation initiale pour trouver l'autre inconnue.
        \end{itemize}
        \item \textbf{Méthode par substitution :} Isoler $x$ (ou $y$) dans une équation, substituer l'expression obtenue dans l'autre équation.
        \item \textbf{Vérification obligatoire :} Remplacer $(x;y)$ dans les \textbf{deux} équations initiales.
        \item \textbf{Cas particuliers (théoriques, discriminant $\delta = a_1b_2 - a_2b_1$) :}
        \begin{itemize}
            \item $\delta \neq 0$ : Solution unique (droites sécantes).
            \item $\delta = 0$ et $\frac{a_1}{a_2} = \frac{b_1}{b_2} = \frac{c_1}{c_2}$ : Infinité de solutions (droites confondues).
            \item $\delta = 0$ et $\frac{a_1}{a_2} = \frac{b_1}{b_2} \neq \frac{c_1}{c_2}$ : Pas de solution (droites parallèles distinctes).
        \end{itemize}
    \end{enumerate}
\end{methode}

\begin{exoresolu}[Système linéaire : mélange de deux alliages]
    On mélange un alliage A à \num{20}\% de cuivre avec un alliage B à \num{50}\% de cuivre pour obtenir \num{100} kg d'un alliage à \num{30}\% de cuivre. Déterminer les masses de A et de B à mélanger.
    \tcblower
    \textbf{Solution détaillée.}
    \begin{enumerate}
        \item \textbf{Modélisation :} Soit $x$ la masse (en kg) de l'alliage A et $y$ la masse (en kg) de l'alliage B.
        \item \textbf{Système :}
        \[\begin{cases}
            x + y = 100 & \text{(masse totale)} \\
            \num{0,20}x + \num{0,50}y = \num{30} & \text{(masse de cuivre : \num{30}\% de 100 = 30 kg)}
        \end{cases}\]
        \item \textbf{Résolution par substitution :}
        D'après la première équation : $y = 100 - x$.
        Substitution dans la seconde :
        \begin{align*}
            \num{0,20}x + \num{0,50}(100 - x) &= 30 \\
            \num{0,20}x + 50 - \num{0,50}x &= 30 \\
            -\num{0,30}x &= -20 \\
            x &= \frac{20}{\num{0,30}} = \frac{200}{3} \approx \num{66,67}
        \end{align*}
        Alors $y = 100 - \frac{200}{3} = \frac{100}{3} \approx \num{33,33}$.
        \item \textbf{Vérification :}
        \begin{itemize}
            \item Masse totale : $\frac{200}{3} + \frac{100}{3} = 100$ kg. \checkmark
            \item Cuivre : $\num{0,20}\times\frac{200}{3} + \num{0,50}\times\frac{100}{3} = \frac{40}{3} + \frac{50}{3} = \frac{90}{3} = 30$ kg. \checkmark
        \end{itemize}
    \end{enumerate}
    \textbf{Conclusion :} Il faut mélanger \textbf{\num{66,67} kg de l'alliage A} et \textbf{\num{33,33} kg de l'alliage B} (soit $\frac{200}{3}$ kg et $\frac{100}{3}$ kg exactement).
\end{exoresolu}

\begin{exoresolu}[Synthèse : Modélisation d'une activité de production]
    Une petite entreprise fabrique et vend des pièces mécaniques. Le coût total de production (en milliers de FCFA) pour $x$ centaines de pièces est $C(x) = 2x^2 - 20x + 100$. Le prix de vente unitaire est fixé à \num{5000} FCFA la pièce (soit \num{500} milliers de FCFA la centaine de pièces).
    \begin{enumerate}
        \item Exprimer le chiffre d'affaires $R(x)$ et le bénéfice $B(x)$ en fonction de $x$.
        \item Déterminer le seuil de rentabilité (point mort) : nombre de pièces pour lequel $B(x)=0$.
        \item Déterminer l'intervalle de production pour lequel l'entreprise est bénéficiaire ($B(x) > 0$).
        \item Trouver la production qui maximise le bénéfice et la valeur de ce maximum.
    \end{enumerate}
    \tcblower
    \textbf{Solution détaillée.}
    \begin{enumerate}
        \item Chiffre d'affaires : $R(x) = \num{500}x$ (en milliers de FCFA).
        Bénéfice : $B(x) = R(x) - C(x) = 500x - (2x^2 - 20x + 100) = -2x^2 + 520x - 100$.
        \item \textbf{Seuil de rentabilité :} $B(x) = 0 \Leftrightarrow -2x^2 + 520x - 100 = 0 \Leftrightarrow x^2 - 260x + 50 = 0$.
        $\Delta = 260^2 - 4\times50 = 67600 - 200 = 67400$. $\sqrt{\Delta} = 10\sqrt{674} \approx \num{259,6}$.
        $x_1 = \frac{260 - 10\sqrt{674}}{2} = 130 - 5\sqrt{674} \approx \num{0,19}$.
        $x_2 = 130 + 5\sqrt{674} \approx \num{259,8}$.
        Seuils : environ \textbf{19 pièces} et \textbf{25 980 pièces}.
        \item \textbf{Zone bénéficiaire :} $a=-2 < 0$, $B(x) > 0$ entre les racines.
        $B(x) > 0 \Leftrightarrow x \in ]130 - 5\sqrt{674} ; 130 + 5\sqrt{674}[$.
        Soit environ \textbf{]0,19 ; 259,8[} centaines de pièces, c'est-à-dire \textbf{]19 ; 25 980[ pièces}.
        \item \textbf{Maximum de bénéfice :} Forme canonique ou sommet $\alpha = -\frac{b}{2a} = -\frac{520}{2\times(-2)} = 130$.
        $B(130) = -2(130)^2 + 520(130) - 100 = -33800 + 67600 - 100 = \num{33700}$ (milliers de FCFA).
        Production optimale : \textbf{13 000 pièces}. Bénéfice maximal : \textbf{33 700 000 FCFA}.
    \end{enumerate}
    \begin{attention}[Unités] Attention aux unités : $x$ en centaines de pièces, $C, R, B$ en milliers de FCFA. Bien convertir pour la conclusion finale.
    \end{attention}
\end{exoresolu}

\newpage

\coursec{Exercices}

\courssub{Je vérifie mes connaissances}

\begin{exercice}[QCM : le discriminant]
Pour chaque question, une seule réponse est exacte. Entoure-la et justifie ton choix.
\begin{enumerate}
\item Le discriminant du trinôme $2x^2 - 5x + 1$ vaut :
\begin{enumerate}
\item[$a.$] $17$ \qquad $b.$ $33$ \qquad $c.$ $-17$
\end{enumerate}
\item Si $\Delta = 0$, l'équation $ax^2 + bx + c = 0$ admet :
\begin{enumerate}
\item[$a.$] deux solutions distinctes \qquad $b.$] une solution double \qquad $c.$] aucune solution réelle
\end{enumerate}
\item L'équation $x^2 + 4x + 4 = 0$ admet pour solution :
\begin{enumerate}
\item[$a.$] $x = 2$ \qquad $b.$ $x = -2$ \qquad $c.$ $x = 4$
\end{enumerate}
\item Le système $\begin{cases} x + y = 5 \\ x - y = 1 \end{cases}$ admet pour solution :
\begin{enumerate}
\item[$a.$] $(3\,;\,2)$ \qquad $b.$ $(2\,;\,3)$ \qquad $c.$ $(4\,;\,1)$
\end{enumerate}
\end{enumerate}
\end{exercice}

\begin{exercice}[Vrai ou faux]
Réponds par Vrai ou Faux en justifiant chaque réponse.
\begin{enumerate}
\item Si $\Delta < 0$, le trinôme $ax^2 + bx + c$ est toujours du signe de $a$.
\item Le trinôme $x^2 - 6x + 9$ se factorise en $(x - 3)(x + 3)$.
\item La forme canonique de $x^2 - 4x + 1$ est $(x - 2)^2 - 3$.
\item Un système de deux équations linéaires à deux inconnues admet toujours une solution unique.
\end{enumerate}
\end{exercice}

\begin{exercice}[Somme et produit des racines]
On considère l'équation $3x^2 - 7x + 2 = 0$.
\begin{enumerate}
\item Sans résoudre l'équation, calculer la somme $S$ et le produit $P$ des racines.
\item Vérifier que $x_1 = 2$ est une solution, puis en déduire l'autre solution $x_2$.
\end{enumerate}
\end{exercice}

\begin{exercice}[Calculs directs]
\begin{enumerate}
\item Calculer le discriminant de chacun des trinômes : $x^2 - 5x + 6$ ; $-x^2 + 2x + 8$ ; $2x^2 + x + 5$.
\item Donner, sans calcul supplémentaire, le nombre de solutions réelles de chaque équation associée.
\end{enumerate}
\end{exercice}

\courssub{Je m'entraîne}

\begin{exercice}[Résolution d'équations du second degré]
Résoudre dans $\mathbb{R}$ les équations suivantes, puis vérifier les résultats à l'aide de la somme et du produit des racines :
\begin{enumerate}
\item $3x^2 - 7x + 2 = 0$ ;
\item $x^2 + 4x + 4 = 0$ ;
\item $2x^2 + x + 5 = 0$.
\end{enumerate}
\end{exercice}

\begin{exercice}[Forme canonique et maximum]
On considère la fonction $f$ définie sur $\mathbb{R}$ par $f(x) = -2x^2 + 8x - 5$.
\begin{enumerate}
\item Mettre $f(x)$ sous forme canonique.
\item En déduire les coordonnées du sommet de la parabole représentant $f$.
\item Quel est le maximum de $f$ sur $\mathbb{R}$ ? Pour quelle valeur de $x$ est-il atteint ?
\end{enumerate}
\end{exercice}

\begin{exercice}[Inéquations du second degré]
\begin{enumerate}
\item Résoudre dans $\mathbb{R}$ l'inéquation $x^2 - 5x + 6 < 0$ en dressant un tableau de signes.
\item Résoudre dans $\mathbb{R}$ l'inéquation $-x^2 + 2x + 8 \geq 0$.
\item Représenter sur une droite graduée l'ensemble des solutions de chaque inéquation.
\end{enumerate}
\end{exercice}

\begin{exercice}[Système de deux équations]
On considère le système $\begin{cases} 3x + 2y = 16 \\ 5x - y = 5 \end{cases}$.
\begin{enumerate}
\item Résoudre ce système par la méthode de substitution.
\item Résoudre ce même système par la méthode de combinaison (addition).
\item Interpréter graphiquement le résultat obtenu (que représentent les deux équations et le couple solution ?).
\end{enumerate}
\end{exercice}

\courssub{Je me prépare au Probatoire}

\begin{exercice}[Rentabilité d'un atelier de menuiserie]
Un atelier de menuiserie de Douala fabrique et vend des tabourets. Le bénéfice mensuel, exprimé en milliers de FCFA, est donné en fonction du nombre $x$ de tabourets vendus par :
\[ B(x) = -x^2 + 40x - 300. \]
\begin{enumerate}
\item Calculer le bénéfice réalisé pour $x = 10$ tabourets vendus, puis pour $x = 35$ tabourets. Commenter.
\item Mettre $B(x)$ sous forme canonique.
\item En déduire le nombre de tabourets à vendre pour que le bénéfice soit maximal, et donner la valeur de ce bénéfice maximal en FCFA.
\item Résoudre l'équation $B(x) = 0$.
\item À l'aide d'un tableau de signes, résoudre l'inéquation $B(x) > 0$.
\item En déduire le nombre entier de tabourets que l'atelier doit vendre pour être rentable (bénéfice strictement positif).
\end{enumerate}
\end{exercice}

\begin{exercice}[Au marché de Mokolo]
Au marché, une ménagère achète $3$ kg de tomates et $2$ kg d'oignons pour un total de \num{2900} FCFA. Le lendemain, aux mêmes prix, elle achète $2$ kg de tomates et $5$ kg d'oignons pour \num{4000} FCFA.
\begin{enumerate}
\item En désignant par $x$ le prix d'un kilogramme de tomates et par $y$ celui d'un kilogramme d'oignons, traduire la situation par un système de deux équations à deux inconnues.
\item Résoudre ce système par la méthode de ton choix, en détaillant les étapes.
\item En déduire le prix d'un kilogramme de tomates et celui d'un kilogramme d'oignons.
\item Combien paiera-t-elle pour $4$ kg de tomates et $3$ kg d'oignons ?
\item \textbf{Bonus.} Deux nombres ont pour somme $26$ et pour produit $165$. Montrer que ces nombres sont les solutions d'une équation du second degré, puis les déterminer.
\end{enumerate}
\end{exercice}

\begin{exercice}[Équation bicarrée et clôture d'un terrain]
\textbf{Partie A.} On considère l'équation $(E)$ : $x^4 - 13x^2 + 36 = 0$.
\begin{enumerate}
\item En posant $X = x^2$, montrer que l'équation $(E)$ se ramène à une équation du second degré d'inconnue $X$.
\item Résoudre cette équation du second degré.
\item En déduire toutes les solutions réelles de l'équation $(E)$.
\end{enumerate}
\textbf{Partie B.} Un technicien doit clôturer un terrain rectangulaire de périmètre $52$ m et d'aire $165$ m$^2$. On note $L$ et $l$ les dimensions du terrain.
\begin{enumerate}
\item Montrer que $L + l = 26$ et $L \times l = 165$.
\item En déduire que $L$ et $l$ sont solutions de l'équation $t^2 - 26t + 165 = 0$.
\item Résoudre cette équation et donner les dimensions du terrain.
\end{enumerate}
\end{exercice}

\newpage

\coursec{Corrigés des exercices}

\coursec{Corrigés des exercices — Équations, inéquations et systèmes}

\begin{corrige}[Exercice 1 — QCM : le discriminant]
\begin{enumerate}
\item \textbf{Réponse $a.$ : $17$.} Pour $2x^2 - 5x + 1$, on a $a = 2$, $b = -5$, $c = 1$.
\[ \Delta = b^2 - 4ac = (-5)^2 - 4 \times 2 \times 1 = 25 - 8 = 17. \]
\item \textbf{Réponse $b.$ : une solution double.} Si $\Delta = 0$, l'équation admet une unique solution (dite \cle{solution double}) $x_0 = -\dfrac{b}{2a}$.
\item \textbf{Réponse $b.$ : $x = -2$.} On reconnaît l'identité remarquable : $x^2 + 4x + 4 = (x+2)^2$. L'équation $(x+2)^2 = 0$ donne $x = -2$. Vérification : $(-2)^2 + 4(-2) + 4 = 4 - 8 + 4 = 0$.
\item \textbf{Réponse $a.$ : $(3\,;\,2)$.} En additionnant les deux équations : $2x = 6$, donc $x = 3$. Alors $y = 5 - x = 2$. Vérification : $3 - 2 = 1$. \checkmark
\end{enumerate}
\end{corrige}

\begin{corrige}[Exercice 2 — Vrai ou faux]
\begin{enumerate}
\item \textbf{Vrai.} Si $\Delta < 0$, le trinôme ne s'annule jamais et garde un signe constant : il est toujours du \cle{signe de $a$} (strictement positif si $a > 0$, strictement négatif si $a < 0$).
\item \textbf{Faux.} $x^2 - 6x + 9 = (x-3)^2 = (x-3)(x-3)$ (identité remarquable $a^2 - 2ab + b^2$). Or $(x-3)(x+3) = x^2 - 9$, ce qui ne correspond pas.
\item \textbf{Vrai.} $x^2 - 4x + 1 = (x^2 - 4x + 4) - 4 + 1 = (x-2)^2 - 3$.
\item \textbf{Faux.} Un tel système peut avoir une solution unique (droites sécantes), aucune solution (droites parallèles distinctes) ou une infinité de solutions (droites confondues). Exemple sans solution : $\begin{cases} x + y = 1 \\ x + y = 3 \end{cases}$.
\end{enumerate}
\end{corrige}

\begin{corrige}[Exercice 3 — Somme et produit des racines]
\begin{enumerate}
\item Pour $3x^2 - 7x + 2 = 0$ ($a = 3$, $b = -7$, $c = 2$), comme $\Delta = 49 - 24 = 25 > 0$, il y a deux racines et :
\[ S = -\frac{b}{a} = \frac{7}{3} \qquad ; \qquad P = \frac{c}{a} = \frac{2}{3}. \]
\item Vérifions que $x_1 = 2$ est solution : $3(2)^2 - 7(2) + 2 = 12 - 14 + 2 = 0$. \checkmark \\
On utilise le produit : $x_1 \times x_2 = P$, donc $2x_2 = \dfrac{2}{3}$, d'où $x_2 = \dfrac{1}{3}$. \\
Contrôle avec la somme : $2 + \dfrac{1}{3} = \dfrac{7}{3} = S$. \checkmark
\end{enumerate}
\end{corrige}

\begin{corrige}[Exercice 4 — Calculs directs]
\begin{enumerate}
\item
\begin{itemize}
\item $x^2 - 5x + 6$ : $\Delta = (-5)^2 - 4 \times 1 \times 6 = 25 - 24 = 1$.
\item $-x^2 + 2x + 8$ : $\Delta = 2^2 - 4 \times (-1) \times 8 = 4 + 32 = 36$.
\item $2x^2 + x + 5$ : $\Delta = 1^2 - 4 \times 2 \times 5 = 1 - 40 = -39$.
\end{itemize}
\item D'après le signe du discriminant :
\begin{itemize}
\item $\Delta = 1 > 0$ : l'équation $x^2 - 5x + 6 = 0$ admet \textbf{deux solutions réelles distinctes} ;
\item $\Delta = 36 > 0$ : l'équation $-x^2 + 2x + 8 = 0$ admet \textbf{deux solutions réelles distinctes} ;
\item $\Delta = -39 < 0$ : l'équation $2x^2 + x + 5 = 0$ n'admet \textbf{aucune solution réelle}.
\end{itemize}
\end{enumerate}
\end{corrige}

\begin{corrige}[Exercice 5 — Résolution d'équations du second degré]
\begin{enumerate}
\item $3x^2 - 7x + 2 = 0$ : $\Delta = (-7)^2 - 4 \times 3 \times 2 = 49 - 24 = 25 > 0$, $\sqrt{\Delta} = 5$.
\[ x_1 = \frac{7 - 5}{6} = \frac{2}{6} = \frac{1}{3} \qquad ; \qquad x_2 = \frac{7 + 5}{6} = \frac{12}{6} = 2. \]
$\mathcal{S} = \left\{ \dfrac{1}{3}\,;\,2 \right\}$. Vérification : $S = \dfrac{1}{3} + 2 = \dfrac{7}{3} = -\dfrac{b}{a}$ et $P = \dfrac{2}{3} = \dfrac{c}{a}$. \checkmark
\item $x^2 + 4x + 4 = 0$ : $\Delta = 16 - 16 = 0$. Solution double : $x_0 = -\dfrac{b}{2a} = -\dfrac{4}{2} = -2$. \\
$\mathcal{S} = \{-2\}$. Vérification : $x^2 + 4x + 4 = (x+2)^2$, donc $S = -4 = -\dfrac{b}{a}$ et $P = 4 = \dfrac{c}{a}$. \checkmark
\item $2x^2 + x + 5 = 0$ : $\Delta = 1 - 40 = -39 < 0$. L'équation n'admet aucune solution réelle : $\mathcal{S} = \emptyset$.
\end{enumerate}
\end{corrige}

\begin{corrige}[Exercice 6 — Forme canonique et maximum]
\begin{enumerate}
\item $f(x) = -2x^2 + 8x - 5$. On factorise par $-2$ les deux premiers termes :
\begin{align*}
f(x) &= -2\left(x^2 - 4x\right) - 5 = -2\left[(x-2)^2 - 4\right] - 5 \\
&= -2(x-2)^2 + 8 - 5 = -2(x-2)^2 + 3.
\end{align*}
\item La parabole a pour sommet le point $S(2\,;\,3)$. Comme $a = -2 < 0$, elle est tournée vers le bas.
\item Le maximum de $f$ sur $\mathbb{R}$ est $3$, atteint pour $x = 2$ (car $-2(x-2)^2 \leq 0$, nul seulement en $x = 2$).
\end{enumerate}
\end{corrige}

\begin{corrige}[Exercice 7 — Inéquations du second degré]
\begin{enumerate}
\item $x^2 - 5x + 6 < 0$. Racines : $\Delta = 1$, $x_1 = \dfrac{5-1}{2} = 2$, $x_2 = \dfrac{5+1}{2} = 3$. Le trinôme ($a = 1 > 0$) est négatif \emph{entre} les racines.
\begin{center}
\begin{tabular}{|c|ccccccc|}\hline
$x$ & $-\infty$ & & $2$ & & $3$ & & $+\infty$ \\ \hline
$x^2 - 5x + 6$ & $+$ & $+$ & $0$ & $-$ & $0$ & $+$ & $+$ \\ \hline
\end{tabular}
\end{center}
$\mathcal{S} = \left]2\,;\,3\right[$ (bornes exclues car inégalité stricte).
\item $-x^2 + 2x + 8 \geq 0$. Racines : $\Delta = 36$, $x_1 = \dfrac{-2-6}{-2} = 4$, $x_2 = \dfrac{-2+6}{-2} = -2$. Le trinôme ($a = -1 < 0$) est positif \emph{entre} les racines.
\begin{center}
\begin{tabular}{|c|ccccccc|}\hline
$x$ & $-\infty$ & & $-2$ & & $4$ & & $+\infty$ \\ \hline
$-x^2 + 2x + 8$ & $-$ & $-$ & $0$ & $+$ & $0$ & $-$ & $-$ \\ \hline
\end{tabular}
\end{center}
$\mathcal{S} = \left[-2\,;\,4\right]$ (bornes incluses car inégalité large).
\item Représentation sur la droite graduée :
\begin{popfigure}
\begin{tikzpicture}[scale=1.1]
\draw[->,popline] (-3,0) -- (5.5,0);
\foreach \x in {-2,-1,0,1,2,3,4,5} \draw (\x,0.08) -- (\x,-0.08) node[below] {\small $\x$};
% Solution 1 : ]2;3[
\draw[very thick,popblue] (2,0.25) -- (3,0.25);
\draw[popblue] (2,0.15) -- (2,0.35); \draw[popblue] (3,0.15) -- (3,0.35);
\draw[popblue,fill=white] (2,0.25) circle (0.06); \draw[popblue,fill=white] (3,0.25) circle (0.06);
\node[popblue,above] at (2.5,0.45) {\small $]2\,;\,3[$};
% Solution 2 : [-2;4]
\draw[very thick,poporange] (-2,-0.55) -- (4,-0.55);
\fill[poporange] (-2,-0.55) circle (0.06); \fill[poporange] (4,-0.55) circle (0.06);
\node[poporange,below] at (1,-0.7) {\small $[-2\,;\,4]$};
\end{tikzpicture}
\legende{En bleu : solutions de la question 1 (bornes exclues) ; en orange : solutions de la question 2 (bornes incluses).}
\end{popfigure}
\end{enumerate}
\end{corrige}

\begin{corrige}[Exercice 8 — Système de deux équations]
\begin{enumerate}
\item \textbf{Méthode par substitution.} De la deuxième équation : $y = 5x - 5$. On reporte dans la première :
\[ 3x + 2(5x - 5) = 16 \iff 3x + 10x - 10 = 16 \iff 13x = 26 \iff x = 2. \]
Alors $y = 5 \times 2 - 5 = 5$. Vérification : $3(2) + 2(5) = 16$ \checkmark\ et $5(2) - 5 = 5$ \checkmark. \\
$\mathcal{S} = \{(2\,;\,5)\}$.
\item \textbf{Méthode par combinaison linéaire.} On multiplie la deuxième équation par $2$ : $10x - 2y = 10$. On additionne membre à membre avec la première :
\[ (3x + 2y) + (10x - 2y) = 16 + 10 \iff 13x = 26 \iff x = 2, \]
puis $y = 5x - 5 = 5$. On retrouve $\mathcal{S} = \{(2\,;\,5)\}$.
\item Chaque équation représente une \cle{droite} dans le plan : $(d_1) : y = -\dfrac{3}{2}x + 8$ et $(d_2) : y = 5x - 5$. Le couple solution $(2\,;\,5)$ est le \cle{point d'intersection} de ces deux droites : elles sont sécantes, d'où l'unicité de la solution.
\end{enumerate}
\end{corrige}

\begin{corrige}[Exercice 9 — Rentabilité d'un atelier de menuiserie]
\textbf{Barème indicatif (sur 10) :} Q1 : 1 pt ; Q2 : 2 pts ; Q3 : 1,5 pt ; Q4 : 2 pts ; Q5 : 2 pts ; Q6 : 1,5 pt.
\begin{enumerate}
\item $B(10) = -100 + 400 - 300 = 0$ : pour 10 tabourets, le bénéfice est nul (seuil de rentabilité). \\
$B(35) = -1225 + 1400 - 300 = -125$ : pour 35 tabourets, l'atelier perd \num{125000} FCFA (surproduction : les coûts dépassent les recettes).
\item $B(x) = -\left(x^2 - 40x\right) - 300 = -\left[(x-20)^2 - 400\right] - 300 = -(x-20)^2 + 100$.
\item Comme $a = -1 < 0$, le maximum est atteint au sommet : pour $x = 20$ tabourets, et vaut $B(20) = 100$ milliers de FCFA, soit \num{100000} FCFA.
\item $B(x) = 0 \iff (x-20)^2 = 100 \iff x - 20 = 10$ ou $x - 20 = -10$, d'où $x_1 = 10$ et $x_2 = 30$. \\
(Contrôle par le discriminant : $\Delta = 1600 - 1200 = 400$, $\sqrt{\Delta} = 20$, $x = \dfrac{-40 \pm 20}{-2}$, soit $10$ et $30$. \checkmark)
\item Trinôme avec $a = -1 < 0$ : il est strictement positif entre les racines.
\begin{center}
\begin{tabular}{|c|ccccccc|}\hline
$x$ & $-\infty$ & & $10$ & & $30$ & & $+\infty$ \\ \hline
$B(x)$ & $-$ & $-$ & $0$ & $+$ & $0$ & $-$ & $-$ \\ \hline
\end{tabular}
\end{center}
$\mathcal{S} = \left]10\,;\,30\right[$.
\item L'atelier est rentable pour $x \in \left]10\,;\,30\right[$, c'est-à-dire pour un nombre \textbf{entier} de tabourets compris entre $11$ et $29$ inclus.
\end{enumerate}
\textbf{Points de vigilance :} citer le signe de $a$ pour justifier la position du maximum et le tableau de signes ; convertir les milliers de FCFA en FCFA dans la réponse finale ; exclure les bornes $10$ et $30$ car le bénéfice doit être \emph{strictement} positif.
\end{corrige}

\begin{corrige}[Exercice 10 — Au marché de Mokolo]
\textbf{Barème indicatif (sur 10) :} Q1 : 2 pts ; Q2 : 3 pts ; Q3 : 1 pt ; Q4 : 1,5 pt ; Q5 (bonus) : 2,5 pts.
\begin{enumerate}
\item Soit $x$ le prix du kg de tomates (en FCFA) et $y$ le prix du kg d'oignons (en FCFA).
\[ \begin{cases} 3x + 2y = 2900 \\ 2x + 5y = 4000 \end{cases} \]
\item \textbf{Résolution par combinaison linéaire (élimination de $x$).} On multiplie la première équation par $2$ et la deuxième par $3$ :
\[ \begin{cases} 6x + 4y = 5800 \\ 6x + 15y = 12000 \end{cases} \]
On soustrait la première de la deuxième : $11y = 6200$, d'où $y = \dfrac{6200}{11}$.
On reporte dans la première équation : $3x + 2\times\dfrac{6200}{11} = 2900 \iff 3x = \dfrac{31900 - 12400}{11} = \dfrac{19500}{11} \iff x = \dfrac{6500}{11}$.
Vérification : $2 \times \dfrac{6500}{11} + 5 \times \dfrac{6200}{11} = \dfrac{13000 + 31000}{11} = \dfrac{44000}{11} = 4000$ \checkmark. \\
$\mathcal{S} = \left\{ \left(\dfrac{6500}{11}\,;\,\dfrac{6200}{11}\right) \right\}$.
\item Le kilogramme de tomates coûte $\dfrac{6500}{11} \approx 591$ FCFA et le kilogramme d'oignons $\dfrac{6200}{11} \approx 564$ FCFA (prix exacts : $\frac{6500}{11}$ FCFA et $\frac{6200}{11}$ FCFA).
\item $4x + 3y = 4\times\dfrac{6500}{11} + 3\times\dfrac{6200}{11} = \dfrac{26000 + 18600}{11} = \dfrac{44600}{11} \approx \num{4055}$ FCFA.
\item \textbf{Bonus.} Deux nombres de somme $S = 26$ et de produit $P = 165$ sont les solutions de l'équation $t^2 - St + P = 0$, soit $t^2 - 26t + 165 = 0$. \\
$\Delta = 676 - 660 = 16$, $\sqrt{\Delta} = 4$ : $t_1 = \dfrac{26 - 4}{2} = 11$, $t_2 = \dfrac{26 + 4}{2} = 15$. \\
Les deux nombres sont $11$ et $15$. Vérification : $11 + 15 = 26$ et $11 \times 15 = 165$. \checkmark
\end{enumerate}
\textbf{Points de vigilance :} bien identifier les inconnues avant de poser le système ; vérifier la solution dans \emph{les deux} équations ; au bonus, citer la propriété « somme $S$, produit $P$ $\Rightarrow$ solutions de $t^2 - St + P = 0$ ».
\end{corrige}

\begin{corrige}[Exercice 11 — Équation bicarrée et clôture d'un terrain]
\textbf{Barème indicatif (sur 10) :} Partie A : 5 pts (Q1 : 1 pt ; Q2 : 2 pts ; Q3 : 2 pts) ; Partie B : 5 pts (Q1 : 1,5 pt ; Q2 : 1,5 pt ; Q3 : 2 pts).

\textbf{Partie A.}
\begin{enumerate}
\item Posons $X = x^2$. Alors $x^4 = (x^2)^2 = X^2$, et $(E)$ devient :
\[ X^2 - 13X + 36 = 0, \]
équation du second degré d'inconnue $X$.
\item $\Delta = (-13)^2 - 4 \times 36 = 169 - 144 = 25 > 0$, $\sqrt{\Delta} = 5$.
\[ X_1 = \frac{13 - 5}{2} = 4 \qquad ; \qquad X_2 = \frac{13 + 5}{2} = 9. \]
\item On revient à $x$ : $x^2 = 4$ donne $x = 2$ ou $x = -2$ ; $x^2 = 9$ donne $x = 3$ ou $x = -3$. \\
$\mathcal{S} = \{-3\,;\,-2\,;\,2\,;\,3\}$. \\
Vérification pour $x = 3$ : $81 - 13 \times 9 + 36 = 81 - 117 + 36 = 0$. \checkmark
\end{enumerate}

\textbf{Partie B.}
\begin{enumerate}
\item Le périmètre vaut $2(L + l) = 52$, donc $L + l = 26$. L'aire vaut $L \times l = 165$ m$^2$.
\item $L$ et $l$ sont deux nombres de somme $S = 26$ et de produit $P = 165$ : ils sont donc solutions de l'équation $t^2 - St + P = 0$, c'est-à-dire $t^2 - 26t + 165 = 0$.
\item $\Delta = 676 - 660 = 16$, $\sqrt{\Delta} = 4$ : $t_1 = \dfrac{26 - 4}{2} = 11$ et $t_2 = \dfrac{26 + 4}{2} = 15$. \\
Le terrain a pour dimensions $L = 15$ m et $l = 11$ m. Vérification : périmètre $= 2(15 + 11) = 52$ m \checkmark\ ; aire $= 15 \times 11 = 165$ m$^2$ \checkmark.
\end{enumerate}
\textbf{Points de vigilance :} dans la Partie A, ne pas oublier de revenir à $x$ et de donner les \emph{quatre} solutions (les valeurs négatives aussi, car $x^2 = X$ avec $X > 0$ donne deux solutions) ; dans la Partie B, justifier le passage « somme et produit $\Rightarrow$ équation $t^2 - St + P = 0$ » et conclure par une phrase avec les unités.
\end{corrige}

\newpage

\coursec{Fiche bilan}

\begin{popfigure}
\begin{tikzpicture}[mindmap, grow cyclic, every node/.style={concept, text=white, font=\small}, root concept/.append style={concept color=popdark, font=\bfseries}, level 1 concept/.append style={level distance=3.6cm, sibling angle=60, font=\scriptsize}]
\node[root concept]{Équations, inéquations, systèmes}
  child[concept color=popblue]{ node {Trinôme $ax^2+bx+c$ et forme canonique $a(x-\alpha)^2+\beta$} }
  child[concept color=poporange]{ node {Discriminant $\Delta=b^2-4ac$} }
  child[concept color=popgreen]{ node {Racines $x_1,x_2$ et factorisation} }
  child[concept color=poppurple]{ node {Signe du trinôme et inéquations} }
  child[concept color=poppink]{ node {Systèmes $2\times 2$ linéaires} }
  child[concept color=popgold]{ node {Équations produit, quotient, $\sqrt{}$} };
\end{tikzpicture}
\legende{Carte mentale de la leçon : les six idées clés à maîtriser.}
\end{popfigure}

\begin{bilan}
\textbf{Définitions clés}
\begin{itemize}
\item \cle{Trinôme du second degré} : expression $f(x)=ax^2+bx+c$ avec $a\neq 0$.
\item \cle{Forme canonique} : $f(x)=a(x-\alpha)^2+\beta$ avec $\alpha=-\dfrac{b}{2a}$ et $\beta=f(\alpha)=-\dfrac{\Delta}{4a}$.
\item \cle{Discriminant} : $\Delta=b^2-4ac$.
\item \cle{Racine} : valeur de $x$ qui annule le trinôme.
\end{itemize}

\textbf{Formules à connaître par cœur}
\begin{center}
\begin{tabularx}{\linewidth}{|l|X|X|}
\hline
\rowcolor{popblueL} \textbf{Cas} & \textbf{Équation $ax^2+bx+c=0$} & \textbf{Factorisation} \\
\hline
$\Delta>0$ & $x_1=\dfrac{-b-\sqrt{\Delta}}{2a}$, $x_2=\dfrac{-b+\sqrt{\Delta}}{2a}$ & $a(x-x_1)(x-x_2)$ \\
\hline
$\Delta=0$ & $x_0=-\dfrac{b}{2a}$ (racine double) & $a(x-x_0)^2$ \\
\hline
$\Delta<0$ & Pas de solution dans $\mathbb{R}$ & Pas de factorisation dans $\mathbb{R}$ \\
\hline
Somme et produit & \multicolumn{2}{X|}{$x_1+x_2=-\dfrac{b}{a}$ \quad et \quad $x_1x_2=\dfrac{c}{a}$} \\
\hline
\end{tabularx}
\end{center}

\textbf{Signe du trinôme} : si $\Delta>0$, le trinôme est du \cle{signe de $a$} à l'extérieur des racines, du signe contraire entre les racines ; si $\Delta\leq 0$, il est toujours du signe de $a$ (nul en $x_0$ si $\Delta=0$).

\textbf{Méthodes express}
\begin{itemize}
\item \cle{Système $2\times 2$} : substitution ou combinaison linéaire ; unique solution si $ab'-a'b\neq 0$.
\item \cle{Équation produit} $(A)(B)=0$ : $A=0$ ou $B=0$.
\item \cle{Équation quotient} $\dfrac{A}{B}=0$ : $A=0$ avec $B\neq 0$ (valeurs interdites d'abord !).
\item \cle{Équation} $\sqrt{A}=B$ : condition $B\geq 0$, puis $A=B^2$.
\item \cle{Inéquation} : se ramener à une comparaison à $0$, étudier le signe, conclure par un intervalle.
\end{itemize}
\end{bilan}

\begin{aretenir}[Checklist avant le Probatoire]
\begin{itemize}
\item[$\square$] Je sais passer de la forme développée à la forme canonique.
\item[$\square$] Je calcule $\Delta=b^2-4ac$ sans erreur de signe.
\item[$\square$] Je connais les formules des racines selon le signe de $\Delta$.
\item[$\square$] Je factorise un trinôme quand $\Delta\geq 0$.
\item[$\square$] Je dresse le tableau de signes d'un trinôme.
\item[$\square$] Je résous une inéquation du second degré et je donne la solution en intervalle.
\item[$\square$] Je résous un système de deux équations linéaires par substitution et par combinaison.
\item[$\square$] Je résous une équation produit nul.
\item[$\square$] Je pense aux valeurs interdites pour un quotient et aux conditions pour $\sqrt{A}$.
\item[$\square$] Je sais modéliser un problème concret (périmètre, coût, recette) par une équation du second degré.
\item[$\square$] Je rédige ma solution en justifiant chaque étape.
\end{itemize}
\end{aretenir}

\findecours

\end{document}
